% concatenated from main.tex
\documentclass[11pt]{article}
\usepackage[utf8]{inputenc}
\usepackage{amssymb,graphicx,mathtools,amsfonts}
\usepackage{amsmath}
\usepackage{amsthm}
\usepackage{hyperref}
\usepackage[dvipsnames]{xcolor}
\usepackage{dsfont}
\renewcommand{\baselinestretch}{1.2}
\setlength\parindent{0pt} %for no indent everywhere
%%%%%%%%%% margins
%\setlength{\topmargin}{-.5cm}
%\setlength{\textheight}{24cm}
%\setlength{\evensidemargin}{1.2cm}
%\setlength{\oddsidemargin}{1.2cm}
%\setlength{\textwidth}{15.1cm}
%%%%%%%%%%%
\usepackage{a4wide}
\usepackage{esint}
\catcode`\@=11
\def\theequation{\@arabic{\c@section}.\@arabic{\c@equation}}
\catcode`\@=12
% \usepackage{refcheck} %% for checking references
\usepackage{enumitem}
\graphicspath{{images/}}
\newtheorem{theorem}{Theorem}[section] %The theorems will be numbered section-wise
\newtheorem{lemma}{Lemma}[section] %The lemmas will be numbered section-wise
\newtheorem{definition}{Definition}[section] %The definitions will be numbered section-wise
\newtheorem{proposition}{Proposition}[section]
\newtheorem{corollary}{Corollary}[section]
\newtheorem{inequality}{Inequality}[section]
\newtheorem{example}{Example}
% \theoremstyle{remark}
\newtheorem*{remark}{Remark}
\newtheorem{result}{Result}[section]
\newcommand{\R}{\mathbb{R}}

%Shortcut of some mathematical symbols
\newcommand{\ld} {\langle}
\newcommand{\rd} {\rangle}
\newcommand{\mc} {\mathcal}
\newcommand{\fa} {\forall}
\newcommand{\e}{\epsilon}
\newcommand{\pa} {\partial}
\newcommand{\al} {\alpha}
\newcommand{\ba} {s}
\newcommand{\de} {\delta}
\newcommand{\ga} {\gamma}
\newcommand{\Ga} {\Gamma}
\newcommand{\om} {\Omega}
\newcommand{\De} {\Delta}
\newcommand{\la} {\lambda}
\newcommand{\La} {\Lambda}
\newcommand{\noi} {\noindent}
\newcommand{\na} {\nabla}
\newcommand{\uline} {\underline}
\newcommand{\oline} {\overline}
\newcommand{\ds} {\displaystyle}
\newcommand{\ra} {\rightarrow}
\newcommand{\real}{\mathbb{R}}
\newcommand{\rn}{\mathbb{R}^{n}}
\newcommand{\rnn}{\mathbb{R}^{N}}
\newcommand{\rtwon}{\mathbb{R}^{2N}}
\newcommand{\lv}{\lVert}
\newcommand{\rv}{\rVert}
\newcommand{\weak}{\rightharpoonup}
\newcommand{\nlap}{\Delta_N}
\newcommand{\sobo}{W_0^{1,p}}
\newcommand{\lb}{\left(}
\newcommand{\rb}{\right)}
\newcommand{\grad}{\nabla}
\newcommand{\ntrl}{\mathbb{N}}
\newcommand{\ka}{\kappa}
\newcommand{\plap}{\Delta_p}
\newcommand{\pfrac}{(-\Delta_p)^s}
\newcommand{\rnnn}{\int_{\mathbb{R}^N}\int_{\mathbb{R}^N}}
\newcommand{\sobop}{W_0^{1,p}}
\newcommand{\fucik}{Fu\v{c}\'{\i}k } % for Fucik name
\newcommand{\X}{\mathcal{X}_0} % space short hand notation
\newcommand{\Xe}{\mathcal{X}_{0},\varepsilon} % space short hand notation
\newcommand{\J}{\mathcal{J}_{\mu,\lambda}} %functional short hand notation
\newcommand{\Je}{\mathcal{J}_{\mu,\lambda}^{\varepsilon}} %functional short hand notation
\newcommand{\I}{\mathcal{I}_{\mu,\lambda}} %functional with second solution
\newcommand{\Ie}{\mathcal{I}_{\mu,\lambda}^{\varepsilon}} %functional with second solution
\newcommand{\Jplus}{\mathcal{J}_{\mu,\lambda}^{+}} %functional short hand notation for positive solutions
\newcommand{\Jtilde}{\widetilde{\mathcal{J}}_{\mu, \la}}
\newcommand{\Jelo}{\mathcal{J}_{\mu, \lambda_0}^{\varepsilon}}
\newcommand{\Jep}{\mathcal{J}_{\mu, \lambda}^{\varepsilon,+}}
\newcommand{\Jepo}{\mathcal{J}_{\mu, \lambda_0}^{\varepsilon,+}}
\newcommand{\Q}{\mathcal{Q}_{\la,\mu}}

%useful commands
\allowdisplaybreaks

\title{Asymptotic behaviour and existence of positive solutions for mixed local nonlocal elliptic equations with Hardy potential}
\author{Shammi Malhotra\footnote{Department of Mathematics, Indian Institute of Technology Delhi, Hauz Khas New Delhi 110016,  India, shammi22malhotra@gmail.com}, Sarika Goyal\footnote{Department of Mathematics, Netaji Subhas University of Technology,
		Dwarka Sector-3, Dwarka, Delhi, 110078, India, sarika1.iitd@gmail.com, sarika@nsut.ac.in},  and K. Sreenadh\footnote{Department of Mathematics, Indian Institute of Technology Delhi, Hauz Khas New Delhi 110016,  India, sreenadh@maths.iitd.ac.in}}
\date{}

\begin{document}
\maketitle

\begin{abstract}
\noindent We investigate the existence and multiplicity of positive solutions to the following problem driven by the superposition of the Laplacian and the fractional Laplacian with Hardy potential
\begin{equation*}
% \label{eq:abs_main_equation}\tag{$\mathcal{P}_{\mu,\lambda}^{\varepsilon}$}
\left\{
\begin{aligned}
    -\Delta u + (-\Delta)^s u - \mu \frac{u}{|x|^2} &= \lambda |u|^{p-2} u + |u|^{2^*-2} u \quad \text{in } \Omega \subset \mathbb{R}^N, \\
    u &= 0 \quad \text{in } \mathbb{R}^N \setminus \Omega,
\end{aligned}
\right.
\end{equation*}
where \( \Omega \subset \mathbb{R}^N \) is a bounded domain with smooth boundary, \( 0 < s < 1 \), \( 1 < p < 2^* \), with \( 2^* = \frac{2N}{N-2} \), \( \lambda > 0 \), and \( \mu \in (0, \bar{\mu}) \) where $\bar \mu = \left( \frac{N-2}{2} \right)^2$. 

\noindent The aim of this paper is twofold. First, we establish uniform asymptotic estimates for solutions of the problem by means of a suitable transformation. Then, according to the value of the exponent \(p\), we analyze three distinct cases and prove the existence of a positive solution. Moreover, in the sublinear regime \(1 < p < 2\), we demonstrate the existence of multiple positive solutions for small perturbations of the fractional Laplacian.\\


\medskip
\noi\textbf{Keywords:} Mixed local nonlocal operator; Hardy potential; critical exponent; multiplicity of positive solutions; asymptotic estimates.  \\
\medskip
\noi\textbf{Mathematics Subject Classification:}  35A21, 35B09, 35B33, 35J20, 35M12.
\end{abstract}


\section{Introduction}
In this paper, we are concerned with the following problem
\begin{equation}\label{eq:main_problem} \tag{$\mathcal{P}_{\mu,\lambda}$}
\left\{
\begin{aligned}
    -\Delta u + (-\Delta)^s u - \mu \frac{u}{|x|^2} &= \lambda |u|^{p-2} u + |u|^{2^*-2} u \quad \text{in } \Omega \subset \mathbb{R}^N, \\
    u &= 0 \quad \text{in } \mathbb{R}^N \setminus \Omega,
\end{aligned}
\right.
\end{equation}
where $\om \subset \rnn$ is a bounded open set of class $C^{1,\al}$ for some $\al \in (0,1)$ such that $0 \in \om$, $\varepsilon \in (0,1]$, $0<s<1<p<2^*$ with $2^* = \frac{2N}{N-2}$, $\la >0$ is a parameter, $(-\Delta)^s$ is the fractional Laplacian operator defined as 
\begin{equation*}
    (-\Delta )^s u (x) = \text{P.V.} \int_{\rnn} \frac{u(x) - u(y)}{|x-y|^{N+2s}} dy, 
\end{equation*}
where P.V. is the Cauchy principle value and $\mu \in (0, \bar \mu)$ where $\bar \mu = \left( \frac{N-2}{2} \right)^2$ is the optimal constant in the Hardy's inequality \cite{adimurthi_hardy_inequality} which is given by
\begin{equation*}
     \int_{\om} \frac{|u|^2}{|x|^2}~ dx \leq \frac{1}{\bar \mu} \int_{\om} |\grad u |^2 ~ dx \qquad \text{for all } u \in C_c^{\infty}(\om).
\end{equation*}
The combination of local and nonlocal operators has recently emerged as a prominent area of research because of its wide-ranging and increasingly recognized applications. In fields such as finance and control theory, modelling often requires incorporating both diffusion and jump components in the underlying Markov processes. This dual nature introduces significant challenges, as the process operates on two distinct scales: the diffusion component dominates at small scales, while the jump component becomes more influential at larger scales. Such operators also play a crucial role in population dynamics. A rigorous mathematical treatment of these models can be found in~\cite{valdinoci_mixed_application}. 
% Moreover, in the presence of strong magnetic fields, mixed operators are used to analyze anisotropic heat transfer, as discussed in~\cite{blazevski_anisotropic_heat_transport}.

The analysis of PDEs involving singular potentials is of intrinsic interest. In quantum mechanics, for instance, the Hardy potential characterizes motion and interactive properties such as repulsion and attraction between charged particles (see~\cite{frank_hardy_applications}). For a comprehensive treatment of the Hardy potential and its wide-ranging applications, we refer the interested reader to the monograph~\cite{peral_hardy_applications}. It is worth noting that in the context of mixed operators, the Hardy potential is typically chosen to coincide with that associated with the local operator. This choice is motivated by the fact that, for the mixed operator, all potentials of the form \( |x|^{-t} \) with \( t \in [2s, 2] \) are admissible; see Lemma~3.1 of~\cite{malhotra2025eigenvalues}. However, among these, the local Hardy potential exhibits the strongest singularity, making it the most natural candidate for perturbative analysis. A rigorous discussion on this can be found in~\cite{biagi2024mixed_hardy}.


Motivated by the foundational role of the Hardy potential, the study of Brezis-Nirenberg-type problems for Laplacian operator with Hardy potential was initiated by Jannelli~\cite{jannelli_hardy_starting}, who investigated the following problem:
\begin{equation}\label{eq:local_problem} 
\left\{
\begin{aligned}
    -\Delta u  - \mu \frac{u}{|x|^2} &= \lambda  u + |u|^{2^*-2} u \quad \text{in } \Omega \subset \mathbb{R}^N, \\
    u &= 0 \quad \text{on } \partial\Omega.
\end{aligned}
\right.
\end{equation}
The author proved that if \( 0 < \mu \leq \bar{\mu} - 1 \), then problem~\eqref{eq:local_problem} admits a positive solution for all \( \lambda \in (0, \lambda_1) \). In contrast, if \( \bar{\mu} - 1 < \mu < \bar{\mu} \) and \( \Omega = B_1(0) \) is the unit ball, there exists a threshold \( \lambda_* \) such that problem~\eqref{eq:local_problem} admits a positive solution if and only if \( \lambda \in (\lambda_*, \lambda_1) \), where \( \lambda_1 \) denotes the first eigenvalue of the Laplacian operator with the Hardy potential. This result highlights that every dimension can become critical when \( \mu \) is close to \( \bar{\mu} \), specifically when \( \mu \in (\bar{\mu} - 1, \bar{\mu}) \). This stands in stark contrast to the classical Brezis-Nirenberg problem with \( \mu = 0 \) (see \cite{brezis_1983}), where only \( N = 3 \) is critical. This phenomenon is explained by the guiding principle in~\cite{jannelli_hardy_starting}, which asserts that a spatial dimension is critical for a linear elliptic operator \( \mathcal{L} \) if and only if \( \mathcal{L} \) admits at least one Green function \( G(x_0, x) \in L^2_{\text{loc}}(\mathbb{R}^N) \). Subsequently, Cao and Peng~\cite{cao_peng2003signchanging} established the existence of sign-changing solutions to problem~\eqref{eq:local_problem} for all \( \lambda \in (0, \lambda_1) \). 
% They constructed a sequence of sign-changing solutions to subcritical problems approximating the critical exponent problem~\eqref{eq:local_problem}, and using a compactness argument, showed that a subsequence converges to a sign-changing solution of the original problem.
In a related direction, Ferrero and Gazzola~\cite{ferrero2001existence} replaced the linear perturbation term \( \lambda u \) with a more general subcritical nonlinearity \( g(x,u) \) and established the existence of solutions under suitable conditions. For further results in the local case, we refer the reader to~\cite{pigong_local_p_hardy} and references therein. \\
Transitioning to the purely nonlocal case, various studies have been conducted to understand perturbations of the Hardy potential in the context of the fractional Laplacian. Dipierro et al.~\cite{dipierro2016fractionalwithHardy} proved the existence of extremals for the fractional Hardy-Sobolev inequality and investigated several of their qualitative properties. In particular, they derived asymptotic estimates via a suitable change of variables and analysis of the resulting transformed equation. Ghoussoub et al.~\cite{ghoussoub_fractional_hardy} considered the following problem involving the fractional Hardy-Schrödinger operator:
\begin{equation}\label{eq:nonlocal_problem}
\left\{
\begin{aligned}
    (-\Delta)^s u  - \mu \frac{u}{|x|^{2s}} &= \lambda u + \frac{|u|^{2^*_s(\alpha)-2} u}{|x|^{\alpha}} \quad \text{in } \Omega \subset \mathbb{R}^N, \\
    u &= 0 \quad \text{on } \mathbb{R}^N \setminus \Omega,
\end{aligned}
\right.
\end{equation}
where \( \mu \in [0, \Lambda_{N,s}] \), \( \lambda \in (0, \lambda_s(\mu)) \) with \( \lambda_{s,\mu} \) denoting the first eigenvalue of the Hardy-Schrödinger operator, \( 0 \leq \alpha < 2s < N \), and \( 2^*_s(\alpha) = \frac{2(N - \alpha)}{N - 2s} \). By introducing the concept of the \emph{internal mass} of the domain, they addressed the critical case and established the existence of least energy solutions to problem~\eqref{eq:nonlocal_problem} under various conditions on the Hardy term. Later on, Shang et al.~\cite{shang_zhang_fractional_laplacian} incorporated weights into the nonlinear terms and obtained results concerning the existence and multiplicity of solutions.\\
Unlike purely local or purely nonlocal cases, the study of mixed local–nonlocal operators involving the Hardy potential is still in the developmental stage. Biagi et al.~\cite{biagi2024mixed_hardy} investigated the existence, uniqueness, and optimal summability of solutions. Malhotra~\cite{malhotra2025eigenvalues} studied the Fučík spectrum and shape optimization problems for the first two eigenvalues of such mixed operators. 

On the other hand, in the absence of the Hardy potential, Brezis–Nirenberg-type results for the mixed linear operator were first established in~\cite{biagi2022brezis} and subsequently generalized to mixed quasilinear operator by Silva et al. in~\cite{silva2024mixed}. 
% To further understand the interaction between the classical Laplacian and the fractional Laplacian, the profile decomposition of Palais–Smale sequences associated with the mixed operator was also analyzed in~\cite{chakraborty2025global}.

Inspired by the aforementioned work, we started our study with the analysis of the following problem
\begin{equation}\label{eq:main_eq}
\left\{
\begin{aligned}
    -\Delta u + \varepsilon(-\Delta)^s u - \mu \frac{u}{|x|^2} &=|u|^{2^*-2} u \quad \text{in } \Omega \subset \mathbb{R}^N, \\
    u &= 0 \quad \text{in } \mathbb{R}^N \setminus \Omega.
\end{aligned}
\right.
\end{equation}
The solutions of the equation \eqref{eq:main_eq} are intimately connected with the minimizers of the following ratio
\begin{equation} \label{eq:spectral_ratio_frac}
S_{\mu, s,\varepsilon}(\Omega) := \inf_{u \neq 0} 
\frac{\displaystyle\int_{\Omega} |\nabla u|^2 - \mu \int_{\Omega} \frac{|u|^2}{|x|^2} + \varepsilon[u]_{s}^2}
{\left(\displaystyle\int_{\Omega} |u|^{2^*} \right)^{\frac{2}{2^*}}}.
\end{equation}
Thus, the study of the problem \eqref{eq:main_eq} is completed once we find the minimizers. Unlike the purely local \cite{terracini_minimizers} or nonlocal case \cite{cotsiolis2004best_fractional}, the minimizers for combination of local and nonlocal terms does not exist due to the lack of common scaling invariance. 
%However, unlike the purely local or nonlocal case where such minimizers exists \cite{terracini_minimizers}, the presence of both terms exhibits lack of scaling invariance that leads to lack of existence of minimizers. 

This can be seen in the following theorem
\begin{theorem} Let $s \in (0,1)$, $\varepsilon \in (0,1]$ be fixed and let $\Omega \subseteq \mathbb{R}^N$ be an open set. Then
\begin{equation}\label{eq:best_constant_equality}
S_{\mu, s, \varepsilon}(\Omega) = S_{\mu}(\Omega) = S_{\mu},
\end{equation}
where
\begin{equation*} 
S_{\mu}(\Omega) := \inf_{u \neq 0} 
\frac{\displaystyle\int_{\Omega} |\nabla u|^2 - \mu \int_{\Omega} \frac{|u|^2}{|x|^2}}
{\left(\displaystyle\int_{\Omega} |u|^{2^*} \right)^{\frac{2}{2^*}}}.
\end{equation*}
Moreover, the optimal constant $S_{\mu, s, \varepsilon}(\Omega)$ in equation \eqref{eq:spectral_ratio_frac} is never attained and independent of $\om$.
\end{theorem}
% We note that the quantity $S_{\mu}(\Omega)$ is independent of $\Omega$. 
The proof of the above theorem follows in the same spirit of Theorem $1.1$ and $1.2$ of \cite[Theorem $1.1$]{biagi2024mixed_hardy}. \\
In fact, a complete classification of positive solutions of the minimization problem for $S_{\mu}(\om)$ is done in \cite{olivia_classfication_localp} via the moving plane method. They proved that all the solutions are radial and radially decreasing about the origin.\\
Since the best constant in the minimization problem \eqref{eq:spectral_ratio_frac} is not achieved. The corresponding problem \eqref{eq:main_eq} does not possesses the groundstate solutions. This motivates the study of the following perturbed problem \eqref{eq:main_problem_ep}
\begin{equation}\label{eq:main_problem_ep} \tag{$\mathcal{P}_{\mu,\lambda}^{\varepsilon}$}
\left\{
\begin{aligned}
    -\Delta u + \varepsilon(-\Delta)^s u - \mu \frac{u}{|x|^2} &= \lambda |u|^{p-2} u + |u|^{2^*-2} u \quad \text{in } \Omega \subset \mathbb{R}^N, \\
    u &= 0 \quad \text{in } \mathbb{R}^N \setminus \Omega,
\end{aligned}
\right.
\end{equation}
where $\varepsilon \in (0,1]$. \\
% However, the study of the perturbed problem greatly varies with the nature of the nonlinearity involved and we study each of the cases separately.\\
We emphasize that introducing the parameter \( \varepsilon \) in front of the fractional Laplacian term in problem~\eqref{eq:main_problem_ep} is essential (at least in the sublinear case) in order to establish the existence of solutions via the standard variational approach, as adapted from~\cite{brezis_1983}. The key idea is to convert the existence question to the analysis of the following minimization problem
\begin{equation*} 
S_{\mu, s,\lambda}(\Omega) := \inf_{u \neq 0} 
\frac{\displaystyle\int_{\Omega} |\nabla u|^2 \, dx - \mu \int_{\Omega} \frac{|u|^2}{|x|^2} \, dx +  [u]_{s}^2 - \lambda \int_{\Omega} |u|^2 \, dx}
{\left(\displaystyle\int_{\Omega} |u|^{2^*} \, dx \right)^{\frac{2}{2^*}}}.
\end{equation*}
A solution to problem~\eqref{eq:main_problem} exists provided that $S_{\mu, s,\lambda}(\Omega) < S_{\mu, s}$, where $ S_{\mu, s}:=S_{\mu, s,1}$. The standard method to verify this inequality involves selecting a suitable test function and evaluating the quotient to show it lies strictly below the threshold. Since \( S_{\mu, s} = S_\mu \), a natural choice is to choose a minimizer for $S_{\mu}$ which is attained by a family of functions \( U_{\varepsilon} \)(defined explicity later). However, one encounters the following asymptotic estimates
\begin{equation*}
    [U_{\varepsilon}]_s^2 = O\left(\varepsilon^{2(1-s)\frac{\sqrt{\bar\mu}}{\sqrt{\bar\mu - \mu}}}\right), \quad 
    \int_{\Omega} |U_{\varepsilon}|^2 \, dx = O\left(\varepsilon^{2\frac{\sqrt{\bar\mu}}{\sqrt{\bar\mu - \mu}}}\right) 
    \quad \text{as } \varepsilon \to 0^+.
\end{equation*}
These estimates reveal that the contribution of the Gagliardo seminorm $[U_{\varepsilon}]_s$ (defined in Section \ref{sec:notations}) becomes non-negligible as \( \varepsilon \to 0^+ \), thereby obstructing the inequality \( S_{\mu, s,\lambda} < S_{\mu, s} \) and making the existence of a minimizer via this approach appear infeasible.


The aim of this paper is twofold. Firstly, the equality in \eqref{eq:best_constant_equality} naturally raises the question of the nature of minimizers of \(S_{\mu, s, \varepsilon}(\Omega)\) with a perturbation, and how they differ from their local counterparts. This motivates the study of the exact asymptotic behavior of solutions to \eqref{eq:main_problem_ep} near the origin.\\
To establish the asymptotic estimates, the key idea is to transform the original problem \eqref{eq:main_problem_ep} into the reformulated problem \eqref{eq:transformed_equation}(defined later), which involves working within radial Sobolev spaces. For this transformed problem, a Harnack inequality and uniform estimates are derived, which play a crucial role in obtaining the lower and upper asymptotics, respectively.

\begin{theorem}\label{thm:asympototic_estimates}
Let $0 < \mu < \bar\mu$ and $1<p<2^*$. Then for any weak solution $u$ of \eqref{eq:main_problem_ep} there admits two positive constants $M_1$ and $M_2$ independent of $\varepsilon$ such that
\begin{equation*}
M_1 |x|^{- \left( \sqrt{\bar{\mu}} - \sqrt{\bar{\mu} - \mu} \right)} \leq u(x) \leq M_2 |x|^{- \left( \sqrt{\bar{\mu}} - \sqrt{\bar{\mu} - \mu} \right)} \quad \text{for all } x \in B_{r_0}(0) \subset \Omega,
\end{equation*}
with some $r_0 > 0$ sufficiently small.
\end{theorem}
% \textcolor{red}{check if the condition $0< p < 1$ can be relaxed for the lower bound and upper bound}
% These asymptotic estimates can be seen as generalisation of asymptotic estimates proved in \cite{terracini_minimizers} for the local case and \cite[Theorem $1.7$]{dipierro2016fractionalwithHardy} in the nonlocal case. 



% Multiplicity results like these are proved in the Laplacian and fractional Laplacian case in \cite{chen2004multiple} and \cite{shen_multiplicity_fractional} respectively.\\

% If we dropped the requirement for solutions to be non-negative we can prove the existence of infinitely many solutions with the help of topological arguments based on Krasnoselskii genus theory. Without the Hardy potential, this type of result was proved in \cite[Theorem $1.1$]{silva2024mixed}.
% \begin{theorem}\label{thm:sublinear_infinite_solutions}
%  Let $0<s<1$, $0<q<1$. There exists $\Lambda > 0$ such that for any $\la \in (0,\La)$ the problem \eqref{eq:main_problem_ep} admits infinitely many solutions with negative energy.
% \end{theorem}

Secondly, we turn our analysis to the effect of perturbations on the existence of solutions. We begin with the case of a linear perturbation, which requires knowledge of the first eigenvalues of the fractional operator and the mixed operator with Hardy potential, defined respectively as
\begin{equation} \label{eq:first_eigenvalue_frac}
   \lambda_{1,s} := 
  \inf\big\{[u]^2_s:\,\text{$u\in C_c^\infty(\Omega)$ and $\|u\|_{L^2(\rnn)} = 1$}\big\};
\end{equation}
\begin{equation} \label{eq:first_eigenvalue}
\lambda_{1} := 
  \inf\left\{ \int_{\om} |\grad u|^2 dx + [u]_s^2 - \mu \int_{\om}\frac{|u|^2}{|x|^2}dx:\,\text{$u\in C_c^\infty(\Omega)$ and $\|u\|_{L^2(\rnn)} = 1$}\right\}.
\end{equation}
For a comprehensive results related with these eigenvalues we refer to \cite[Remark $4.4$]{biagi2022brezis}. Then we have the following theorem concerning the solutions of the linear problem \eqref{eq:main_problem}
\begin{theorem}\label{thm:existence_linear_case}
Let $\Omega \subseteq \mathbb{R}^N$, $\varepsilon = 1$ and $p = 2$. Then the following holds
\begin{enumerate}
    \item For every $0 < \lambda \leq \lambda_{1, s}$, there does not exists a solution $u \in \mathcal{B} \subset L^{2^*}(\mathbb{R}^N)$ of \eqref{eq:main_problem}, where
    \begin{equation}\label{eq:linear_ball_NE}
    \mathcal{B} := \left\{ u \in L^{2^*}(\mathbb{R}^N) : \| u \|_{L^{2^*}} \leq  S_{\mu}^{\frac{N-2}{4}} \right\}.
    \end{equation}
    \item There exists a parameter $\lambda^* \in [\lambda_{1, s}, \lambda_1)$ such that the problem \eqref{eq:main_problem} possesses at least one solution if $\lambda \in (\lambda^*, \lambda_1)$. 
    \item There does not exist a solution to the problem \eqref{eq:main_problem} if $\lambda > \lambda_1$.
\end{enumerate}
\end{theorem}

% Even though there are no groundstate solutions of the linear problem for $\la > \la_1$, we can look for high energy solutions. With the help of the linking theorem \cite[Theorem $2.5$]{cerami_multiplicity}, the existence of high energy solutions are proved for suitable values of $\la$.
% \begin{theorem}\label{thm:high_energy_linear}
% Let $\varepsilon =1$ and $p=2$. Also, assume that $\mu \in \left( \bar{\mu} - 1, \bar{\mu} \right)$, and there exists $\lambda_k \in \sigma_\mu$ such that
% \begin{equation} \label{eq:lambda-range}
% \lambda \in \left( \lambda_k - S_{\mu} |\Omega|^{-2/N}, \, \lambda_k \right),
% \end{equation}
% where $\sigma_\mu$ denotes the spectrum associated with the operator $ -\Delta  + (-\Delta)^s  - \mu \frac{1}{|x|^2}$, and $S_{\mu}$ is the best constant as in Theorem \ref{thm:optimal_frac_constant}.\\
% Then, the corresponding linear problem admits $\eta_k$ pairs of nontrivial solutions, where $\eta_k$ is the multiplicity of the eigenvalue $\lambda_k$.
% \end{theorem}
For the superlinear case, i.e., $p \in (2,2^*)$, we are interested in the existence of a positive solution of the problem \eqref{eq:main_problem}. We apply the mountain pass theorem to prove the existence of a solution. However, the presence of nonlocal term added another realm of difficulties. A perturbation close to linear power is not sufficient to guarantee the existence of the solution. In fact we require a perturbation with higher exponent to tackle the presence of fractional term. For this purpose we define the following crucial parameter
\begin{equation}\label{eq:superlinear_parameter}
    \beta_{\mu,N,s} = \min \left\{ N-2, \frac{2(1-s)\sqrt{\bar\mu}}{\sqrt{\bar\mu - \mu}} \right\}.
\end{equation}
Thus, we have the following existence result.
\begin{theorem}\label{thm:superlinear_pb}
Let $p \in (2,2^*)$. Then there exists a $\la_0 >0$ such that the problem \eqref{eq:main_problem} has a positive solution in the following cases
\begin{enumerate}
    \item $ 0 < \la < \la_0$ and $\beta_{\mu,N,s} > \frac{\sqrt{\bar\mu}(N - p \sqrt{\bar\mu})}{\sqrt{\bar\mu - \mu}}$,
    \item $\la \geq \la_0$,
\end{enumerate}
where $\beta_{\mu,N,s}$ is defined as in \eqref{eq:superlinear_parameter}.
\end{theorem}
As we note that the first case in the above only occurs when the exponent $p$ is sufficiently large enough. For perturbations with lower order exponent, the existence of solutions are proved only for large values of $\la$.\\

We end our study by adding the perturbation of sublinear nature, i.e., $p \in (1,2)$. This nonlinearity greatly influences the topology of the associated functional of the problem \eqref{eq:main_problem_ep} and leads to the existence of two nonnegative multiple solutions in the spirit of \cite{abc_laplacian}. However, the solutions are not expected to be bounded in a small neighbourhood of zero. This requires us to work with minimal solutions. Once the minimal solutions are obtained, we are able to prove the following multiplicity result with the help of the Mountain pass theorem. 
\begin{theorem}\label{thm:sublinear_two_solutions}
Suppose \(0 < \mu < \bar{\mu} - 1\). Then there exists \(\Lambda > 0\) such that  
\begin{enumerate}
    \item for all \(\lambda \in (0, \Lambda)\), the problem \eqref{eq:main_problem_ep} admits a minimal solution \(u_{\lambda}\), and these minimal solutions are increasing with respect to \(\lambda\);  
    \item for \(\lambda = \Lambda\), the problem \eqref{eq:main_problem_ep} admits at least one weak solution;  
    \item for \(\lambda > \Lambda\), the problem \eqref{eq:main_problem_ep} does not admit any solution.  
\end{enumerate}
Moreover, if \(\lambda \in (0, \Lambda)\), the problem admits a second positive solution distinct from the first one.
\end{theorem}

One of the main novelties of this paper is the derivation of blow-up estimates for the Gagliardo norm of minimizers of \(S_{\mu}\). These estimates are obtained by employing certain inequalities from \cite{giovanni_fractional_book}. Another difficulty arises in establishing the Harnack inequality, where obtaining estimates for the Hardy term proves challenging; this issue is addressed by applying a crucial lemma from \cite{stein_weiss1957fractional}. Finally, the presence of the Hardy term results in the lack of equivalence of minimizers in the \(C^1\) topology, which necessitates a delicate analysis in order to prove the existence of local minima of solutions in the sublinear case.



The paper is organized as follows. In Section~\ref{sec:notations}, we begin by introducing the basic notations and preliminaries required throughout the paper. In addition, several auxiliary lemmas are stated and proved to support the main results. Then we derive asymptotic estimates for solutions of problem~\eqref{eq:main_problem_ep} by employing a suitable transformation in Section~\ref{sec:asymptotics} and end it by proving a strong maximum principle. 
% This section concludes with a proof of the existence of infinitely many solutions.
In Section~\ref{sec:linear_sublinear}, we start with examining the case of linear perturbations and establish both existence and nonexistence results, depending on the value of the parameter relative to the first eigenvalue of the underlying operator. Then we consider the case of superlinear perturbations. Applying the Mountain Pass Theorem, we show the existence of solutions under appropriate conditions on the exponent of the nonlinearity. Finally the section~\ref{sec:sublinear} is devoted to the analysis of sublinear perturbations, where the existence of two positive solutions is established. 



%%%%%%%%%%%%%%%%%%%%%%%%%%%%%%%%%%%%%%%%%%%%%%%%%%%%%%%%%%%%%%%%%%
%%%%%%%%%%%%%%%%   Notations and Preliminaries          %%%%%%%%%%%%%%%%%
%%%%%%%%%%%%%%%%%%%%%%%%%%%%%%%%%%%%%%%%%%%%%%%%%%%%%%%%%%%%%%%%%%
\section{Notations and Preliminaries}\label{sec:notations}
%%%% Functional Formulation}

In this section we setup the function spaces which are used to study our problem \eqref{eq:main_problem_ep} and state some lemmas and prove some propositions required in our analysis. First of all, let us fix some notations
\begin{enumerate}
    \item[$(i)$] The constants will be denoted by $C$ and they are allowed to vary within a single line or formula. 
    \item[$(ii)$] If $A$ is any measurable set in $\rnn$, then $|A|$ denotes the $N$-dimensional Lebesgue measure.
    \item[$(iii)$] $B_R(x)$ denotes the $N$-dimensional ball in $\rnn$ centered at $x$ having radius $R$. If it is centered at $0$, then we denote it by $B_R$.
    \item[$(iv)$] For $1 \leq p < \infty$, we denote and define the $L^p(\om)$ norm as $\lv u \rv_{L^p(\om)}^p := \int_{\om} |u|^p dx$.
    \item[$(v)$] $C_c^\infty(\om):= \{ u \in C^{\infty}(\om) : Supp(u) \text{ is compact in } \om \}$ and $C_{c,rad}^\infty(\om) := \{ u \in C_c^\infty(\om) : u \text{ is radial}\}$.
\end{enumerate}

Now, let $\om \subset \rnn$ be a bounded domain and $\varepsilon \in (0,1]$. Then we define the following function space
\begin{equation*}
\X(\Omega) := \left\{ u \in H^1(\mathbb{R}^N) \,:\, u_{|\Omega} \in H_0^1(\Omega) \text{ and } u \equiv 0 \text{ a.e. in } \mathbb{R}^N \setminus \Omega \right\}.
\end{equation*}
It is a Hilbert space endowed with the norm 
\begin{equation*}
    \lv u \rv_{\X,\varepsilon} := \left( \int_{\om} |\grad u|^2 dx +  \varepsilon [ u ]_s^2\right)^{\frac{1}{2}},
\end{equation*}
where $[u]_s^2 = \iint_{\rtwon} \frac{|u(x) - u(y)|^2}{|x-y|^{N +2s}}dy dx$ is known as Gagliardo seminorm induced from the inner product $\langle u, v\rangle_s = \iint_{\rtwon} \frac{(u(x) - u(y))(v(x) - v(y))}{|x - y|^{N  + 2s}}dy dx$. We note that $\lv \cdot \rv_{\X,\varepsilon}$ is equivalent to the usual norm $\lv \cdot \rv_{\X} \left( := \lv \cdot \rv_{\X,1} \right)$. Moreover, $\X(\om)$ is reflexive and separable with respect to the norm $\lv \cdot \rv_{\X,\varepsilon}$. \\
A function $u \in \X(\om)$ is said to be a weak solution of the problem \eqref{eq:main_problem_ep} if for every $v \in \X(\om)$, it holds 
\begin{equation*}
    \int_{\om} \grad u \grad v \, dx + \varepsilon \langle u , v\rangle_s - \mu \int_{\om} \frac{u v}{|x|^2}dx = \int_{\om} |u|^{2^* -2}u v dx + \int_{\om} |u|^{p-2}u v dx.
\end{equation*}
In addition, if the equality is replaced with $\geq(\leq)$ then it is known as weak supersolution(subsolution). 
The weak solution to problem \eqref{eq:main_problem_ep} corresponds to a critical point of the energy functional $\Je: \X(\Omega) \to \mathbb{R}$ given by
\begin{equation*}
\Je(u) = \frac{1}{2} \int_{\Omega} |\nabla u|^2 - \frac{\mu}{2} \int_{\Omega} \frac{|u|^2}{|x|^2} dx
+ \frac{\varepsilon}{2} [u]_{s}^2 
- \frac{1}{2^*} \int_{\Omega} |u|^{2^*} dx
- \frac{\lambda}{p} \int_{\Omega} |u|^{p}dx.
\end{equation*}
It is clear that $\Je \in C^1(\X,\real)$. 

\begin{definition}
    A sequence $\{ u_n \} \subset \X(\Omega)$ is said to be a Palais-Smale $(PS)$ sequence for a functional $\mathcal{J}$ at level $\beta$, if $\mathcal{J}(u_n) \to \beta$ in $\R$ and $\mathcal{J}^{\prime}(u_n) \to 0$ in $(\X(\Omega))'$ as $n \to \infty$. The function $\mathcal{J}$ is said to satisfy $(PS)$ condition at level $\beta$, if every $(PS)$ sequence for $\mathcal{J}$ at level $\beta$ admits a convergent subsequence.
\end{definition}


Also, we would like to mention a very crucial lemma from \cite{stein_weiss1957fractional}.
\begin{lemma}[Lemma $3.5$ of \cite{stein_weiss1957fractional}]\label{lem:stein_strong_lemma}
Let 
\begin{equation*}
    V_{\de}(f) = |x|^{-N + \delta} \int_{\{|y| < |x|\}} |y|^{- \delta} f(y) \, dy
\end{equation*}
where $1< q < \infty$ and $\de < \frac{N(q-1)}{q}$. Then $|V_{\de}(f)| \leq C |x|^{-\frac{N}{q}} \lv f \rv_{L^q}$ for some $C>0$.
\end{lemma}

Finally, we end this section by mentioning some useful inequalities borrowed from \cite{brasco_second_eigenvalue}
\begin{lemma}\label{ineq:A1}
    Let $f$ be a convex function, then 
    \begin{equation*}
        (a-b)[A f^{\prime}(a) - B f^{\prime}(b)] \geq (f(a) - f(b))(A - B)
    \end{equation*}
for all $a, b \in \real$ and $A, B \geq 0$.
\end{lemma}

\begin{lemma}\label{ineq:A2}
Let $g : \real \to \real $ be an increasing function and $G(t) = \int_0^t g^{\prime} (\tau)^{\frac{1}{2}} d\tau$. Then
\begin{equation*}
    (a - b)(g(a) - g(b)) \geq |G(a) - G(b)|^2.
\end{equation*}
\end{lemma}

\begin{lemma}\label{ineq:A3}
Let $g : \real \to \real $ be an decreasing function and $H(t) = -\int_0^t g^{\prime} (\tau)^{\frac{1}{2}} d\tau$. Then
\begin{equation*}
    (a - b)(g(b) - g(a)) \geq |H(a) - H(b)|^2.
\end{equation*}
\end{lemma}
%%%%%%%%%%%%%%%%{Boundedness and Positivity Result}


%%%%%%%%%%%%%%%%%%%%%%%%%%%%%%%%%%%%%%%%%%%%%%%%%%%%%%%%%%%%%%%%%%
%%%%%%%%%%%%%%%%    Asymptotic Estimates         %%%%%%%%%%%%%%%%%
%%%%%%%%%%%%%%%%%%%%%%%%%%%%%%%%%%%%%%%%%%%%%%%%%%%%%%%%%%%%%%%%%%
\section{Qualitative Properties}\label{sec:asymptotics}

In this section we study the behaviour of solutions of \eqref{eq:main_problem_ep} near the origin. We start with transforming the problem \eqref{eq:main_problem_ep} into a suitable problem. After that Harnack inequality and uniform estimates are obtained. Finally, we end the section by proving a strong maximum principle.\\
We start this section by introducing the Sobolev inequality and Poincar\'{e} inequality in the radial case.


% Moreover we will be working with radial weights to get asymptotic estimates, for that we require Sobolev inequality and Poincar\'{e} inequality. 
\begin{lemma}\label{lem:weighted_sobolev_inequality}
Let $d \in  (0, \frac{N-2}{2})$. Then there exists a positive constant $C>0$ such that for all $u \in C_c^{\infty}(B_R)$ the following inequality holds
\begin{equation*}
        \left( \frac{1}{|B_R|_{d\mu}}\int_{B_R} \frac{|u|^q}{|x|^{2d}} dx\right)^{\frac{1}{q}} \leq C \left( \frac{R^2}{|B_R|_{d \mu}} \int_{B_R} \frac{|\grad u |^2}{|x|^{2d}}dx\right)^{\frac{1}{2}},
\end{equation*}
where $q = \frac{2 (N- 2d)}{N - 2d - 2}$ and $|B_R|_{d \mu} = \int_{B_R} \frac{dx}{|x|^{2d}}$.
\end{lemma}
We note that the above inequality also holds for lower exponents \( 1 \leq q_1 < q \), which is a direct consequence of Jensen's inequality. \\
% \begin{proof}
% This can be deduced from by taking following choice of parameters in \cite{caffarelli1984interpolation_sobolev}
% \begin{equation*}
%     \sigma =  \gamma = \beta = \frac{- \delta(N- \delta-2)}{2(N- \delta)}, \al = - \frac{\delta}{2}, r = q = \frac{2(N- \delta)}{N- \delta-2}, a = \frac{1}{2}, \text{ and } p=2,
% \end{equation*}
% for any $\delta \in \left(0, \frac{N-2}{2}\right)$.
% \end{proof}
To prove the above inequality we require the following lemma.
\begin{lemma}\label{lem:radial_estimate}
For all $u \in C_{c, \mathrm{rad}}^{\infty}(B_1)$, the following holds:
\begin{enumerate}
    \item[$(i)$]  $
|u(x)| \;\leq\; C \left( \int_{B_1} \frac{|\nabla u|^2}{|y|^{2d}} \, dy \right)^{\tfrac{1}{2}}
\, r^{-\tfrac{(N-2d-2)}{2}},$
\item[$(ii)$] $
\left( \int_{B_1} \frac{|u|^q}{|x|^{2d}} \, dx \right)^{1/q} 
\;\leq\; C \left( \int_{B_1} \frac{|\nabla u|^2}{|x|^{2d}} \, dx \right)^{1/2},$
\end{enumerate}
where $d = \sqrt{\bar{\mu}} - \sqrt{\bar{\mu} - \mu}, \, r = |x| < 1$ and $q = \frac{2(N-2d)}{N-2d-2}$.
\end{lemma}

\begin{proof}
Since $u \in C_{c,\mathrm{rad}}^{\infty}(B_1)$, we have
\begin{equation*}
-u(x) = - u(r) = u(1) - u(r) = \int_{r}^{1} u'(s) \, ds .
\end{equation*}
This together with the H\"{o}lder inequality yields
\begin{equation*}
\begin{aligned}
|u(x)| 
&\leq \int_{r}^{1} |u'(s)| \, ds \\
&\leq \left( \int_{r}^{1} |u'(s)|^2 \, s^{N-2d-1} \, ds \right)^{\tfrac{1}{2}}
\left( \int_{r}^{1} s^{-(N-2d-1)} \, ds \right)^{\tfrac{1}{2}} \\
& \leq \omega_N^{-1/2} 
\left( \int_{B_1} \frac{|\nabla u|^2}{|y|^{2d}} \, dy \right)^{\tfrac{1}{2}}
\left( \frac{r^{-(N-2-2d)}}{N-2-2d} \right)^{\tfrac{1}{2}}\\
& =\; C 
\left( \int_{B_1} \frac{|\nabla u|^2}{|y|^{2d}} \, dy \right)^{\tfrac{1}{2}} 
r^{-\tfrac{(N-2-2d)}{2}} .
\end{aligned}
% \label{eq:radial_splitting}
\end{equation*}
For the second part, we apply Cauchy-Schwarz inequality and part $(i)$ of Lemma~\ref{lem:radial_estimate} to obtain
\begin{equation*}
\begin{aligned}
\int_{B_1} \frac{|u|^q}{|x|^{2d}} \, dx 
&= \omega_N \int_{0}^{1} |u(r)|^q \, r^{N-1-2d} \, dr \\
&= -\omega_N \int_{0}^{1} q |u(r)|^{q-2} u(r) u'(r) \, \frac{r^{N-2d}}{N-2d} \, dr \\
&\leq \frac{q \, \omega_N}{N-2d} 
\int_{0}^{1} \big( |u'(r)| r^{\tfrac{N-2d-1}{2}} \big)
\big( |u(r)|^{q-1} r^{\tfrac{N-2d+1}{2}} \big) \, dr \\
& \leq \frac{q \, \omega_N}{N-2d} 
\left( \int_{0}^{1} |u'(r)|^2 r^{N-2d-1} \, dr \right)^{1/2}
\left( \int_{0}^{1} |u(r)|^{2(q-1)} r^{(N-2d+1)} \, dr \right)^{1/2}\\
& \leq \frac{C \, q }{N-2d} 
\left( \int_{B_1} \frac{|\nabla u|^2}{|x|^{2d}} \, dx \right)^{q/4}
\left( \int_{B_1} \frac{|u|^q}{|x|^{2d}} \, dx \right)^{1/2}.
\end{aligned}
% \label{eq:radial_q_integral}
\end{equation*}
Rearranging terms, we finally deduce
\begin{equation*}
\left( \int_{B_1} \frac{|u|^q}{|x|^{2d}} \, dx \right)^{1/q}
\;\leq\; \left( \frac{C \, q}{N-2d} \right)^{2/q}
\left( \int_{B_1} \frac{|\nabla u|^2}{|x|^{2d}} \, dx \right)^{1/2}. \qedhere
% \label{eq:final_weighted_sobolev}
\end{equation*}

\end{proof}

% \begin{lemma}\label{lem:radial_sobolev}
% Let $q = \frac{2(N-2d)}{N-2d-2}, d = \sqrt{\bar{\mu}} - \sqrt{\bar{\mu} - \mu}$. Then there exists a constant $C > 0$ such that
% \begin{equation}
% \left( \int_{B_1} \frac{|u|^q}{|x|^{2d}} \, dx \right)^{1/q} 
% \;\leq\; C \left( \int_{B_1} \frac{|\nabla u|^2}{|x|^{2d}} \, dx \right)^{1/2}.
% \label{eq:weighted_sobolev_inequality}
% \end{equation}
% for all $u \in C_{c,\mathrm{rad}}^{\infty}(B_1)$.
% \end{lemma}
% \begin{proof}
% By applying Cauchy-Schwarz inequality and Lemma~\ref{lem:radial_estimate}, we obtain
% \begin{equation*}
% \begin{aligned}
% \int_{B_1} \frac{|u|^q}{|x|^{2d}} \, dx 
% &= \omega_N \int_{0}^{1} |u(r)|^q \, r^{N-1-2d} \, dr \\
% &= -\omega_N \int_{0}^{1} q |u(r)|^{q-2} u(r) u'(r) \, \frac{r^{N-2d}}{N-2d} \, dr \\
% &\leq \frac{q \, \omega_N}{N-2d} 
% \int_{0}^{1} \big( |u'(r)| r^{\tfrac{N-2d-1}{2}} \big)
% \big( |u(r)|^{q-1} r^{\tfrac{N-2d+1}{2}} \big) \, dr \\
% & \leq \frac{q \, \omega_N}{N-2d} 
% \left( \int_{0}^{1} |u'(r)|^2 r^{N-2d-1} \, dr \right)^{1/2}
% \left( \int_{0}^{1} |u(r)|^{2(q-1)} r^{(N-2d+1)} \, dr \right)^{1/2}\\
% & \leq \frac{C \, q }{N-2d} 
% \left( \int_{B_1} \frac{|\nabla u|^2}{|x|^{2d}} \, dx \right)^{q/4}
% \left( \int_{B_1} \frac{|u|^q}{|x|^{2d}} \, dx \right)^{1/2}.
% \end{aligned}
% % \label{eq:radial_q_integral}
% \end{equation*}
% Rearranging terms, we finally deduce
% \begin{equation*}
% \left( \int_{B_1} \frac{|u|^q}{|x|^{2d}} \, dx \right)^{1/q}
% \;\leq\; \left( \frac{C \, q}{N-2d} \right)^{2/q}
% \left( \int_{B_1} \frac{|\nabla u|^2}{|x|^{2d}} \, dx \right)^{1/2}.
% % \label{eq:final_weighted_sobolev}
% \end{equation*}
% \end{proof}

\begin{proof}[Proof of Lemma \ref{lem:weighted_sobolev_inequality}]
Let $u \in C_{c, \mathrm{rad}}^{\infty}(B_R)$ be radial. Take $v(x) = u(Rx)$ and apply part $(ii)$ of Lemma \ref{lem:radial_estimate} to get 
\begin{equation*}
    \left( \int_{B_R} \frac{|u|^q}{|x|^{2d}} \, dx \right)^{1/q} 
\;\leq\; C \left( \int_{B_R} \frac{|\nabla u|^2}{|x|^{2d}} \, dx \right)^{1/2}.
\end{equation*}
Now, let $  u \in C_c^{\infty}(B_R)$ and $u^*$ denotes its symmetric rearrangement, then using the Rearrangement inequality (see Theorem $3.4$ and property $(v)$ in section $3.3$ of \cite{liebloss2001analysis}) and Theorem $8.1$ of \cite{alvino2017radial_polya} with parameters $l = 2d$ and $ k = 0$ (this choice satisfies condition $(ii)$ of Theorem $1.1$), we deduce
\begin{equation*}
\begin{aligned}
 \left( \int_{B_R} \frac{|u|^q}{|x|^{2d}} \, dx \right)^{1/q}  & \leq    \left( \int_{B_R} (|u|^q)^* \left(\frac{1}{|x|^{2d}} \right)^* \, dx \right)^{1/q}  = \left( \int_{B_R} |u^*|^q \frac{1}{|x|^{2d}}  \, dx \right)^{1/q} \\  &\leq C \left( \int_{B_R} \frac{|\nabla u^*|^2}{|x|^{2d}} \, dx \right)^{1/2} \leq C \left( \int_{B_R} \frac{|\nabla u|^2}{|x|^{2d}} \, dx \right)^{1/2}.
\end{aligned}
\end{equation*}
From this the result follows.
\end{proof}



Further, since the weights $|x|^{-\delta}$ are admissible weights in the sense of \cite{heinonen_potential_theory}, we state the following Poincar\'{e} inequality
\begin{lemma}\label{lem:weighted_poincare_inequality}
 Let $\delta > 0$. Then there exists a constant $C > 0$ such that for all $u \in C^{\infty}(B_R)$ the following inequality holds
 \begin{equation*}
     \int_{B_R} | u - (u)_{B_R}|^2 \frac{dx}{|x|^\delta} \leq C R^2 \int_{B_R} |\grad u|^2 \frac{dx}{|x|^{\delta}},
 \end{equation*}
 where $(u)_{B_R} = \frac{1}{|B_{R}|_{d\delta}} \int_{B_{R}} u \frac{dx}{|x|^{\delta}}$ with $|B_R|_{d \delta} = \int_{B_R} \frac{dx}{|x|^{\delta}}$.
\end{lemma}
For general domains, we have the following Sobolev inequality with weights
\begin{lemma}
Let $d \in \left(0, \frac{N-2}{2} \right)$. Then there exists a positive constant $C>0$ such that for all $u \in C_c^{\infty}(B_R)$ the following inequality holds  
\begin{equation}\label{ineq:sobolev_inequality_weight}
 \left( \int_{\om} |u |^{2^*} \frac{dx}{|x|^{2^*d}}\right)^{2/2^*}   \leq  C \int_{\om}  |\grad u |^2 \frac{dx}{|x|^{2d}}.
\end{equation}
\end{lemma}
\begin{proof}
The proof follows from picking the following parameters in \cite{caffarelli1984interpolation_sobolev}
\begin{equation*}
a = \frac{1}{2}, \; \gamma = \beta = \alpha = \sigma = - d, \; r = q = 2^*, \text{ and } p=2. \qedhere
\end{equation*}
\end{proof}




Now we state the ground state representation borrowed from \cite[Formula $(4.3)$]{frank2008hardy} and refined in \cite[Lemma $3.2$]{ghoussoub_fractional_hardy}.\\
\textbf{Ground state representation:} For $0 < s < 1$, $N > 2$, and $0 < d < \frac{N - 2}{2}$, for $u \in C_c^\infty(\mathbb{R}^N \setminus \{0\})$, we define
\begin{equation*}
    w(x) = |x|^{d} u(x).
\end{equation*}
\begin{equation*}
\begin{aligned}
    \frac{C_{N,2s}}{2} \iint_{\rtwon} \frac{|u(x) - u(y)|^2}{|x - y|^{N + 2s}} \,dx\,dy 
    &= \psi_{N,2s}(d) \int_{\mathbb{R}^N} \frac{u^2(x)}{|x|^{2s}} \,dx + \frac{C_{N,2s}}{2} \iint_{\rtwon} \frac{|w(x) - w(y)|^2}{|x - y|^{N + 2s}} \frac{dx}{|x|^{d}} \frac{dy}{|y|^{d}},
\end{aligned}
\end{equation*}
where
\begin{equation*}
    \psi_{N,2s}(d) = 2^{2s} \frac{
        \Gamma\left( \frac{N  {-d}}{2} \right) 
        \Gamma\left( \frac{2s +{d}}{2} \right)
    }{
        \Gamma\left( \frac{N  {-d} - 2s}{2} \right)
        \Gamma\left( \frac{d}{2} \right)
    }, \text{ and } \quad 
     C_{N,2s} = \frac{
        2^{2s} \Gamma\left( \frac{N + 2s}{2} \right)
    }{
        \pi^{N/2} \left| \Gamma\left( -\frac{2s}{2} \right) \right|
    }.
\end{equation*}
Now, from \cite{chen2004multiple}, we have
\begin{equation*}
    \int_{\Omega}  |\nabla w(x)|^2 \,\frac{dx}{|x|^{2 d}} = \int_{\Omega} |\nabla u|^2 \, dx - \mu \int_{\Omega} \frac{u^2}{|x|^2} \,dx.
\end{equation*}
Substituting these into the identity obtained in \eqref{eq:main_problem_ep} after testing with $u$, we obtain the following ground state representation
\begin{equation*}
\begin{aligned}
    \int_{\Omega}  \frac{|\nabla w|^2}{|x|^{2d}} \,dx 
    &+ \varepsilon\frac{2 \psi_{N,2s}(d)}{C_{N,2s}} \int_{\Omega} \frac{ w^2(x)}{|x|^{2s}} \,\frac{dx}{|x|^{2 d}} \\
    &+ \varepsilon\iint_{\mathbb{R}^{2N}} \frac{|w(x) - w(y)|^2}{|x - y|^{N + 2s}} \frac{dx}{|x|^{d}} \frac{dy}{|y|^{d}}
    = \int_{\Omega} \frac{ w^{2^*}}{|x|^{2^* d}} \,dx + \lambda \int_{\Omega}  \frac{w^{q+1}}{|x|^{(q+1) d}} \,dx.
\end{aligned}
\end{equation*}
This holds for all $w(x) = |x|^{d} u(x)$ with $u \in C_c^\infty(\mathbb{R}^N)$. In fact, by Lemma $4.4$ of \cite{dipierro2016fractionalwithHardy}, it also holds for all $u \in H^1(\mathbb{R}^N) \cap H^s(\mathbb{R}^N)$.\\

In view of this ground state representation, if $u$ is a solution of \eqref{eq:main_problem_ep} then we make a change of variable defined by
\begin{equation*}
    u(x) = |x|^{-d} w(x) \quad \text{ with } d = \sqrt{\bar\mu} - \sqrt{\bar\mu - \mu},
\end{equation*}
where $\bar\mu= \left( \frac{N-2}{2} \right)^2$ being the critical Hardy constant. Then we deduce that $w$ satisfies the following equation as a weak solution
\begin{equation}\label{eq:transformed_equation}\tag{$\mathcal{T_{\mu,\la}^{\varepsilon}}$}
\left\{
\begin{aligned}
    -\mathrm{div}\left( \frac{\nabla w(x)}{|x|^{2 d}}\right) + \varepsilon C_c \frac{w}{|x|^{2s + 2d}} + \varepsilon (-\Delta_{d})^s w &= \frac{w^{2^* - 1}}{|x|^{2^* d}} + \lambda  \frac{w^q}{|x|^{(q+1) d}} \quad \text{in } \Omega,\\
    w &= 0 \quad \text{ in } \rnn \setminus \om,
\end{aligned}
\right.
\end{equation}
where the constant \( C_0 = \frac{2 \psi_{N,2s}(d)}{C_{N,2s}} \). \\
The operator \( (-\Delta_{d})^s \) is defined via the duality pairing
\begin{equation*}
    \langle (-\Delta_{d})^s v, \varphi \rangle_{s,d} = 
    \iint_{\rtwon} \frac{(v(x) - v(y))(\varphi(x) - \varphi(y))}{|x - y|^{N + 2s}}  \frac{dx}{|x|^{d}}  \frac{dy}{|y|^{d}}
\end{equation*}
for any \( \varphi \in \dot{H}^{s,d}(\mathbb{R}^N) \), where \( \dot{H}^{s,d}(\mathbb{R}^N) \) is the closure of \( C_c^\infty(\mathbb{R}^N) \) with respect to the norm
\begin{equation*}
    \|\varphi\|_{\dot{H}^{s,d}(\mathbb{R}^N)} = 
    \left( 
        \int_{\mathbb{R}^N} \frac{|\varphi(x)|^{2^*_s}}{|x|^{d 2^*_s}} \,dx 
    \right)^{\frac{1}{2^*}} 
    + 
    \left( 
        \iint_{\rtwon} \frac{|\varphi(x) - \varphi(y)|^2}{|x - y|^{N + 2s}}  \frac{dx}{|x|^{d}}  \frac{dy}{|y|^{d}} 
    \right)^{1/2},
\end{equation*}
where $2^*_s = \frac{2N}{N - 2s}$. Moreover, we need the space $ \dot{H}(\rnn,|x|^{-2d})$ to work with the transformed equation. It is defined as the closure of \( C_c^\infty(\mathbb{R}^N) \) with respect to the norm 
\begin{equation*}
    \|\varphi\|_{\dot{H}(\rnn,|x|^{-2d})} = 
    \left( 
        \int_{\mathbb{R}^N} \frac{|\varphi(x)|^{2^*}}{|x|^{d 2^*}} \,dx 
    \right)^{\frac{1}{2^*}} 
    + 
    \left( \int_{\rnn} | \grad \varphi |^2 dx \right)^2.
\end{equation*}
Having transformed the original equation, we now analyze the qualitative properties of its solutions. One of the key advantages of the transformed formulation is that boundedness of solutions becomes accessible. In order to establish this, we derive a weak Harnack inequality for solutions of the transformed problem~\eqref{eq:transformed_equation}. To simplify notations, we define
\begin{equation*}
    d\mu := \frac{dx}{|x|^{2d}}, \quad 
    d\nu := \frac{dx\,dy}{|x - y|^{N + 2s} |x|^{d} |y|^{d}},
\end{equation*}
and $|E|_{d\mu}$ is the measure of set $E$ with respect to measure $\mu$ given as $|E|_{d\mu} = \int_{E} \frac{1}{|x|^{2 d}} dx$.
\begin{theorem}\label{thm:harnack_inequality_transformed}
Let $d \in (0,\sqrt{\bar{\mu}})$ and $w \in \dot{H}^{s,d}(\mathbb{R}^N) \cap \dot{H}(\rnn,|x|^{-2d})$ be a weak solution of \eqref{eq:transformed_equation}. Then, there exist $r_0 > 0$ such that for $\eta \in [1,\frac{N-2}{2})$ the following inequality holds
\begin{equation*}
        \left( \frac{1}{|B_r|_{d\mu}}\int_{B_r} w^{\eta} \, d\mu \right)^{1/\eta} \leq C \inf_{B_{3r/2}} w,
\end{equation*}
for all $r < r_0$.
\end{theorem}
To prove this Harnack inequality, we first prove a series of lemmas. The first result in the direction is the following lemma, known as the \textit{propagation of positivity}

\begin{lemma}[Propagation of Positivity]\label{lem:propagation_of_positivity}
Let \( w > 0 \) in \( B_R(0) \), with \( 0 < R \leq 1 \), be a supersolution to equation~\eqref{eq:transformed_equation}. Let \( k > 0 \), and suppose that for some \( \sigma \in (0,1] \), we have
\[
    |B_r \cap \{w \geq k\}|_{d\mu} \geq \sigma |B_r|_{d\mu}
\]
with \( 0 < r < \frac{R}{16} \). Then, there exists a constant \( C = C(N, s) \) such that
\[
    |B_{6r} \cap \{w \leq 2\delta k\}|_{d\mu} \leq \frac{C}{\sigma \log\left( \frac{1}{2\delta} \right)} |B_{6r}|_{d\mu}
\]
for all \( \delta \in \left(0, \frac{1}{4}\right) \).
\end{lemma}
\begin{proof}
Choose a cut-off function \( \varphi \in C_c^{\infty}(\rnn)\) such that 
\[
    \text{Supp}(\varphi) \subset B_{7r}, \quad 0 \leq \varphi \leq 1 \quad \text{and} \quad \varphi \equiv 1 \text{ in } B_{6r}, \quad |\nabla \varphi| \leq \frac{C}{r}.
\]
Choosing test function \( \eta = w^{-1} \varphi^2 \) in \eqref{eq:transformed_equation} and Young's inequality give
\begin{align*}
    0 \leq & \int_{\om} \grad w \nabla \left( w^{-1} \varphi^2 \right) d\mu + \varepsilon C_0 \int_{\om} \frac{\varphi^2}{|x|^{2s}} d\mu   +\varepsilon \iint_{\rtwon} \left( w(x) - w(y) \right) \left(  \frac{\varphi^2(x)}{w(x)} - \frac{\varphi^2(y)}{w(y)} \right) d\nu \\
    \leq & -\frac{1}{2} \int_{\om}  \varphi^2 w^{-2} |\nabla w|^2\, d\mu + 2 \int_{\om} |\nabla \varphi|^2 \,d\mu + \varepsilon C_0 \int_{\om} \frac{\varphi^2}{|x|^{2s}} d\mu   +\varepsilon I_3 \\
    \leq & -\frac{1}{2} \int_{\om}  \varphi^2 w^{-2} |\nabla w|^2\, d\mu + C  |B_{6r}|_{d\mu} r^{-2} + C |B_{6r}|_{d\mu} r^{-2s} + I_3.
\end{align*}
For \( I_3 \), we break the integral as follows
\begin{align*}
    I_3 = \iint_{B_{8r} \times B_{8r}} d\nu + \iint_{B_{8r} \times \mathbb{R}^N \setminus B_{8r}} d\nu + \iint_{\mathbb{R}^N \setminus B_{8r} \times B_{8r}} d\nu = I_{3,1} + I_{3,2} + I_{3,3}.
\end{align*}
Using the same idea of \cite[Lemma $1.3$]{palatucci_castro2016local}, we obtain
\begin{align*}
    I_{3,1} & \leq -\frac{1}{C} \iint_{B_{8r} \times B_{8r}} \left| \log \left(\frac{w(x)}{w(y)}\right) \right|^2 \varphi^2(y) d\nu + C \iint_{B_{8r} \times B_{8r}} \left| \varphi(x) - \varphi(y) \right|^2 d\nu \\
    & \leq -\frac{1}{C} \iint_{B_{6r} \times B_{6r}} \left| \log \left(\frac{w(x)}{w(y)}\right) \right|^2 d\nu + C \iint_{B_{8r} \times B_{8r}} \left| \varphi(x) - \varphi(y) \right|^2 d\nu \\
    & \leq -\frac{1}{C} \iint_{B_{6r} \times B_{6r}} \left| \log \left(\frac{w(x)}{w(y)}\right) \right|^2 d\nu + C \lv \grad \varphi \rv_{L^{\infty}}^2 \iint_{B_{8r} \times B_{8r}} \frac{|x-y|^{2-N-2s}}{|x|^{d}|y|^{d}} \, dy \, dx \\
    & \leq -\frac{1}{C} \iint_{B_{6r} \times B_{6r}} \left| \log \left(\frac{w(x)}{w(y)}\right) \right|^2 d\nu + \frac{C}{r^2} \left( \int_{B_{8r}} |x|^{-2d N/(N+2-2s)}dx\right)^{(N+2-2s)/N} \\
    & \leq -\frac{1}{C} \iint_{B_{6r} \times B_{6r}} \left| \log \left(\frac{w(x)}{w(y)}\right) \right|^2 d\nu  +  C |B_{6r}|_{d\mu} r^{-2s},
\end{align*}
where the inequality in the penultimate line follows from the HLS inequality \cite{liebloss2001analysis}. Now
\begin{align*}
I_{3,2} = I_{3,3} &= \iint_{ B_{8r} \times \mathbb{R}^N \setminus B_{8r} } (w(x) - w(y)) w^{-1}(x) \varphi^2(x) \, d\nu \\
& \leq \iint_{B_{7r} \times \mathbb{R}^N \setminus B_{8r}} \frac{1}{|x - y|^{N + 2s}} \frac{1}{|x|^{d}}  \frac{1}{|y|^{d}} \, dx \, dy \\
&\leq \int_{B_{7r}} \frac{1}{|x|^{d}} \, dx \int_{\mathbb{R}^N \setminus B_{8r}} \frac{1}{|y|^{d}}  \frac{1}{( |y|-7r)^{N + 2s}} \, dy\\
& \leq C |B_{6r}|_{d\mu} r^{-2s}.
\end{align*}
Combining all estimates and the fact that \( r^{-2s} < r^{-2} \) for $r < 1$, we deduce that
% \begin{align*}
% 0 \leq &  ~I_1 +  I_2 +  I_3 \\
% \leq &~  \frac{-1}{2} \int_{\om}  \varphi^2 w^{-2} |\nabla w|^2\, d\mu + 2 \int_{\om} |\nabla \varphi|^2 \,d\mu +  C \, |B_{6r}|_{d\mu} \, r^{-2s}  - \frac{1}{C} \iint_{B_{6r} \times B_{6r}} \left| \log \left( \frac{w(x)}{w(y)} \right) \right|^2 d\nu
% \end{align*}
% which can be rewritten as
\begin{align*}
\int_{B_{6r}} |\nabla(\log w)|^2 \, d\mu  &\leq \int_{\Omega}  \frac{\varphi^2}{w^2} |\nabla w|^2 \, d\mu +  \iint_{B_{6r} \times B_{6r}} \left| \log \left( \frac{w(x)}{w(y)} \right) \right|^2 d\nu \\
&\leq C \left[ |B_{6r}|_{d\mu} r^{-2} +  |B_{6r}|_{d\mu} r^{-2s} \right] \leq C \, r^{-2} \, |B_{6r}|_{d\mu}.
\end{align*}
% Using above and the fact that \( r^{-2s} < r^{-2} \) for $r < 1$, we get
% \[
% \int_{B_{6r}} |\nabla(\log w)|^2 \, d\mu \leq C \, r^{-2} \, |B_{6r}|_{d\mu}.
% \]
For any \( \delta \in (0, \tfrac{1}{4}) \), define $v = \left[ \min \left\{ \log \left( \frac{1}{2\delta} \right),~ \log \left( \frac{k}{w} \right) \right\} \right]_+$ and
% \[
% v = \left[ \min \left\{ \log \left( \frac{1}{2\delta} \right),~ \log \left( \frac{k}{w} \right) \right\} \right]_+
% \]
denote $(v)_{B_{6r}} = \frac{1}{|B_{6_r}|_{d\mu}} \int_{B_{6r}}v \, dx$. Then Lemma \ref{lem:weighted_poincare_inequality} yields
\begin{align*}
\int_{B_{6r}} |v(x) - (v)_{B_{6r}}| \, d\mu & \leq \left( \int_{B_{6r}} |v(x) - (v)_{B_{6r}}|^2 \, d\mu \right)^{1/2}   |B_{6r}|_{d\mu}^{1/2} \\
&\leq  C \, r \left( \int_{B_{6r}} |\nabla v|^2 \, d\mu \right)^{1/2}   |B_{6r}|_{d\mu}^{1/2} \\
& \leq C \, r \left( \int_{B_{6r}} |\nabla \log w|^2 \, d\mu \right)^{1/2}   |B_{6r}|_{d\mu}^{1/2} \\
& \leq  C \, |B_{6r}|_{d\mu}.
\end{align*}
Now, we know $ \{ v = 0 \} =  \{ w \geq k \}$ then $    |B_{6r} \cap \{ v = 0 \}|_{d\mu} \geq \frac{\sigma}{6^{N+2d}} |B_{6r}|_{d\mu}$.  Moreover,
% \begin{equation*} 
%     |B_{6r} \cap \{ v = 0 \}|_{d\mu} \geq \frac{\sigma}{6^{N+2d}} |B_{6r}|_{d\mu}.
% \end{equation*}
% Now,

\small{\begin{align*}
    \log \left( \frac{1}{2\delta} \right) = \frac{1}{|B_{6r} \cap \{ v = 0 \}|} \int_{B_{6r} \cap \{ v = 0 \}} \left( \log \left( \frac{1}{2\delta} \right) - v(x) \right) \, dx \leq \frac{6^{N+2d}}{\sigma} \left[ \log \left( \frac{1}{2\delta} \right) - (v)_{B_{6r}} \right].
\end{align*}}

Integrating the above inequality with respect to the measure $\mu$, we have
\begin{equation*}
    \left|\{ v = \log \left( \tfrac{1}{2\delta} \right) \} \cap B_{6r}\right|_{d\mu}   \log \left( \tfrac{1}{2\delta} \right)
    \leq \frac{6^{N+2d}}{\sigma} \int_{B_{6r}} |v(x) - (v)_{B_{6r}}| \, d\mu
    \leq \frac{C}{\sigma} |B_{6r}|_{d\mu}.
\end{equation*}
Therefore, for all $\delta \in \left( 0, \tfrac{1}{4} \right)$, we obtain the estimate
\begin{equation*} 
    |B_{6r} \cap \{ w \leq 2\delta k \}|_{d\mu} 
    \leq \frac{C}{\sigma \log \left( \frac{1}{2\delta} \right)} |B_{6r}|_{d\mu}. \qedhere
\end{equation*}
\end{proof}
\medskip

\begin{lemma}\label{lem:lower_bound_Harnack}
Assuming the hypothesis of Lemma \ref{lem:propagation_of_positivity}, there exists $\delta \in (0, \tfrac{1}{4})$ such that
\begin{equation*} 
    \inf_{B_{4r}} w \geq \delta k.
\end{equation*}
\end{lemma}
\begin{proof}
Take a smooth function $\varphi \in C_c^{\infty}(\om)$ with $\text{supp}(\varphi) \subseteq B_{\rho}$ with $\rho \in [r,6r]$. Testing the equation with the function $\eta = v_{\ell} \varphi^2$, where $v_{\ell} = (\ell - w)_+$. Then taking $\Omega_\ell = \Omega \cap \{x :  w(x) < \ell \}$ and using Young's inequality, we obtain
\begin{align*}
    0 &\leq \int_{\Omega} \nabla w   \nabla \eta  \, d\mu +\varepsilon C_0 \int_{\Omega} \frac{w \eta}{|x|^{2s}} \, d\mu 
    + \varepsilon\iint_{\mathbb{R}^{2N}} (w(x) - w(y)) (\eta(x) - \eta(y)) \, d\nu \\
    & \leq  -\frac{1}{2} \int_{\om_{\ell}} |\nabla v_{\ell}|^2 \varphi^2 \, d\mu 
+ 2 \int_{\om_{\ell}} |\nabla \varphi|^2 v_{\ell}^2 \, d\mu  -C_0 \int_{\om_{\ell}} \frac{v_{\ell}^2 \varphi^2}{|x|^{2s}} \, d\mu 
+ C_0 \int_{\om_{\ell}} \frac{\ell v_{\ell} \varphi^2}{|x|^{2s}} \, d\mu + \varepsilon I_3 \\
& \leq  -\frac{1}{2} \int_{\om_{\ell}} |\nabla v_{\ell}|^2 \varphi^2 \, d\mu 
+ 2 \int_{\om_{\ell}} |\nabla \varphi|^2 v_{\ell}^2 \, d\mu -\frac{C_0}{2} \int_{\om_{\ell}} \frac{v_{\ell}^2 \varphi^2}{|x|^{2s}} \, d\mu 
+ \frac{C_0}{2} \int_{\om_{\ell}} \frac{\ell^2 \varphi^2}{|x|^{2s}} \, d\mu + I_3.
\end{align*}
% Let $\Omega_\ell = \Omega \cap \{x :  w(x) < \ell \}$. Then following as in previous lemma, we get
% \begin{equation*}
%     I_1 = \int_{\Omega} \nabla w   \nabla (v_{\ell} \varphi^2) \, d\mu \leq  -\frac{1}{2} \int_{\om_{\ell}} |\nabla v_{\ell}|^2 \varphi^2 \, d\mu 
% + 2 \int_{\om_{\ell}} |\nabla \varphi|^2 (v_{\ell})^2 \, d\mu
% \end{equation*}
% Using Young's inequality in $I_2$ yields
% \begin{equation*}
% \begin{aligned}
% I_2 = -C_0 \int_{\om_{\ell}} \frac{v_{\ell}^2 \varphi^2}{|x|^{2s}} \, d\mu 
% + C_0 \int_{\om_{\ell}} \frac{\ell v_{\ell} \varphi^2}{|x|^{2s}} \, d\mu \leq  -\frac{C_0}{2} \int_{\om_{\ell}} \frac{v_{\ell}^2 \varphi^2}{|x|^{2s}} \, d\mu 
% + \frac{C_0}{2} \int_{\om_{\ell}} \frac{\ell^2 \varphi^2}{|x|^{2s}} \, d\mu.
% \end{aligned}
% \end{equation*}
For $I_3$, we proceed as in \cite[Lemma $3.2$]{castro2014nonlocal_harnack}, leading to
\begin{equation*} 
\begin{aligned}
2 I_3 &\leq  \iint_{B_\rho \times B_\rho} \left( - \left| v_{\ell}(x) \varphi(x) - v_{\ell}(y) \varphi(y) \right|^2 + C \left( \max\{ v_{\ell}(x), v_{\ell}(y) \} \right)^2 \left| \varphi(x) - \varphi(y) \right|^2  \right)\, d\nu \\
&\quad + C \left( \sup_{x \in \text{supp}(\varphi)} \int_{\mathbb{R}^N \setminus B_{\rho}} \frac{v_{\ell}(y)}{|x - y|^{N+2s}} \, \frac{dy}{|y|^{d}} \right)\int_{B_{\rho}} \frac{v_{\ell}(x) \varphi(x)^2}{|x|^{d}} \, dx .
\end{aligned}
\end{equation*}
Combining all, we obtain
\begin{align}
  &\int_{\om_{\ell}} |\nabla v_{\ell}|^2 \varphi^2 \, d\mu +  C_0 \int_{\om_{\ell}} \frac{v_{\ell}^2 \varphi^2}{|x|^{2s}} \, d\mu + \iint_{B_\rho \times B_\rho} \left| v_{\ell}(x) \varphi(x) - v_{\ell}(y) \varphi(y) \right|^2\, d\nu \notag\\
  &\leq  2 \int_{\om_{\ell}} |\nabla \varphi|^2 v_{\ell}^2 \, d\mu 
+ \frac{C_0}{2} \int_{\om_{\ell}} \frac{\ell^2 \varphi^2}{|x|^{2s}} \, d\mu +
C \iint_{B_\rho \times B_\rho}  \left( \max\{ v_{\ell}(x), v_{\ell}(y) \} \right)^2 \left| \varphi(x) - \varphi(y) \right|^2  d\nu \notag\\
&\quad + C \left( \sup_{x \in \text{supp}(\varphi)} \int_{\mathbb{R}^N \setminus B_{\rho}} \frac{v_{\ell}(y)}{|x - y|^{N+2s}} \, \frac{dy}{|y|^{d}} \right)\int_{B_{\rho}} \frac{v_{\ell}(x) \varphi(x)^2}{|x|^{d}} \, dx \notag\\
& \leq C(K_1 + K_2 + K_3 + K_4). \label{eq:lower_bd_harnack_eq1}
\end{align}
Now, we recast the inequality \eqref{eq:lower_bd_harnack_eq1} into a format suitable for applying a iteration lemma. To this end, let
\begin{equation*}
\begin{aligned}
\ell \equiv \ell_j &= \delta k + 2^{-j} \delta k, \qquad
\rho = \rho_j :=4r +  2^{1-j}r,& \qquad 
\tilde{\rho}_j := \frac{\rho_{j+1} + \rho_j}{2}.
\end{aligned}
\end{equation*}
We have $\ell_j - \ell_{j+1} = 2^{-j-1} \delta k$ which implies $\ell_j - \ell_{j+1} \geq 2^{-j-2} \ell_j$. Define the function
\begin{equation}\label{eq:vj_lower_bound}
v_{\ell} \equiv v_j := (\ell_j - w)_+ \geq (\ell_j - \ell_{j+1}) \chi_{\{ w < \ell_{j+1} \}} \geq 2^{-j-2} \ell_j \chi_{\{ w < \ell_{j+1} \}}.
\end{equation}
Define the balls $B_j$ and cut-off functions $\varphi_j$ for all \( j = 0, 1, \dots \) as
\begin{equation*}
B_j := B_{\rho_j}(0), \quad 
\varphi_j \in C_c^\infty(B_{\tilde{\rho}_j}), 
\quad \text{with} \quad 
0 \leq \varphi_j \leq 1, \quad 
\varphi_j \equiv 1 \text{ in } B_{{j+1}}, \quad 
|\nabla \varphi_j| \leq C\frac{2^{j}}{r},
\end{equation*}
and denote $\widetilde{B}_j = B_j \cap \{ w < \ell_j \}$.
Now, estimate $K_1$ in terms of $\ell_j$ as
\begin{equation*}
K_1 = \int_{B_j} |\nabla \varphi_j|^2 v_j^2 \, d\mu 
\leq C \ell_j^2   \frac{2^{2j}}{r^2} 
\left| B_j \cap \{ w < \ell_j \} \right|_{d\mu} = C \ell_j^2   \frac{2^{2j}}{r^2} \widetilde{B}_j.
\end{equation*}
Next, we estimate the Hardy term by noting that
\begin{align*}
\int_{\widetilde{B}_j} \frac{1}{|x|^{2s}} \, d\mu 
&= \frac{1}{|B_j|} \int_{B_j} \left( \int_{\widetilde{B}_j} 
\frac{1}{|x |^{2s + 2d}} \, dx \right) dy = \frac{1}{|B_j|} (J_1 + J_2 + J_3),
\end{align*}
where $J_i = \iint_{R_i \cap (\widetilde{B_j} \times B_j)} \frac{1}{|x|^{2s + 2d}} dx dy$ for $i = 1,2,3$ with regions are defined as follows
\begin{equation}\label{eq:region_breaking}
R_1 = \left\{ \frac{1}{2} \leq  \frac{|x|}{|y|} < 2 \right\}, \quad
R_2 = \left\{ |x| < \frac{1}{2}|y| \right\}, \quad
R_3 = \left\{ |y| \leq \frac{1}{2}|x| \right\}.
\end{equation}
For the first term
\begin{equation*}
J_1 \leq C \iint_{R_1 \cap (\widetilde{B_j} \times B_j)} \frac{1}{|y|^{2s}} 
  \frac{1}{|x|^{2d}} \, dx \, dy 
\leq C \left| B_j \cap \{ w < \ell_j \} \right|_{d\mu}   r^{N - 2s}.
\end{equation*}
Similarly, for the third term
\begin{equation*}
J_3 \leq C \left| B_j \cap \{ w < \ell_j \} \right|_{d\mu}   r^{N - 2s}.
\end{equation*}
Finally, for $J_2$, Lemma \ref{lem:stein_strong_lemma} with parameters $\delta = -2s + \frac{N}{2}$ and $p=2$ gives
\begin{align*} 
J_2 
&\leq C \int_{\widetilde{B}_j} \frac{1}{|x|^{2 s  + 2 d - N + \delta}}   |x|^{-N + \delta} 
\left( \int_{|y| \leq |x|} |y|^{-\delta}   \frac{1}{|y|^{-\delta}} \, dy \right) dx\\
&\leq C \int_{\widetilde{B}_j} \frac{1}{|x|^{2 s  + 2 d - N + \delta}}   |x|^{-N/2}
\left( \int_{B_j} \frac{1}{|y|^{-2\delta}} \, dy \right)^{1/2} dx\\
&\leq C   r^{N - 2s}   \left| B_j \cap \{ w < \ell_j \} \right|_{d\mu}.
\end{align*}
Thus we obtain that
\begin{equation*}
K_2
\leq C   r^{ - 2s}   \left| B_j \cap \{ w < \ell_j \} \right|_{d\mu}.
\end{equation*}
For the fractional term $K_3$, we have
\begin{align*} 
K_3 &
\leq  \ell_j^2 \iint_{B_j \times \widetilde{B}_j} 
\frac{\| D\varphi_j \|^2_\infty}{|x - y|^{N + 2s-2}} 
\frac{dy}{|y|^{d}}   \frac{dx}{|x|^{d}}
\leq C \ell_j^2 2^{2j} r^{-2} 
\iint_{B_j \times \widetilde{B}_j} 
\frac{1}{|x - y|^{N + 2s-2}} 
  \frac{dy}{|y|^{d}}   \frac{dx}{|x|^{d}} \\
&= C \ell_j^2 2^{2j} r^{-2} \left( L_1 + L_3 + L_2 \right),
\end{align*}
where $L_i = \iint_{R_i \cap [B_j \times (B_j \cap \{ w < \ell_j \})]}$ with regions $R_1, R_2$ and $R_3$ are same as in \eqref{eq:region_breaking}.

We now estimate each region. For \( L_1 \), we obtain
\begin{equation*}
L_1 \leq C \iint_{R_1 \cap [B_j \times (B_j \cap \{ w < \ell_j \})]} \frac{1}{|x - y|^{N + 2s - 2}}   \frac{dy}{|y|^{2d}}   dx
\leq C | B_j \cap \{ w < \ell_j \} |_{d\mu}   r^{2 - 2s}.
\end{equation*}
Similarly, for \( L_2 \), we have
\begin{equation*}
L_2 \leq C |B_j \cap \{ w < \ell_j \} |_{d\mu}   r^{2 - 2s}.
\end{equation*}
For \( L_3 \), we invoke the Lemma \ref{lem:stein_strong_lemma} with parameters \( q \) and \( \de \) satisfying the relations $q = \frac{N}{N+d - \delta}$, $ \delta < \frac{N(q-1)}{q}$, and $q > \frac{N}{N-d}$ and noting that in $R_2$ we have $|x-y| \geq |y| - |x| \geq |x|$, we estimate
\begin{align*} 
L_3 &\leq C \int_{B_j \cap \{ w < \ell_j \}} \frac{1}{|y|^{d}} 
\left( \int_{\{ 2|x| \leq |y| \}} \frac{1}{|x|^{d}}   \frac{1}{|x |^{N + 2s - 2}} \, dx \right) dy\\
& \leq C \int_{\widetilde{B}_j} \frac{1}{|y|^{d - N + \delta}} |y|^{-N+\delta}
\left( \int_{ \{|x| \leq |y|\} } \frac{|x|^{-\delta}}{|x |^{- \delta + d+ N + 2s - 2}} \, dx \right) dy\\
& \leq C \int_{\widetilde{B}_j} \frac{1}{|y|^{2d}} 
\left( \int_{B_j} \left( \frac{1}{|x|^{- \delta + d + N + 2s - 2}} \right)^q \, dx \right)^{1/q} dy\\
& \leq C \left| B_j \cap \{ w < \ell_j \} \right|_{d\mu}   r^{2 - 2s}.
\end{align*}
Therefore, we have 
\begin{equation*}
    K_3 \leq C \ell_j^2   2^{2j}   r^{-2s} |\widetilde{B}_j|_{d\mu}.
\end{equation*}
% \begin{equation*}
% \iint_{B_j \times B_j} 
% \left( \max \left\{ v_j(x), v_j(y) \right\} \right)^2 
% \left( \varphi_j(x) - \varphi_j(y) \right)^2 \, d\nu 
% \leq C \ell_j^2   2^{2j}   r^{-2s} 
% \left| B_j \cap \{ w < \ell_j \} \right|_{d\mu}.
% \end{equation*}
We now estimate $K_4$,
\begin{align*}
K_4 
% \int_{B_j} v_j(x) \varphi_j^2(x) \frac{dx}{|x|^{d}} 
% &\left( \sup_{x \in \text{supp}(\varphi_j)} \int_{\rnn \setminus B_j} 
% \frac{v_j(y)}{|x - y|^{N+2s}}   \frac{dy}{|y|^{d}} \right) \\
\leq C \ell_j^2 \int_{\widetilde{B}_j} \int_{\rnn \setminus B_j} 
\frac{2^{j(N+2s)}}{|y|^{N+2s}}   
\frac{dy}{|y|^{d}}   \frac{dx}{|x|^{d}}
\leq C \ell_j^2   2^{j(N+2s)}   r^{-2s} \left| \widetilde{B}_j \right|_{d\mu}.
\end{align*}



Using values obtained of $K_i$ in \eqref{eq:lower_bd_harnack_eq1} and \eqref{eq:vj_lower_bound}, we obtain the following estimate
\begin{align*}
 (\ell_j -  \ell_{j+1})^2 
&\left( \frac{| B_{j+1} \cap \{ w < \ell_{j+1}  \}|_{d\mu}}{|B_{j+1}|_{d\mu}}   \right)^{\frac{2}{2^*}} \leq \left( \frac{1}{|B_{j+1}|_{d\mu}} \int_{B_{j+1}} (v_j \varphi_j)^{2^*} \, d\mu \right)^{\frac{2}{2^*}} \\
&   \leq  C\left( \frac{1}{|B_{j}|_{d\mu}} \int_{B_{j}} (v_j \varphi_j)^{2^*} \, d\mu \right)^{\frac{2}{2^*}} \\
& \leq C \left( \frac{r^2}{|B_j|_{d\mu}} \int_{B_j} |\nabla(v_j \varphi_j)|^2 \, d\mu \right)\\
&\leq \frac{C r^2}{|B_j|_{d\mu}} \left[ 
\int_{B_j} |\nabla(v_j \varphi_j)|^2 \, d\mu 
+   \int_{\widetilde{B}_j} \frac{(v_j \varphi_j)^2}{|x|^{2s}} \, d\mu 
\right. \\
& \quad \left.+ \iint_{B_j \times B_j} \frac{|v_j(x) \varphi_j(x) - v_j(y) \varphi_j(y)|^2}
{|x - y|^{N+2s}} \, \frac{dy}{|y|^{d}} \frac{dx}{|x|^{d}}
\right]\\
& \leq \frac{C r^2}{|B_j|_{d\mu}} \left[ 
\ell_j^2   \frac{2^{2j}}{r^2}   |\widetilde{B}_j|_{d\mu} 
+  \ell_j^2   \frac{|\widetilde{B}_j|_{d\mu}}{r^{2s}} 
+  \ell_j^2   2^{j(N+2s)}   \frac{|\widetilde{B}_j|_{d\mu}}{  r^{2s}} 
\right]\\
& \leq C   2^{j(N+2s + 2)}   \ell_j^2   
\frac{\left| \widetilde{B}_j  \right|_{d\mu}}{|B_j|_{d\mu}},
\end{align*}
where $C$ is a constant independent of $\varepsilon$.
Now, denote $A_j = \frac{\left| \widetilde{B}_j  \right|_{d\mu}}{|B_j|_{d\mu}}$, then we obtain
\begin{equation*}
A_{j+1}^{\frac{2}{2^*}} \leq C   \frac{\ell_j^2   2^{j(N+2s + 2)}}{(\ell_j - \ell_{j+1})^2}   A_j \leq C 2^{j(N+2s + 4)}   A_j.
\end{equation*}
Following the iteration scheme(\cite[Lemma $7.1$]{giusti2003iteration_lemma}), we conclude $\lim_{j \to \infty} A_j = 0$. Using this completes the proof.
\end{proof}
% Reverse H\"{o}lder Inequality
Next, we derive a reverse H\"{o}lder inequality for the solutions of \eqref{eq:transformed_equation}.

\begin{lemma}[Reverse H\"{o}lder Inequality]
Let \( w \) be a supersolution to \eqref{eq:transformed_equation}. Then, for all \( 0 < \gamma_1 < \gamma_2 < \frac{N}{N - 2} \), we have:
\begin{equation} \label{lem:reverse_holder}
\left( \frac{1}{|B_r|_{d\mu}} \int_{B_r} w^{\gamma_2} \, d\mu \right)^{1/\gamma_2}
\leq C \left( \frac{1}{|B_{3r/2}|_{d\mu}} \int_{B_{3r/2}} w^{\gamma_1} \, d\mu \right)^{1/\gamma_1}.
\end{equation}
\end{lemma}
\begin{proof}
Let \( q \in (1,2) \) and \( n \in \ntrl\). Define \( \widetilde{w} = w + \frac{1}{n} \), and let \( \psi \) be a cut-off function with \( \operatorname{Supp}(\psi) \subseteq B_{\tau r} \), such that
\begin{equation*}
    \psi \equiv 1 \text{ in } B_{\tau' r} \quad \text{and} \quad |\nabla \psi| \leq \frac{C}{(\tau - \tau')r}, \quad \text{where } \frac{1}{2} \leq \tau' < \tau < \frac{3}{2}.
\end{equation*}
Taking the test function \( \eta = \widetilde{w}^{1-q} \psi^2 \), we deduce the following inequality
\begin{equation}\label{eq:rev_holder_eq1}
\begin{aligned}
0 \leq 
& \int_{B_{\tau r}} \nabla w  \nabla\left( \widetilde{w}^{1-q} \psi^2 \right) \, d\mu 
+ \varepsilon C_0 \int_{B_{\tau r}} \frac{w\widetilde{w}^{1-q} \psi^2}{|x|^{2s}} \, d\mu \\
&+ \varepsilon \iint_{\mathbb{R}^{2N}} \left( \widetilde{w}(x) - \widetilde{w}(y) \right) 
\left( \frac{\psi^2(x)}{\widetilde{w}^{q-1}(x)} - \frac{\psi^2(y)}{\widetilde{w}^{q-1}(y)} \right) \, d\nu\\
= &~ I_1 + \varepsilon C_0 I_2 + \varepsilon I_3 \leq I_1 + C_0 I_2 + I_3.
\end{aligned}
\end{equation}
To estimate \( I_1 \), we apply Young's inequality with \( \epsilon = \frac{(q - 1)}{2} \), yielding
\begin{equation}
\begin{aligned}
I_1 
&= - (q - 1) \int_{B_{\tau r}} |\nabla w|^2 \, \widetilde{w}^{-q} \psi^2 \, d\mu 
+ \int_{B_{\tau r}} \nabla w  \nabla \psi \, \widetilde{w}^{1-q} \psi \, d\mu \\
&\leq -\frac{(q - 1)}{2} \int_{B_{\tau r}} |\nabla w|^2 \, \widetilde{w}^{-q} \psi^2 \, d\mu 
+ \frac{4}{(q - 1)} \int_{B_{\tau r}} |\nabla \psi|^2 \, \widetilde{w}^{2 - q} \, d\mu \\
&\leq -\frac{(q - 1)}{2} \int_{B_{\tau r}} |\nabla w|^2 \, \widetilde{w}^{-q} \psi^2 \, d\mu 
+ \frac{C}{(\tau - \tau')^2 r^2} \int_{B_{\tau r}} \widetilde{w}^{2 - q} \, d\mu.
\end{aligned}
\label{eq:I1_estimate}
\end{equation}
For \( I_2 \), breaking the region into subregions as in \eqref{eq:region_breaking}, and estimating, we obtain
\begin{equation}
    I_2 \leq C \, r^{-2s} \int_{B_{\tau r}} w \widetilde{w}^{1 - q} \, d\mu.
    \label{eq:I2_estimate}
\end{equation}
For \( I_3 \), we follow the approach used in \cite[Lemma $3.7$]{peral2016CZ}, to get
\begin{equation}
\begin{aligned}
    I_3 \leq & -C \iint_{B_{\tau r} \times B_{\tau r}} 
    \left( \frac{\psi}{\widetilde{w}^{(q - 2)/2}}(x) - \frac{\psi}{\widetilde{w}^{(q - 2)/2}}(y) \right)^2 \, d\nu + C \, \frac{r^{-2s}}{(\tau - \tau')^2} \int_{B_{\tau r}} \widetilde{w}^{2 - q} \, d\mu.
\end{aligned}
\label{eq:I3_estimate}
\end{equation}


Using estimates \eqref{eq:I1_estimate}, \eqref{eq:I2_estimate}, and \eqref{eq:I3_estimate} in \eqref{eq:rev_holder_eq1}, and $r^{-2s} \leq r^{-s}$ for $r \leq 1$, we deduce the inequality
\small{\begin{equation*}
\begin{aligned}
    \int_{B_{\tau r}} \left| \nabla\left( \widetilde{w}^{(2 - q)/2} \right) \right|^2 \psi^2 \, d\mu + \iint_{B_{\tau r} \times B_{\tau r}} 
    \left( \frac{\psi}{\widetilde{w}^{(q - 2)/2}}(x) - \frac{\psi}{\widetilde{w}^{(q - 2)/2}}(y) \right)^2 \, d\nu \leq \frac{C r^{-2}}{(\tau - \tau')^2} \int_{B_{\tau r}} \widetilde{w}^{2 - q} \, d\mu.
    % & \leq \frac{C r^{-2}}{(\tau - \tau')^2} \int_{B_{\tau r}} \widetilde{w}^{2 - q} \, d\mu 
    % + \frac{C r^{-2s}}{(\tau - \tau')^2} \int_{B_{\tau r}} \widetilde{w}^{2 - q} \, d\mu \\
    % & \leq \frac{C r^{-2}}{(\tau - \tau')^2} \int_{B_{\tau r}} \widetilde{w}^{2 - q} \, d\mu.
\end{aligned}
\end{equation*}}


Now, from the above and applying the weighted Sobolev inequality in Lemma \ref{lem:weighted_sobolev_inequality}, we deduce
\begin{equation*}
\begin{aligned}
\left( \frac{1}{|B_{\tau' r}|_{d\mu}} \int_{B_{\tau' r}} \widetilde{w}^{(2 - q)2^*/2 } \, d\mu \right)^{\frac{N - 2}{N}}
&\leq C \left( \frac{1}{|B_{\tau r}|_{d\mu}} \int_{B_{\tau r}} \left( \widetilde{w}^{(2 - q)/2} \psi \right)^{2^*} \, d\mu \right)^{\frac{N - 2}{N}}\\
&\leq \frac{C \tau^2 r^2}{|B_{\tau r}|_{d\mu}} \int_{B_{\tau r}} 
\left| \nabla\left( \widetilde{w}^{(2 - q)/2} \psi \right) \right|^2 \, d\mu \\
&\leq \frac{C}{|B_{\tau r}|_{d\mu} (\tau - \tau')^2} \int_{B_{\tau r}} \widetilde{w}^{2 - q} \, d\mu.
\end{aligned}
\end{equation*}
Applying the Monotone Convergence Theorem and letting \( n \to \infty \). Since \( 1 < q < 2 \) is arbitrary, H\"{o}lder's inequality implies the desired estimate.
\end{proof}


The next two lemmas follows directly from \cite[Lemma $4.10$ and $4.11$]{dipierro2016fractionalwithHardy}.
\begin{lemma} \label{lemma:measure_enlargement}
Assume that \( E \subset B_r(x_0) \) is a measurable set. For \( \bar{\delta} \in (0,1) \), we define the \emph{enlargement} of \( E \) as
\begin{equation*} 
[E]_{\bar{\delta}} := \bigcup_{\rho > 0} \left\{ B_{3\rho}(x) \cap B_r(x_0) \,:\, x \in B_r(x_0),\; |E \cap B_{3\rho}(x)|_{d\mu} > \bar{\delta} |B_{\rho}(x)|_{d\mu} \right\}.
\end{equation*}
Then, there exists a constant \( \tilde{C} \), depending only on \( N \), such that one of the following holds
\begin{enumerate}
    \item \( |[E]_{\bar{\delta}}|_{d\mu} \geq \dfrac{\tilde{C}}{\bar{\delta}} |E|_{d\mu}, \) or
    \item \( [E]_{\bar{\delta}} = B_r(x_0). \)
\end{enumerate}
\end{lemma}
\begin{lemma}\label{lemma:main_supersolution}
Assume that \( w \) is a nonnegative supersolution to \eqref{eq:transformed_equation}. Then, there exists \( \eta \in (0,1) \), depending only on \( N \), such that the following inequality holds
\begin{equation*} 
\left( \frac{1}{|B_r|_{d\mu}} \int_{B_r} w^{\eta} \, d\mu(x) \right)^{1/\eta} \leq C \inf_{B_r} w.
\end{equation*}
\end{lemma}
Note that the constant in the above lemma is independent of $\varepsilon$, as it follows from Lemmas~\ref{lem:lower_bound_Harnack} and~\ref{lemma:measure_enlargement}, whose constants are also independent of $\varepsilon$. Consequently, the constant in the subsequent proof of the Harnack inequality is likewise independent of $\varepsilon$.
% Harnack Inequality
\begin{proof}[Proof of Theorem \ref{thm:harnack_inequality_transformed}(Harnack Inequality):]
Applying Lemma \ref{lem:reverse_holder} with $\gamma_1 = \eta$ and $\gamma_2 = q$ and Lemma \ref{lemma:main_supersolution}, we get 
\begin{equation*}
    \left( \frac{1}{|B_r|_{d\mu}} \int_{B_r} w^{q} \, d\mu(x) \right)^{1/q} \leq C  \left( \frac{1}{|B_{3r/2}|_{d\mu}} \int_{B_{3r/2}} w^{\eta} \, d\mu(x) \right)^{1/\eta} \leq C \inf_{B_{3r/2}} u.
\end{equation*}
This concludes the proof. 
\end{proof}

% Moser Iteration

Next, we obtain a uniform estimate for the solutions of \eqref{eq:transformed_equation} using the Moser iteration technique. For this we require the following two lemmas



\begin{lemma}\label{lem:moser_iter_lem1}
 If $w$ satisfies 
 \begin{equation*}
      -\mathrm{div}\left( \frac{\nabla w(x)}{|x|^{2 d}}\right) + \varepsilon C_0 \frac{w}{|x|^{2s + 2d}} + \varepsilon (-\Delta_{d})^s w  = F(x,w(x))
 \end{equation*}
  where $F \in L^q(\om)$ for $q > N/2$, then $w \in L^{\infty}(\om)$.   
\end{lemma}
\begin{proof}
For \( 0 < \varepsilon < 1 \), consider the regularized function $f_\varepsilon(t) = \left( \varepsilon^2 + t^2 \right)^{1/2}$. Now, take the test function
\begin{equation*}
    \varphi = \psi \, f_\varepsilon'(w),
\end{equation*}
where \( \psi \in C_c^\infty(\Omega) \) is a positive function.
Then, using the weak formulation, we obtain
\begin{equation*}
\begin{aligned}
    \int_\Omega \frac{\nabla w  \nabla \varphi}{|x|^{2d}} \, dx 
    + \varepsilon C_0 \int_\Omega \frac{w \varphi}{|x|^{2s + 2d}} \, dx 
    + \varepsilon \iint_{\rtwon} (w(x) - w(y))(\varphi(x) - \varphi(y)) \, d\nu 
    \leq \int_\Omega |F| \left| f_\varepsilon'(w) \right| |\psi| \, dx.
\end{aligned}
\end{equation*}
By using Lemma \ref{ineq:A1} and invoking Fatou's Lemma, we deduce the inequality
\begin{equation*}
\begin{aligned}
    \int_\Omega \frac{\nabla |w| \,  \nabla \psi}{|x|^{2d}} \, dx 
    + \varepsilon \iint_{\mathbb{R}^{2N}} \left( |w(x)| - |w(y)| \right) ( \psi(x) - \psi(y) ) \, d\nu 
    \leq \int_\Omega |F| \, \psi \, dx,
\end{aligned}
\end{equation*}
for all positive test functions \( \psi \in C_c^\infty(\Omega) \). Hence, by density, the above inequality holds for all $ \psi \in \dot{H}^{s,d}(\Omega) \cap \dot{H}^1(\Omega, |x|^{-2d})$. \\

Let us define \( w_M = \min\{ |w|, M \} \) for some \( M > 0 \). Let \( \beta > 0 \) and \( \delta > 0 \), and consider the test function
\begin{equation*}
    \psi = (w_M + \delta)^\beta - \delta^\beta.
\end{equation*}
Using Lemma \eqref{ineq:A3}, we obtain the following estimate
\begin{equation}
\begin{aligned}
    \frac{4\beta}{(\beta + 1)^2} \int_{\om}  | \grad  (w_M + \delta)^{\frac{\beta + 1}{2}}|^2 &\frac{dx}{|x|^{2d}} 
    \leq 
    \beta \int_\Omega (w_M + \delta)^{\beta - 1} |\nabla w_M|^2 \, \frac{dx}{|x|^{2d}} \\
    &\quad + \varepsilon \iint_{\mathbb{R}^{2N}} (|w(x)| - |w(y)|) \left[ (w_M(x) + \delta)^\beta - (w_M(y) + \delta)^\beta \right] \, d\nu \\
    &\leq \int_\Omega |F| (w_M + \delta)^\beta \, dx.
\end{aligned}
\label{eq:main_energy_estimate}
\end{equation}
Now, invoking the Sobolev inequality with weights \eqref{ineq:sobolev_inequality_weight} and Young's inequality
% \begin{equation*}
%     \left( a^{1/2} - b^{1/2} \right)^2 \geq \frac{a}{2} - b,
% \end{equation*}
% we derive
\begin{equation*}
\begin{aligned}
S\int_{\om}  | \grad  (w_M + \delta)^{\frac{\beta + 1}{2}}|^2 \frac{dx}{|x|^{2d}} 
&\geq 
\left( \int_\Omega \left[ (w_M + \delta)^{\frac{\beta + 1}{2}} - \delta^{\frac{\beta + 1}{2}} \right]^{2^*} \frac{dx}{|x|^{ 2^* d}} \right)^{\frac{2}{2^*}}, \\
    &\geq  \left( \int_\Omega \left( \frac{(w_M + \delta)^{\beta + 1}}{2} - \delta^{{\beta + 1}} \right)^{\frac{2^*}{2}} \frac{dx}{|x|^{2^* d}} \right)^{\frac{2}{2^*}} \\
    &\geq \frac{1}{2} \left( \int_\Omega (w_M + \delta)^{\frac{(\beta + 1) 2^*}{2}} \frac{dx}{|x|^{2^* d}} \right)^{\frac{2}{2^*}} - \delta^{{\beta + 1}} \left( \int_\Omega \frac{dx}{|x|^{ 2^* d}} \right)^{\frac{2}{2^*}},
\end{aligned}
\end{equation*}
where \( S \) is the best constant in the Sobolev embedding.
Substituting the above estimate into inequality \eqref{eq:main_energy_estimate}, we get
\begin{equation}
\begin{aligned}
    \frac{\delta}{2} \left( \int_\Omega (w_M + \delta)^{\frac{\beta 2^*}{2}} \frac{dx}{|x|^{2^* d}} \right)^{\frac{2}{2^*}}
    &\leq \frac{1}{2} \left( \int_\Omega (w_M + \delta)^{\frac{(\beta + 1) 2^*}{2}} \frac{dx}{|x|^{2^* d}} \right)^{2/2^*} \\
    &\leq \delta^{\beta + 1} \left( \int_\Omega \frac{dx}{|x|^{2^* d}} \right)^{\frac{2}{2^*}} + \frac{(\beta + 1)^2}{4\beta S} \int_\Omega |F| (w_M + \delta)^\beta dx.
\end{aligned}
\label{eq:improved_estimate_main}
\end{equation}
Moreover, we also have the following inequality
\begin{equation}
\begin{aligned}
    \delta^\beta \left( \int_\Omega \frac{dx}{|x|^{2^* d}} \right)^{2/2^*}
    &\leq \frac{1}{\beta} \left( \frac{\beta + 1}{2} \right)^2 
    \left( \int_\Omega \frac{dx}{|x|^{2^* d}} \right)^{\frac{2}{2^*} - \frac{1}{q'}} 
    \left( \int_\Omega \frac{(w_M + \delta)^{\beta q'}}{|x|^{2^* d}} dx \right)^{1/q'}
\end{aligned}
\label{eq:holder_embedding_estimate}
\end{equation}
where the exponents satisfy the usual duality relation $\frac{1}{q} + \frac{1}{q'} = 1$.

Substituting the estimate \eqref{eq:holder_embedding_estimate} into \eqref{eq:improved_estimate_main}, we deduce the following refined inequality
\begin{equation*}
\begin{aligned}
    \left( \int_\Omega (w_M + \delta)^{\frac{\beta 2^*}{2}} \frac{dx}{|x|^{2^* d}} \right)^{\frac{2}{2^*}}
    &\leq C \left( \frac{\beta + 1}{2} \right)^2 \frac{1}{\beta}
    \left[ \left( \int_\Omega \frac{(w_M + \delta)^{\beta q'}}{|x|^{2^* d}} dx \right)^{\frac{1}{q'}} \right. \\
    &\quad \left. \left( \int_\Omega \frac{dx}{|x|^{2^* d}} \right)^{\frac{2}{2^*} - \frac{1}{q'}} + \frac{2}{\delta} \left\| F  |x|^{2^* d/q'} \right\|_{L^q} \right]
\end{aligned}
\end{equation*}
We now choose the parameter \( \delta \) to balance the terms as follows
\begin{equation*}
    \delta = 2 \left\| F  |x|^{2^* d/q'} \right\|_{L^q} 
    \left( \int_\Omega \frac{dx}{|x|^{2^* d}} \right)^{\frac{1}{q'} - \frac{2}{2^*}}
\end{equation*}
Letting \( \ell = \beta q' \) and $ \chi = \frac{N}{N - 2}  \frac{1}{q'} > 1$, we obtain
\begin{equation*}
    \left( \int_\Omega (w_M + \delta)^{\chi \ell} \frac{dx}{|x|^{2^* d}} \right)^{\frac{1}{\chi \ell}}
    \leq \left( C \left[ \int_\Omega \frac{dx}{|x|^{2^* d}} \right]^{\frac{2}{2^*} - \frac{1}{q'}} \right)^{\frac{q'}{\ell}}
    \left( \frac{q'}{\ell} \right)^{\frac{q'}{\ell}}
    \left( \frac{q' + \ell}{2\ell} \right)^{\frac{2}{\ell}}  \left( \int_\Omega \frac{(w_M + \delta)^{\ell}}{|x|^{2^* d}} dx \right)^{\frac{1}{\ell}}.
\end{equation*}
Now, applying the Moser-type iteration as detailed in \cite[Theorem $3.1$]{brasco_second_eigenvalue}, we obtain the desired result.
\end{proof}

The following lemma is motivated from Theorem $3.3$ of \cite{squassina_regularity}.
\begin{lemma}\label{lem:moser_iter_lem2}
If $w$ satisfies equation \eqref{eq:transformed_equation} then $w \in L^q(\om)$ for any $q \geq 1$.    
\end{lemma}
\begin{proof}
Let $g_\beta(t) = \text{sgn}(t)\,|t|\,|t_k|^\beta$, where $\beta > 0$,  and $t_k = \min\{t, k\}$. Also, define the auxiliary function
\begin{equation*}
G_\beta(t) = \int_0^t g_\beta'(t)^{1/2} \, dt \geq \frac{2(\beta + 1)^{1/2}}{\beta + 2} g_{\beta/2}(t).
\end{equation*}
Then, by Lemma \ref{ineq:A2}, we have
\begin{equation*}
\langle (-\Delta)^s_d w, g_\beta(w) \rangle_{s,d} \geq \left[ G_\beta(w) \right]_{s,d}^2.
\end{equation*}
Moreover, we have the following identity
\begin{equation*}
\int_\Omega \frac{w |w|\, |w_k|^\beta \, \text{sgn}(w)}{|x|^{2s + 2 d}}\, dx = \int_\Omega \frac{|w|^2\, |w_k|^\beta}{|x|^{2s + 2 d}}\, dx \geq 0.
\end{equation*}
Finally, considering the energy estimate with \( g_\beta(w) \), we compute
\begin{equation*}
\begin{aligned}
\int_\Omega \nabla w  \nabla \big( g_\beta(w) \big)\, d\mu 
&= \int_\Omega |\nabla w|^2 |w_k|^\beta\, d\mu + \beta \int_\Omega |w_k|^\beta |\nabla w_k|^2 \, d\mu \\
&\geq \left(1 + \frac{\beta}{4}\right)^{-1} \int_\Omega \left( |\nabla w|^2 |w_k|^\beta + \left( \frac{\beta^2}{4} + \beta \right) |w_k|^\beta |\nabla w_k|^2 \right) d\mu \\
&= \left(1 + \frac{\beta}{4} \right)^{-1} \int_\Omega |\nabla (w |w_k|^{\beta/2})|^2\, d\mu.
\end{aligned}
\end{equation*}
Combining all previous results and using \( g_\beta(w) \) as a test function along with Sobolev inequality, we deduce
\begin{equation}\label{eq:moser_estimate1}
\begin{aligned}
S \left( 1 + \frac{\beta}{4} \right)^{-1}
\left( \int_\Omega \frac{|w|^{2^*} |w_k|^{\beta \frac{2^*}{2}}}{|x|^{2^* d}} \, dx \right)^{\frac{2}{2^*}}
&\leq \int_\Omega \nabla w  \nabla (g_\beta(w)) \, d\mu + \varepsilon \left[ G_\beta(w) \right]_{s,d}^2 + \varepsilon C_0\int_\Omega \frac{|w|^2 |w_k|^\beta}{|x|^{2s + 2 d}}\, dx \\
&\leq \int_\Omega \frac{|w|^{2^*} |w_k|^\beta}{|x|^{2^* d}} \, dx 
+ \lambda \int_\Omega \frac{|w|^{q+1} |w_k|^\beta}{|x|^{(p+1)d}}\, dx \\
&= I_1 + I_2.
\end{aligned}
\end{equation}
We now estimate each term separately. First, we handle \( I_1 \)
\begin{equation}\label{eq:I_1_estimate_moser}
\begin{aligned}
I_1 &\leq \int_{\{|w| < K\}} \frac{|w|^{2^*} |w_k|^\beta}{|x|^{2^* d}} \, dx 
+ \int_{\{|w| \geq K\}} \frac{|w|^{2^*} |w_k|^\beta}{|x|^{2^*d}} \, dx \\
&\leq K^\beta \int_{\{|w| < K\}} \frac{|w|^{2^*}}{|x|^{2^* d}} \, dx 
+ \left( \int_{\{|w| > K\}} \frac{|w|^{2^*}}{|x|^{2^* d}} \, dx \right)^{1 - \frac{2}{2^*}} 
\left( \int_{\{|w| > K\}} \frac{|w|^{2^*}|w_k|^{\beta 2^*/2}}{|x|^{2^* d}} \, dx \right)^{\frac{2}{2^*}}.
\end{aligned}
\end{equation}

Similarly, for \( I_2 \), we obtain
\begin{equation}\label{eq:I_2_estimate_moser}
\begin{aligned}
I_2 &\leq K^\beta \int_{\{|w| < K\}} \frac{|w|^{p + 1}}{|x|^{(p+1)d}} \, dx 
+ \left( \int_{\{|w| > K\}} \frac{|w_k|^{\beta 2^*/2}}{|x|^{2^* d}} \, dx \right)^{2/2^*}
\left( \int_{\{|w| > K\}} \left(\frac{|w|^{q + 1}}{|x|^{(q-1)d}} \right)^{N/2} \, dx \right)^{2/N}.
\end{aligned}
\end{equation}
Choosing \( K \) sufficiently large so that the terms in \eqref{eq:I_1_estimate_moser} and \eqref{eq:I_2_estimate_moser} are absorbed into the right-hand side of \eqref{eq:moser_estimate1}, we deduce the following estimate
\begin{equation}
\left( \int_{\Omega} \frac{|w|^{2^*} |w_k|^{\beta \frac{2^*}{2}}}{|x|^{2^* d}} \, dx \right)^{\frac{2}{2^*}}
\leq 
C \left[
K^\beta \int_{\Omega} \frac{|w|^{2^*}}{|x|^{2^* d}} \, dx 
+ K^\beta \int_{\om} \frac{|w|^{p + 1}}{|x|^{(p+1)d}} \, dx 
\right].
\label{eq:estimate_absorbed}
\end{equation}
Now, applying Fatou's Lemma to the inequality \eqref{eq:estimate_absorbed}, we conclude the following for any \( \beta \geq 1 \)

\begin{equation*}
\int_{\Omega} \frac{|w|^{2^*(1 + \beta/2)}}{|x|^{2^* d}} \, dx < \infty. \qedhere
\end{equation*}
\end{proof}

\begin{theorem}[Uniform Boundedness] \label{thm:uniform_boundedness}
If \( w \) is a positive solution of \eqref{eq:transformed_equation}, then $w \in L^\infty(\om)$.
\end{theorem}
\begin{proof}
The proof follows directly from Lemmas \ref{lem:moser_iter_lem1} and \ref{lem:moser_iter_lem2}.
\end{proof}
% \begin{proof}
% \textcolor{red}{change the proof}From the Harnack inequality, we know that there exists \( r_0 > 0 \) such that $w(x) \geq C_0 > 0$ in $B_{r_0}$. Now, consider the test function
% \begin{equation*}
% \Psi = \eta^2   w   \min \left\{ w^{2\tau}, \ell^2 \right\}
% \end{equation*}
% where \( \eta \) is a cut-off function in \( B_{r}(0) \) with \( r \leq r_0 \), and \( \tau \) is a positive parameter.

% Using the inequality (see reference \eqref{fractional_minimum_inequality}), we get
% \begin{equation*}
% \begin{split}
% (w(x) - &w(y)) (\Psi(x) - \Psi(y)) \\
% & = \begin{cases}
%   (w(x) - w(y)) \left( \eta^2(x) w(x) - \eta^2(y) w(y) \right) \ell^2, & w^{2\tau}(x) > \ell^2, \ w^{2\tau}(y) > \ell^2  \\
%   (w(x) - w(y)) \left( \eta^2(x) w^{2\tau+1}(x) - \eta^2(y) w(y)\ell^2 \right), \quad& w^{2\tau}(x) < \ell^2, \ w^{2\tau}(y) > \ell^2\\
%   (w(x) - w(y)) \left( \eta^2(x) w(x)\ell^2 - \eta^2(y) w^{2\tau +1}(y) \right), \quad& w^{2\tau}(x) > \ell^2, \ w^{2\tau}(y) < \ell^2\\
%   (w(x) - w(y)) \left( \eta^2(x) w^{2\tau+1}(x) - \eta^2(y) w^{2\tau+1}(y) \right), \quad& w^{2\tau}(x)< \ell^2, w^{2\tau}(y) < \ell^2
% \end{cases}\\
% & \geq \begin{cases}
%     \frac{4\tau \ell^2}{(\tau+1)^2}   \left| (\eta^2(x)w(x))^{(1+\tau)/2} - (\eta^2(y) w(y))^{(1+\tau)/2} \right|^2, \quad& w^{2\tau}(x) > \ell^2, w^{2\tau}(y) > \ell^2 \\
%     \ell^2(w(y) - w(x))(\eta^2(y) w(y) - \eta^2(x) w(x)), & w^{2 \tau} (x) < \ell^2, \ w^{2\tau}(y) > \ell^2\\
%     \ell^2(w(x) - w(y))(\eta^2(x) w(x) - \eta^2(y) w(y)), & w^{2 \tau} (x) > \ell^2, \ w^{2\tau}(y) < \ell^2\\
%     \frac{(2\tau +1)^2}{(\tau +1)^2} \left| \eta^{\frac{2(\tau +1)}{2 \tau +1}}(x) w^{\tau +1}(x) - \eta^{\frac{2(\tau +1)}{2 \tau +1}}(y) w^{\tau +1}(y) \right|^2, & w^{2\tau}(x)< \ell^2, w^{2\tau}(y) < \ell^2
% \end{cases}\\
% & \geq0.
% \end{split}
% \end{equation*}

% Therefore, we obtain
% \begin{align*} 
% \int_{\Omega}  \nabla w   \nabla \Psi \frac{dx}{|x|^{2 d}}
% &\leq \int_{\Omega}  \nabla w   \nabla \Psi \frac{dx}{|x|^{2d}}
% + \varepsilon C_0 \int_{\Omega} \frac{ w \Psi}{|x|^{2s}} \frac{dx}{|x|^{2d}}
% + \varepsilon  \langle  w, \Psi \rangle_{s,d} \\
% &\leq \int_{\Omega}  w^{2^*-1} \Psi \frac{dx}{|x|^{2^*d}} + \lambda \int_{\Omega} w^{q} \Psi \frac{dx}{|x|^{(q+1)d}}.
% \end{align*}
% Now, the $L^{\infty}$-estimate follows from the Proposition $3.2$ of \cite{chen2004multiple}.
% \end{proof}


\begin{proof}[Proof of Theorem \ref{thm:asympototic_estimates}]
 The proof follows directly from Theorems \ref{thm:harnack_inequality_transformed} and \ref{thm:uniform_boundedness}.   
\end{proof}

Finally, we have the following strong maximum principle. 
\begin{theorem}\label{thm:uniform_estimate_and_smp}
Assume that there exists a nonnegative solution $w_0 \in \X(\Omega)$ to the problem \eqref{eq:main_problem_ep} for some $\lambda \geq 0$ and $p \in [1,2^*-1)$. Then $w_0 >0$ almost everywhere in $\Omega$.
\end{theorem}
\begin{proof}
Consider the inner product
\begin{equation*}
\langle w_0, \varphi \rangle + \varepsilon\langle w_0, \varphi \rangle_{s} = \int_{\Omega} \frac{w_0}{|x|^2} \varphi + \int_{\Omega} |w_0|^{2^*-2} w_0 \varphi + \lambda \int_{\Omega} |w_0|^{p-2} w_0 \varphi \geq 0.
\end{equation*}
Since $w_0 \geq 0$, applying Theorem 8.4 of \cite{garain_higher_holder_regularity}, we conclude that $w_0 > 0$.
\end{proof}

%%%%%%%%%%%%%%%%%%%%%%%%%%%%%%%%%%%%%%%%%%%%%%%%%%%%%%%%%%%%%%%%%%
%%%%%%%%%%%%%%%%    End of Asyptotic estimates            %%%%%%%%
%%%%%%%%%%%%%%%%%%%%%%%%%%%%%%%%%%%%%%%%%%%%%%%%%%%%%%%%%%%%%%%%%%

%%%%%%%%%%%%%%%%%%%%%%%%%%%%%%%%%%%%%%%%%%%%%%%%%%%%%
%%%%%%% Linear Problem
%%%%%%%%%%%%%%%%%%%%%%%%%%%%%%%%%%%%%%%%%%%%%%%%%%%%%
\section{Linear and Superlinear Case}\label{sec:linear_sublinear}
\subsection{Linear Case (\texorpdfstring{$p = 2$})}\label{sec:linear}
In this subsection, we consider the case of linear perturbation in problem \eqref{eq:main_problem}. The analysis relies heavily on the spectral properties of the underlying operator, particularly the behaviour of its first eigenvalue. These properties play a crucial role in establishing the existence of solutions.

% \textbf{Proof of Theorem \ref{thm:nonexistence_star_shaped}}
% Suppose, on the contrary, that a solution $u_0$ exists for \eqref{eq:main_problem_ep}. Then, by Theorem \ref{thm:uniform_estimate_and_smp}, we obtain:
% \begin{equation*}
% u_0 \notin L^\infty_{\text{loc}}(\mathbb{R}^N \setminus \{0\}).
% \end{equation*}

% By applying the \textcolor{red}{check this} **Serra-Miranda Pohozaev identity** (which provides **non-existence results of nontrivial bounded solutions**), we conclude that a contradiction arises. Thus, no such solution $u_0$ can exist.

%%%%%%%%%%%%{Spectral Properties and Further Existence Results}
Let $V(\om) = \left\{ u \in \X(\Omega) \mid \int_{\Omega} |u|^{2^*} dx= 1 \right\}$. Then we take the restriction of the functional $\Je$ on $V(\om)$ defined as 
\begin{equation*}
    \Q(u) := \int_{\Omega} |\nabla u|^2dx - \mu \int_{\Omega} \frac{|u|^2}{|x|^2} dx + [u]_s^2 - \lambda \int_{\Omega} |u|^2dx.
\end{equation*}
Now define
\begin{equation*}
S_{\mu}(\lambda) := \inf_{u \in V(\om)} \Q(u).
% \label{eq:infimum_functional}
\end{equation*}
% Let $\lambda_{1, s}$ and $\la_1$ denote the first eigenvalues defined in \eqref{eq:first_eigenvalue} and \eqref{eq:first_eigenvalue_frac} respectively.
Let us state and prove some of the important properties of $S_{\mu}(\lambda)$.
\begin{lemma}[Properties of $S_{\mu}(\lambda)$]\label{lem:best_constant_being_constant}
The function $S_{\mu}(\lambda)$ satisfies the following properties
\begin{enumerate}
    \item It satisfies the lower bound
    \begin{equation*}
    S_{\mu}(\lambda) \geq S_{\mu} - \lambda |\Omega|^{2/N}.
    \end{equation*}
    % Let $u \in V(\om)$, then 
    % $$\mathcal{Q}(u) \geq \inf_{u \in V(\om)} \left( \rho(u)^2 - \mu \int_{\om} \frac{|u|^2}{|x|^2} \right) - \la \int_{\om} |u|^2 \geq S_{\mu} - \la |\om|^{2/N}.$$
    
    \item For every $\lambda > 0$, we have $S_{\mu}(\lambda) \leq S_{\mu}.$

    \item If $0 < t_2 < t_{1}$, then $ S_{\mu}(t_2) \geq S_{\mu}(t_1)$.
    \item $S_{\mu}(\la) \geq 0$ if and only if $0 < \la \leq \la_1$, where $\la_1$ is the first eigenvalue defined in \eqref{eq:first_eigenvalue}.
    \item For $0 < \lambda \leq \lambda_{1, s}$, we have $S_{\mu}(\lambda) = S_{\mu} > 0$, where $\la_{1,s}$ is the first eigenvalue defined in \eqref{eq:first_eigenvalue_frac}.
    \item The function $\lambda \mapsto S_{\mu}(\lambda)$ is continuous on $(0, \infty)$.
\end{enumerate}
\end{lemma}
\begin{proof}
Part $1 - 4$ follows directly. \\
Part $5$:  For $u \in  V(\om)$, we estimate
\begin{equation*}
\Q(u) \geq \int_{\Omega} |\nabla u|^2 - \mu \int_{\Omega} \frac{|u|^2}{|x|^2} dx + \lambda_{1,s} \int_{\Omega} |u|^2 - \lambda \int_{\Omega} |u|^2 \geq \int_{\Omega} |\nabla u|^2 - \mu \int_{\Omega} \frac{|u|^2}{|x|^2}d x.
\end{equation*}
This implies $S_{\mu}(\lambda) \geq  S_{\mu}.$ Using this with part $2$, we conclude $S_{\mu}(\lambda) = S_{\mu}$ for $0 < \lambda \leq \lambda_{1, s}$.\\
Finally part $6$ follows similarly to Lemma $4.8$ in \cite{biagi2022brezis}.
\end{proof}

% We recall that, given the bounded open set $\Omega$, one defines
%  \begin{itemize}
%   \item[(i)] the first Dirichlet eigenvalue of $(-\Delta)^s$ in $\Omega$ as
%   \begin{equation} \label{eq:deflambda1s}
%    \lambda_{1,s} := 
%   \inf\big\{[u]^2_s:\,\text{$u\in C_0^\infty(\Omega)$ and $\|u\|_{L^2(\rnn)} = 1$}\big\};
%   \end{equation}
%   \item[(ii)] the first Dirichlet eigenvalue of $\mathcal{L}$ in $\Omega$ as
%   \begin{equation} \label{eq:deflambda1}
%    \lambda_{1} := 
%   \inf\big\{\int_{\om} |\grad u|^2 dx + [u]_s^2:\,\text{$u\in C_0^\infty(\Omega)$ and $\|u\|_{L^2(\rnn)} = 1$}\big\}.
%   \end{equation}
%  \end{itemize}
%  Clearly, both $\lambda_{1,s}$ and $\lambda_1$ depend on the open set $\Omega$;
%  however, in order to simplify the notation, we avoid to make explicit this dependence
%  in what follows.
 
%  For the sake
%  of completeness (and for a future reference), 
%  we then present in the next remark a short
%  overview of the main properties
%  of $\lambda_{1,s}$ and $\lambda_{1}$.
%  \begin{remark}[Properties of $\lambda_{1,s}$ and $\lambda_{1}$] \label{rem:problambdalambdas}
%  As regards $\lambda_{1,s}$ we first
%   observe that, since the map 
%   $u\mapsto [u]_s =: \mathcal{N}(u)$ is
%  a norm on $C^\infty_0(\Omega)$
%  which is equivalent to the full $H^s$-norm
%  (see, e.g., \cite[Theorem~6.5]{di2012hitchhiker} and recall that $\Omega$ is bounded), one has
%  $$\lambda_{1,s} = 
%  \inf\big\{[u]^2_s:\,\text{$u\in \mathcal{D}_0^{s,2}(\Omega)$ and $\|u\|_{L^2(\rnn)} = 1$}\big\},$$
%  where $\mathcal{D}_0^{s,2}(\Omega)\subseteq L^2(\rnn)$ is the completion of
%  $C_0^{\infty}(\Omega)$ with respect to $\mathcal{N}$.
%  Moreover, since the embedding $\mathcal{D}_0^{s,2}(\Omega)\hookrightarrow L^2(\rnn)$
%  is compact, it is not difficult to recognize that $\lambda_{1,s}$
%  is \emph{actually achieved} in this larger space $\mathcal{D}_0^{s,2}(\Omega)$, that is,
%  $$\exists\,\,\varphi_0\in\mathcal{D}_0^{s,2}(\Omega):\,\text{$\|\varphi_0\|_{L^2(\rnn)} = 1$
%  and $[\varphi_0]^2 = \lambda_{1,s} > 0$}.$$
%  As a matter of fact, the above function $\varphi_0$ can be chosen to be \emph{strictly positive
%  a.e.\,in $\Omega$} (see, e.g., \cite[Proposition~9]{servadei2013variational}); 
%  in addition, since $\varphi_0$ is a constrained minimizer of the functional $u\mapsto [u]_s^2$,
%  by the Lagrange Multiplier Rule we easily see that
%  $$\iint_{\R^{2n}}\frac{(\varphi_0(x)-\varphi_0(y))(v(x)-v(y))}{|x-y|^{n+2s}}\,dx\,dy
%  = \lambda_{1,s}\int_{\Omega}\varphi_0v\,d x\qquad\forall\,\,v\in\mathcal{D}_0^{s,2}(\Omega),$$
%  and this shows that $\varphi_0$ is weak solution to the \emph{eigenvalue problem}
%  \begin{equation*}
%   \begin{cases}
%  (-\Delta)^su = \lambda_{1,s} u & \text{in $\Omega$}, \\
%  u \not\equiv 0 & \text{in $\Omega$}, \\
%  u\equiv 0 & \text{in $\rnn\setminus\Omega$}.
%  \end{cases}
%  \end{equation*}
%  As regards $\lambda_{1}$, we have completely analogous
%  results (with the obvious modifications): first of all, since $\rho( )$
%  defines
%  a norm on $C^\infty_0(\Omega)$
%  which is e\-qui\-va\-lent to the full $H^1$-norm in $\rnn$
%  (as $\Omega$ is bounded), we have
%  $$\lambda_{1} = 
%  \inf\big\{\int_{\om} |\grad u|^2 dx + [u]_s^2:\,\text{$u\in \X(\Omega)$ and $\|u\|_{L^2(\rnn)} = 1$}\big\},$$
%  Moreover, since the embedding $\X^(\Omega)\hookrightarrow L^2(\rnn)$
%  is compact, it is not difficult to see that $\lambda_{1}(\Omega)$
%  is \emph{actually achieved} in this larger space $\X(\Omega)$, that is,
%  $$\exists\,\,\psi_0\in\X(\Omega):\,\text{$\|\psi_0\|_{L^2(\rnn)} = 1$
%  and $\rho(\psi_0)^2 = \lambda_{1} > 0$}.$$
%  As a matter of fact, the above function $\psi_0$ can be chosen to be \emph{strictly positive
%  a.e.\,in $\Omega$}
%  (see, e.g., \cite[Proposition~5.1]{di2012hitchhiker}); in addition, 
%  since $\psi_0$ is a constrained minimizer of the functional
%  $\rho( )^2$, by the Lagrange Multiplier Rule we easily see that
%  \begin{align*}
%   & \int_{\Omega}\nabla \psi_0 \nabla v\,dx 
%   + \iint_{\R^{2n}}\frac{(\psi_0(x)-\psi_0(y))(v(x)-v(y))}{|x-y|^{n+2s}}\,dx\,dy \\
%   & \qquad\quad = \lambda_{1}\int_{\Omega}\psi_0v\,d x
%   \qquad \forall\,\,v\in\X(\Omega),
%   \end{align*}
%  and this shows that $\psi_0$ is weak solution to the \emph{eigenvalue problem}
%  \begin{equation*}
%   \begin{cases}
%  \mathcal{L} u = \lambda_{1} u & \text{in $\Omega$}, \\
%  u \not\equiv 0 & \text{in $\Omega$}, \\
%  u\equiv 0 & \text{in $\rnn\setminus\Omega$}.
%  \end{cases}
%  \end{equation*}
%  \end{remark}




% Before establishing Theorem \ref{thm:existence_linear_case} rigorously, we require the following lemma

% \begin{lemma}\label{lem:best_constant_being_constant}
% For $0 < \lambda \leq \lambda_{1, s}$, we have $S_{\mu}(\lambda) = S_{\mu} > 0$.
% \end{lemma}
% \begin{proof}
% Since $\lambda > 0$, we have $S_{\mu}(\lambda) \leq S_{\mu}.$ For $u \in C_c^\infty(\Omega) \cap S(\om)$, we estimate
% \begin{equation*}
% \Q(u) \geq \int_{\Omega} |\nabla u|^2 - \mu \int_{\Omega} \frac{|u|^2}{|x|^2} dx + \lambda_{1,s} \int_{\Omega} |u|^2 - \lambda \int_{\Omega} |u|^2 \geq \int_{\Omega} |\nabla u|^2 - \mu \int_{\Omega} \frac{|u|^2}{|x|^2}d x.
% \end{equation*}
% This implies $S_{\mu}(\lambda) \geq  S_{\mu}.$ Thus, for $0 < \lambda \leq \lambda_{1, s}$, we conclude $S_{\mu}(\lambda) = S_{\mu}$.
% \end{proof}

%%%%%%%{Existence of Solutions}

% \begin{lemma}\label{lem:existence_of_soln}
% Suppose that $S_{\mu}(\lambda)$, defined as in \eqref{eq:infimum_functional}, satisfies $S_{\mu}(\lambda) > 0$ and is achieved by some function $w$. Then, the problem \eqref{eq:main_problem_ep} admits a solution.
% \end{lemma}
% \begin{proof}
% Let $w \in V(\Omega)$ be a minimizer of the functional $\Q$, i.e., $\Q(w) = S_{\mu}(\lambda)$.
% Since, $\Q(|w|) \leq  \Q(w)$. Without loss of generality, we assume $w \geq 0$ almost everywhere in $\Omega$.
% By the method of Lagrange multipliers, there exists a scalar $\theta$ such that
% \begin{equation*}
% \nabla w  \nabla \varphi + \langle w, \varphi \rangle_{s} - \mu \int_{\Omega} \frac{w \varphi}{|x|^2}  - \lambda \int_{\Omega} w \varphi = \theta \int_{\Omega} |w|^{2^*-2}w \varphi, \quad \forall \varphi \in \X(\Omega).
% \end{equation*}
% Choosing $\varphi = w$ gives $\theta = S_{\mu}(\lambda)$.
% Now, letting $u = S_{\mu}(\lambda)^{(N-2)/4}w$, we obtain
% \begin{equation*}
% \int_{\Omega} \nabla u \nabla \varphi + \langle u, \varphi \rangle_{s} - \mu \int_{\Omega} \frac{u \varphi}{|x|^2} \varphi = \int_{\Omega} |u|^{2^*-2} u \varphi + \lambda \int_{\Omega} u \varphi.
% \end{equation*}
% Thus, $u$ is a solution to the problem \eqref{eq:main_problem_ep}.
% \end{proof}

%%%%%%%%{Analysis of \texorpdfstring{$S_{\mu}(\lambda)$}{S(B)(lambda)}}

% \begin{lemma}
% The function $\lambda \mapsto S_{\mu}(\lambda)$ is continuous on $(0, \infty)$.
% \end{lemma}
% \begin{proof}
% This result follows similarly to Lemma $4.8$ in \cite{biagi2022brezis}.    
% \end{proof}


\begin{proof}[Proof of Theorem \ref{thm:existence_linear_case}]
Define
\begin{equation*}
\lambda^* = \sup \left\{ \lambda_0 > 0 \mid S_{\mu}(\lambda) = S_{\mu} \mbox{ for all } 0 < \la \leq \la_0 \right\}.
\end{equation*}
By part $5$ of Lemma \ref{lem:best_constant_being_constant}, it is evident that $\lambda^* \geq \lambda_{1, s}$. Furthermore, since $S_{\mu}(\lambda)$ is a continuous function, we obtain $S_{\mu}(\lambda^*) = S_{\mu}$. Moreover, from the properties of $S_{\mu}(\lambda)$, we deduce
\begin{equation*}
S_{\mu}(\lambda) 
\begin{cases}
    \geq 0, & \text{if } \lambda \leq \lambda_1, \\
    \leq 0, & \text{if } \lambda > \lambda_1.
\end{cases}
\end{equation*}
This implies $S_{\mu}(\lambda_1) = 0$. Thus, we conclude $\lambda^* \in [\lambda_{1, s}, \lambda_1)$.

\noi \textbf{Case 1:} $0 < \lambda \leq \lambda_{1, s}$ \\
In this range, there does not exist a solution to the problem $( \mathcal{P}_{\mu,\la})$ within the closed ball $\mathcal{B}$ defined in \eqref{eq:linear_ball_NE}. On the contrary, $u$ be the solution to $( \mathcal{P}_{\mu,\la})$. Let $ v = \frac{u}{\| u \|_{L^{2^*}(\Omega)}}$. Then
\begin{equation*}
\Q(v)  = \frac{1}{\| u \|_{L^{2^*}}^2} \left( \int_{\om} |\grad u|^2 dx + [u]_s^2 - \mu \int_{\Omega} \frac{u^2}{|x|^2} dx - \lambda \| u \|_{L^2}^2 \right) = \frac{1}{\| u \|_{L^{2^*}}^2} \| u \|_{L^{2^*}}^{2^*}\leq S_{\mu}.
\end{equation*}
However, by part $5$ of Lemma \ref{lem:best_constant_being_constant} yield $\Q(u) \geq S_{\mu}.$ Thus, we conclude that $\Q(v) = S_{\mu}$. Further, we analyze
\begin{equation*}
S_{\mu}  \leq \| \nabla v \|_{L^2}^2 - \mu \int_{\Omega} \frac{v^2}{|x|^2} = \Q(v) - \left( [v]^2_{s} - \lambda \| v \|_{L^2}^{2} \right)\leq \Q(v)-\left( \lambda_{1,s} - \lambda \right) \| v \|_{L^2}^{2} \leq \Q(v) \leq S_{\mu}.
\end{equation*}
But $S_{\mu}$ is never achieved in $\Omega \subsetneq \mathbb{R}^N$ (bounded domain), we arrive at a contradiction.

\noi \textbf{Case 2:} $\lambda^* < \lambda < \lambda_1$ \\
For this range, by definition of $\la^*$, we have $0 \leq S_{\mu}(\lambda) < S_{\mu}$. By Lemma 1.2 of \cite{brezis_1983}, $S_{\mu}(\la)$ is achieved. Thus, there exists a nonzero function $w \in V(\Omega)$ such that $\Q(w) = S_{\mu}(\lambda)$. Since $\lambda < \lambda_1$, we establish $S_{\mu}(\lambda)  \geq (\lambda_1 - \lambda) \| w \|_{L^2}^2 > 0$. Since, $\Q(|w|) \leq  \Q(w)$. Without loss of generality, we assume $w \geq 0$ almost everywhere in $\Omega$.
Applying the method of Lagrange multipliers, there exists a scalar $\theta$ such that
\begin{equation*}
\nabla w  \nabla \varphi + \langle w, \varphi \rangle_{s} - \mu \int_{\Omega} \frac{w \varphi}{|x|^2}  - \lambda \int_{\Omega} w \varphi = \theta \int_{\Omega} |w|^{2^*-2}w \varphi, \quad \forall \varphi \in \X(\Omega).
\end{equation*}
Choosing $\varphi = w$ gives $\theta = S_{\mu}(\lambda)$.
Now, letting $u = S_{\mu}(\lambda)^{(N-2)/4}w$, we obtain
\begin{equation*}
\int_{\Omega} \nabla u \nabla \varphi + \langle u, \varphi \rangle_{s} - \mu \int_{\Omega} \frac{u \varphi}{|x|^2} \varphi = \int_{\Omega} |u|^{2^*-2} u \varphi + \lambda \int_{\Omega} u \varphi.
\end{equation*}
Thus, $u$ is a solution to the problem $( \mathcal{P}_{\mu,\la})$.

% Applying Lemma \ref{lem:existence_of_soln}, we conclude that the solution exists.

\noi \textbf{Case 3:} $\lambda \geq \lambda_1$ where $\la_1$ is the eigenvalue of \eqref{eq:first_eigenvalue} corresponding to the eigenfunction $\phi_0$\\
Since $S_{\mu}(\lambda)$ is a non-increasing function, we obtain $S_{\mu}(\lambda) \leq 0 < S_{\mu}$. By Lemma 1.2 of \cite{brezis_1983}, $S_{\mu}(\lambda)$ is achieved.
Now, on contrary, let $u$ be a positive solution of $(\mathcal{P}_{\mu,\la})$. 
% We know that there exists an eigenvalue $\lambda_1 \in \mathbb{R}$ corresponding to the eigenfunction $\varphi_0$ satisfying the equation
% \begin{equation*}
% \int_{\Omega} \nabla \varphi_0   \nabla v + \langle \varphi_0, v \rangle_s - \mu \int_{\Omega} \frac{\varphi_0 v}{|x|^2} = \lambda_1 \int_{\Omega} \varphi_0 v, \quad \forall v \in \X(\Omega).
% \end{equation*}
Then using the strong maximum principle in Theorem \ref{thm:uniform_estimate_and_smp}, we write
\begin{equation*}
0 < \int_{\Omega} (u^{2^*-1} + \lambda u) \varphi_0 = \int_{\Omega} \nabla \varphi_0   \nabla u + \langle \varphi_0, u \rangle_s - \mu \int_{\Omega} \frac{\varphi_0 u}{|x|^2} =  \lambda_1 \int_{\Omega} \varphi_0 u.
\end{equation*}
So, we conclude that $\lambda_1 > \lambda$, which is a contradiction.
This completes the proof of Theorem \eqref{thm:existence_linear_case}.
\end{proof}

%%%%%%%%%%%%%%%%%%%%%%%%%%%%%%%%%%%%%%%%%%%%%%%%%%%%%%%%%%%%%%%%%%%%%%%%%%%%%%%%
%%%%%%%%%%%%%%%%%%%%%%%%%%%%%%%%%%%%%%%% HIGH ENERGY LINEAR PROBLEM 
%%%%%%%%%%%%%%%%%%%%%%%%%%%%%%%%%%%%%%%%%%%%%%%%%%%%%%%%%%%%%%%%%%%%%%%%%%%%%%%%



% \subsection{High Energy Solution }

% In this subsection, we prove the existence of high energy solutions with the help of the following theorem borrowed from \cite{cerami_multiplicity}.
% \begin{theorem}[Theorem $2.5$ of \cite{cerami_multiplicity}]\label{thm:high_energy_cerami}
%  Let $(H,\lv .\rv)$ be a Hilbert space and $\mathcal{I} \in C^1(H,\real)$ be a functional satifying the following properties
%  \begin{enumerate}
%      \item $I(u) = I(-u)$, $I(0) = 0$;
%      \item $I$ satisfies $(PS)$ condition in $(0,\beta)$ for some $\beta > 0$;
%      \item There exists two closed subspaces $V, W  \subset H$ and positive constants $\rho, \delta, \beta^{\prime}$, with $\delta < \beta^{\prime} < \beta$ such that 
%      \begin{enumerate}
%          \item $I(u) \leq \beta^{\prime}$ for any $u \in W$
%          \item $I(u) \geq \delta $ for any $u \in V, \, \lv u \rv = \rho$
%          \item codim $V < \infty$ and dim $ W \geq $ codim $V$. 
%      \end{enumerate}
%  \end{enumerate}
%  Then there exists at least $ \text{dim }W - \text{codim }V$ pair of critical points of $I$ with critical values belonging to the interval $[\delta, \beta^{\prime}]$.
% \end{theorem}
% \medskip
% Let us define 
% \begin{equation*}
% \lambda_+ := \min \{ \lambda_j \in \sigma_\mu : \lambda < \lambda_j \}
% \end{equation*}
% and the spaces
% \begin{equation*}
% M^+ = \overline{\bigoplus_{\lambda_j \geq \lambda_+} M(\lambda_j)}, \quad
% M^- = \bigoplus_{\lambda_j \leq \lambda_+} M(\lambda_j),
% \end{equation*}
% where $M(\lambda_j)$ denotes the eigenspace corresponding to the eigenvalue $\lambda_j$.
% \medskip

% \begin{lemma}
% Let $\la_+$, $M^+$, and $M^-$ be as above and assume that $\lambda_+ - \lambda < S_{\mu} |\Omega|^{-2/N}$. Then we have
% \begin{equation} \label{eq:sup-inequality}
% \beta_\lambda = \sup_{u \in M^-} \J(u) \leq (\lambda_+ - \lambda)^{N/2} \frac{|\Omega|}{N} < \frac{1}{N} S_{\mu}^{N/2}.
% \end{equation}
% Moreover, there exist constants $t_0 > 0$ and $\delta_\lambda \in (0, \beta_\lambda)$ such that
% \begin{equation*} 
% \J(u) \geq \delta_\lambda
% \quad \text{for all } u \in M^+ \text{ with } \|u\|_{\X} = t_0.
% \end{equation*}
% \end{lemma}
% \begin{proof}
% For any $u \in M^-$, we have
% \begin{equation*} 
% \int_{\om} |\grad u|^2 dx + [u]_s^2 - \mu \int_{\Omega} \frac{|u|^2}{|x|^2} \, dx \leq \lambda_+ \int_{\Omega} |u|^2 \, dx.
% \end{equation*}
% Applying H\"{o}lder's inequality and the definition of $M^-$, we obtain
% \begin{equation*}
% \J(u) \leq \frac{1}{2} (\lambda_+ - \lambda) |\Omega|^{2/N} \|u\|_{L^{2^*}}^2 - \frac{1}{2^*} \|u\|_{L^{2^*}}^{2^*},
% \end{equation*}
% where $2^* = \frac{2N}{N-2}$ is the critical Sobolev exponent. Define the auxiliary function
% \begin{equation*} 
% g(t) = \frac{1}{2} (\lambda_+ - \lambda) |\Omega|^{2/N} t^2 - \frac{1}{2^*} t^{2^*}.
% \end{equation*}
% Then, it follows that
% \begin{equation*} 
% \max_{t > 0} g(t) = \frac{1}{N} (\lambda_+ - \lambda)^{N/2} |\Omega| < \frac{1}{N} S_{\mu}^{N/2}.
% \end{equation*}
% Thus, from the choice of $\lambda$ as in \eqref{eq:lambda-range}, the inequality \eqref{eq:sup-inequality} holds.
% Next, for $u \in M^+$, we observe
% \begin{equation*} 
% \lambda_+ \int_{\om} |u|^2 dx \leq \int_{\om} |\grad u|^2 dx + [u]_s^2 - \mu \int_{\Omega} \frac{|u|^2}{|x|^2} dx.
% \end{equation*}
% Then using the above relation, we obtain
% \begin{equation*} 
% \J(u) \geq \frac{1}{2} \left(1 - \frac{\lambda}{\lambda_+}\right) \left[\int_{\om} |\grad u|^2 dx + [u]_s^2 - \mu \int_{\Omega} \frac{|u|^2}{|x|^2}\right] - \frac{1}{2^*} \int_{\Omega} |u|^{2^*},
% \end{equation*}
% which simplifies further to
% \begin{align*} 
% \J(u) & \geq \frac{1}{2} \left(1 - \frac{\lambda}{\lambda_+}\right) \left[\int_{\om} |\grad u|^2 dx + [u]_s^2  - \mu \int_{\Omega} \frac{|u|^2}{|x|^2}\right] \\
% & \quad- \frac{1}{2^*} S_{\mu}^{-2^*/2} \left[\int_{\om} |\grad u|^2 dx + [u]_s^2 - \mu \int_{\Omega} \frac{|u|^2}{|x|^2}\right]^{2^*/2} \\
% & \geq \frac{1}{2} \left(1 - \frac{\lambda}{\lambda_+}\right) \left( 1 - \frac{\mu}{\bar{\mu}} \right) \lv u \rv_{\X}^2  - \frac{1}{2^*} S_{\mu}^{-2^*/2} \lv u \rv_{\X}^{2^*}.
% \end{align*}
% Again define the auxiliary function $h:[0,\infty) \to \real$ as
% \begin{equation*}
% h(t) = \frac{1}{2} \left(1 - \frac{\lambda}{\lambda_+}\right) \left( 1 - \frac{\mu}{\bar{\mu}} \right) t^2 - \frac{1}{2^*} S_{\mu}^{-2^*/2} t^{2^*}.
% \end{equation*}
% Then, the function $h(t)$ achieves its maximum at $t = t_0$, where
% \begin{equation*} 
% t_0 = \left[\left(1 - \frac{\lambda}{\lambda_+}\right) \left( 1 - \frac{\mu}{\bar{\mu}} \right) S_{\mu}^{2^*/2}\right]^{(N-2)/4}.
% \end{equation*}
% Moreover,
% \begin{equation*}
% \max_{t \geq 0} h(t) = h(t_0) = \frac{1}{N} \left(1 - \frac{\lambda}{\lambda_+}\right)^{N/2} \left( 1 - \frac{\mu}{\bar{\mu}} \right)^{N/2}S_{\mu}^{N/2}.
% \end{equation*}
% Thus, we can conclude that
% \begin{equation*}
% \J(u) > \delta_\lambda \quad \text{for all} \quad u \in M^+ \cap \partial B_{t_0},
% \end{equation*}
% where $\delta_{\lambda} < h(t_0)$.
% Finally, we also have $\delta_\lambda < \beta_\lambda$. Indeed, since $M(\lambda_+) = M^+ \cap M^-$, it follows that
% \begin{equation*}
% M^+ \cap M^- \cap B_{t_0} \neq \emptyset,
% \end{equation*}
% and for any $u \in M^+ \cap M^- \cap B_{t_0}$, we have
% \begin{equation*}
% \delta_\lambda < \J(u) \leq \sup_{u \in M^-} \J(u) = \beta_\lambda. \qedhere
% \end{equation*}
% \end{proof}
% \medskip

% \begin{proof}[Proof of Theorem \ref{thm:high_energy_linear}]
% Using the Theorem \ref{thm:high_energy_cerami} with the identifications $H = \X(\om), V = M^+, W = M^-, \beta = S_{\mu}^{N/2}/N, \beta' = \beta_\lambda, \delta = \delta_\lambda, \rho = \beta_\lambda$ and in addition, we note that
% \begin{equation*}
% \dim V - \text{codim} \, W = \nu_k + 1,
% \end{equation*}
% where $\nu_k$ is the multiplicity of the eigenvalue $\lambda_k$. This completes our proof.
% \end{proof}



%%%%%%%%%%%%%%%%%%%%%%%%%%%%%%%%%%%%%%%%%%%%%%%%%%%%%
%%%%%%% Superlinear Problem
%%%%%%%%%%%%%%%%%%%%%%%%%%%%%%%%%%%%%%%%%%%%%%%%%%%%
\subsection{Superlinear Case (\texorpdfstring{$2 < p < 2^*$})}\label{sec:superlinear}

Let \( 2 < p < 2^* \). Since we are concerned with positive solutions of \eqref{eq:main_problem}, it is not necessary that \(|u|\) is also a solution whenever \( u \) is a solution. To address this issue, we consider the following modified functional
\begin{equation*}
\Jplus(u) = \frac{1}{2} \int_{\om} |\grad u|^2 + \frac{1}{2} [u]_s^2 - \frac{\mu}{2} \int_{\Omega} \frac{|u|^2}{|x|^2} dx 
- \frac{1}{2^*} \int_{\Omega} (u^+)^{2^*} - \frac{\lambda}{p} \int_{\Omega} (u^+)^{p}.
\end{equation*}
\noi Any critical point of the functional \( \Jplus \) is a solution of the following equation
\begin{equation}\tag{$\mathcal{P}_{\mu, \la}^{+}$}
\begin{cases}
-\Delta u + (-\Delta)^s u - \mu \frac{u}{|x|^2} = (u^+)^{2^*-1} + \lambda (u^+)^{p-1}, & \text{in } \Omega, \\
u = 0, & \text{in } \mathbb{R}^N \setminus \Omega.
\end{cases}
\label{eq:critical_equation}
\end{equation}
To link the solutions of equation \eqref{eq:critical_equation} with the solutions of problem \eqref{eq:main_problem}, we use the following Weak Maximum Principle.
%Theorem 1.2 (Weak Maximum Principle)}
\begin{theorem}[Weak Maximum Principle]\label{thm:weak_maximum_principle}
Let \( u \in \X(\Omega) \) be a supersolution of \eqref{eq:critical_equation}. Then \( u \geq 0 \) almost everywhere in \( \Omega \).
\end{theorem}
\begin{proof}
On contrary assume that $u \not \geq 0 $ everywhere in $\om$. Then there exists a set \( E \) of positive measure such that \( u < 0 \) in \( E \). 
Take \( u^- \) as a test function and Hardy inequality, we have
\begin{equation}\label{eq:nonnegative_ip1}
\begin{aligned}
0  &\leq \int_{\Omega} \nabla u  \nabla u^- + \langle u, u^- \rangle_{s} - \mu \int_{\Omega} \frac{u u^-}{|x|^2}  = -\int_{\Omega} |\nabla u^-|^2 + \mu \int_{\Omega} \frac{(u^-)^2}{|x|^2} + \langle u, u^-\rangle_s \\
& \leq - \left(1 - \frac{\mu}{\bar{\mu}} \right) \int_{\Omega} |\nabla u^-|^2 +  \langle u, u^-\rangle_s \leq \langle u , u^- \rangle_s.
\end{aligned}
\end{equation}
% Noting that
% \begin{equation*}
% (u^+(x) - u^+(y))(u^-(x) - u^-(y)) = -u^+(x)u^-(y) - u^-(y)u^+(x) \leq 0,
% \end{equation*}
Moreover, using inner product properties, we obtain
\begin{equation*}
\langle u^-, u^- \rangle_{s} = \iint_{\mathbb{R}^{2N}} \frac{|u^-(x) - u^-(y)|^2}{|x - y|^{N+2s}} dy\, dx
\geq \iint_{E \times (\mathbb{R}^N \setminus \Omega)} \frac{|u^-(x)|^2}{|x - y|^{N+2s}} dy\, dx > 0.
\end{equation*}
Thus, we deduce
\begin{equation*}
\langle u, u^- \rangle_{s} = \langle u^+, u^- \rangle_{s} - \langle u^-, u^- \rangle_{s}
< \langle u^+, u^- \rangle_{s} \leq 0.
\end{equation*}
Combining this with \eqref{eq:nonnegative_ip1}, we arrive at a contradiction. 
\end{proof}

% To link the solutions of equation \eqref{eq:critical_equation} with the solutions of problem \eqref{eq:main_problem_ep}, we use the following Weak Maximum Principle.
% %Theorem 1.2 (Weak Maximum Principle)}
% \begin{theorem}[Weak Maximum Principle]
% Let \( u \in \X(\Omega) \) be a supersolution of \eqref{eq:critical_equation}. Then \( u \geq 0 \) almost everywhere in \( \Omega \).
% \end{theorem}
% \begin{proof}
% On contrary assume that $u \not \geq 0 $ everywhere in $\om$. Then there exists a set \( E \) of positive measure such that \( u < 0 \) in \( E \). 
% Take \( u^- \) as a test function, we have
% \begin{align}\label{eq:nonnegative_ip1}
% 0  &\leq \int_{\Omega} \nabla u  \nabla u^- + \langle u, u^- \rangle_{s} - \mu \int_{\Omega} \frac{u u^-}{|x|^2}  = -\int_{\Omega} |\nabla u^-|^2 + \mu \int_{\Omega} \frac{(u^-)^2}{|x|^2} + \langle u, u^-\rangle_s \leq \langle u , u^- \rangle_s,
% \end{align}
% where in the last inequality we have used the Hardy's Inequality as
% \begin{equation*}
% \int_{\Omega} |\nabla u^-|^2 - \mu \int_{\Omega} \frac{|u^-|^2}{|x|^2} \geq \left(1 - \frac{\mu}{\bar{\mu}} \right) \int_{\Omega} |\nabla u^-|^2 \geq 0.
% \end{equation*}
% Noting that
% \begin{equation*}
% (u^+(x) - u^+(y))(u^-(x) - u^-(y)) = -u^+(x)u^-(y) - u^-(y)u^+(x) \leq 0,
% \end{equation*}
% Moreover, using inner product properties, we obtain
% \begin{equation*}
% \langle u^-, u^- \rangle_{s} = \iint_{\mathbb{R}^{2N}} \frac{|u^-(x) - u^-(y)|^2}{|x - y|^{N+2s}} dy\, dx
% \geq \iint_{E \times (\mathbb{R}^N \setminus \Omega)} \frac{|u^-(x)|^2}{|x - y|^{N+2s}} dy\, dx > 0.
% \end{equation*}
% Thus, we deduce
% \begin{equation*}
% \langle u, u^- \rangle_{s} = \langle u^+, u^- \rangle_{s} - \langle u^-, u^- \rangle_{s}
% < \langle u^+, u^- \rangle_{s} \leq 0.
% \end{equation*}
% Combining this with \eqref{eq:nonnegative_ip1}, we arrive at a contradiction. 
% \end{proof}

Applying the Weak Maximum Principle in Theorem \ref{thm:weak_maximum_principle}, we obtain \( u_0 \geq 0 \) almost everywhere in \( \mathbb{R}^N \), with $u_0$ being solution of equation \eqref{eq:critical_equation} holds. This implies $(u_0)_+ \equiv u_0$.
Thus, \( u_0 \) is a solution to problem \eqref{eq:main_problem}.
Hence, it is now sufficient to establish the existence of a nonzero critical point of the functional \( \Jplus \). To this end, we first prove that the functional $\J^+$ satisfies the mountain pass geometry.

%{Mountain Pass Geometry and Palais-Smale Condition}

\begin{lemma}\label{lem:mountain_pass_superlinear}
The functional $\J^+$ satisfies the following
\begin{enumerate}
    \item[$(i)$] There exists $\al, \rho > 0$ such that for any $u \in \X(\om)$ with $\lv u \rv_{\X} = \rho$ we have $\J^+(u) \geq \al$;
    \item[$(ii)$] for any $u \in \X(\om)$ we have $\J^+(tu) \to - \infty$ as $ t \to \infty$.
\end{enumerate}
\end{lemma}
\begin{proof}
Sobolev inequality and Hardy inequality yield
\begin{equation*}
    \J^+(u) \geq \frac{1}{2} \left( 1 - \frac{\mu}{\bar \mu} \right) \lv u \rv_{\X}^2 - C \lv u \rv_{\X}^{2^*} - C \la \lv u \rv_{\X}^p.
\end{equation*}
Then $(i)$ follows directly. On the other hand
\begin{equation*}
    \J^+(tu) \leq \frac{t^2}{2} \lv u \rv_{\X}^2 - \frac{t^{2^*}}{2^*} \int_{\om} (u^+)^{2^*} - \la \frac{ t^p}{p} \int_{\om} (u^+)^p \to - \infty \quad \text{as } t \to \infty.
\end{equation*}
This concludes the proof of $(ii)$.
\end{proof}

\begin{lemma}\label{lem:ps_holds_superlinear}
The functional \( \Jplus \) satisfies the (PS)\(_c\) condition for every \( c < \frac{1}{N} S_{\mu}^{N/2} \). 
\end{lemma}
\begin{proof}
We begin by evaluating
\begin{equation*}
    \Jplus(u_n) - \frac{1}{2^*} \langle (\Jplus)^{\prime}(u_n), u_n \rangle = \frac{1}{N} \left( \int_{\Omega} |\nabla u_n|^2 \, dx + [u_n]^2_{s}  - \mu \int_{\Omega} \frac{|u_n|^2}{|x|^2} \, dx \right) 
    - \lambda \left( \frac{1}{p} - \frac{1}{2^*} \right) \|u_n\|_{L^p}^p.
\end{equation*}
Since \( \{u_n\} \) is a (PS)\(_c\) sequence, we obtain
\begin{equation*}
    C \left(1 + \|u_n\|_{{\X}}\right) \geq \frac{1}{N} \left(1 - \frac{\mu}{\bar{\mu}} \right) \|u_n\|_{\X}^2 
    - \lambda \left( \frac{1}{p} - \frac{1}{2^*} \right) \|u_n\|_{\X}^p,
\end{equation*}
which implies that the sequence \( \{u_n\} \) is bounded. Thus, there exists \( u \in \X(\Omega) \) such that $u_n \rightharpoonup u$ weakly in $\X(\om)$, $ u_n \rightarrow u$ strongly in $L^r(\om)$ for $r \in [1, 2^*)$, and $u_n(x) \rightarrow u(x)$ pointwise a.e. in $\om$.
\noi Now,
\begin{align*}
    &o_n(1)  = \langle (\Jplus)^{\prime}(u_n), u_n - u \rangle \\
    & = \left( \|\nabla u_n\|^2_2 - \|\nabla u\|^2_2 \right) +  \left( [u_n]^2_{s} - [u]^2_{s} \right)
    - \mu \int_{\Omega} \frac{|u_n|^2 - |u|^2}{|x|^2} \, dx - (\|u_n^+\|_{L^{2^*}}^{2^*} - \|u^+\|_{L^{2^*}}^{2^*}) + o_n(1) \\
    & = \|\nabla u_n - \nabla u\|^2_2 +  [u_n - u]^2_{s} 
    - \mu \int_{\Omega} \frac{|u_n - u|^2}{|x|^2} \, dx - \|u_n^+ - u^+\|_{L^{2^*}}^{2^*} + o_n(1),
\end{align*}
where the final equality follows from the Brezis–Lieb Lemma. Thus, we conclude that
\begin{equation*}
    \lim_{n \to \infty} \left( \|u_n - u\|^2_{\X} - \mu \int_{\Omega} \frac{|u_n - u|^2}{|x|^2} \, dx \right)
    = \lim_{n \to \infty} \|u_n^+ - u^+\|_{L^{2^*}}^{2^*} = \ell  (\text{say}).
\end{equation*}
By \eqref{eq:spectral_ratio_frac}, we also have
\begin{equation*}
   S_{\mu} \|u_n^+ - u^+\|_{L^{2^*}}^2 \leq S_{\mu} \|u_n - u\|_{L^{2^*}}^2 \leq \|u_n - u\|^2_{\X} - \mu \int_{\Omega} \frac{|u_n - u|^2}{|x|^2} \, dx.
\end{equation*}
This implies $ S_\mu \, \ell^{\frac{2}{2^*}} \leq \ell$.  If \( \ell > 0 \), then $\ell \geq S_\mu^{\frac{N}{2}}$. Next, consider
\begin{align*}
    \Jplus(u_n) - \frac{1}{2} \langle (\Jplus)^{\prime}(u_n), u_n \rangle 
    &= \left( \frac{1}{2} - \frac{1}{2^*} \right) \|u_n^+\|_{L^{2^*}}^{2^*} + \lambda \left( \frac{1}{2} - \frac{1}{p} \right)\int_{\Omega} |u_n^+|^p  dx \geq \frac{1}{N}\|u_n^+\|_{L^{2^*}}^{2^*}.
\end{align*}
This gives $c \geq \frac{1}{N}S_\mu^{\frac{N}{2}}$.
\end{proof}
% \begin{proof}
% Let \( (u_n) \) be a (PS)\(_c\) sequence for \( \Jplus \) in \( \X(\Omega) \). Then
% \begin{equation}
% \Jplus(u_n) \to c \quad \text{as } n \to \infty, \quad
% (\Jplus)^{\prime}(u_n) \to 0 \quad \text{as } n \to \infty.
% \label{eq:PS_derivative}
% \end{equation}
% Note that 
% \begin{equation}\label{eq:ps_lp_term_bound}
%     C(1+ \lv u_n\rv_{\X} ) \geq \J(u_n) - \frac{1}{p} \langle \J^{\prime}(u_n),u_n\rangle  \geq \left( \frac{1}{2} - \frac{1}{p}\right) \int_{\om} |u_n|^p.
% \end{equation}
% Consider
% \begin{align*}
%     C(1+ \lv u_n\rv_{\X} ) &\geq \J(u_n) - \frac{1}{2^*} \langle \J^{\prime}(u_n),u_n\rangle  \\
%      &\geq \frac{1}{N} \left( 1 - \frac{\mu}{\bar{\mu}} \right) \lv u_n \rv_{\X}^2 - \la \left( \frac{1}{p} - \frac{1}{2^*}\right) \int_{\om} |u_n|^p.
% \end{align*}
% From above and \eqref{eq:ps_lp_term_bound} we deduce that \( (u_n) \) is a bounded sequence. Thus there exists \( u_0 \in \X(\Omega) \) such that
% \begin{equation*}
% u_n \rightharpoonup u_0 \quad \text{in } \X(\Omega).
% \end{equation*}
% % This implies
% % \begin{equation*}
% % \int_{\Omega} \nabla u_n  \nabla \varphi + \langle u_n, \varphi \rangle_{s} 
% % \to \int_{\Omega} \nabla u_0   \nabla \varphi + \langle u_0, \varphi \rangle_{s} \quad \mbox{ as }  n \to \infty.
% % \end{equation*}
% % Moreover, as \( \frac{u_n(x)}{|x|} \to \frac{u_0(x)}{|x|} \) almost everywhere in \( \Omega \), and \( \left\{ \frac{u_n}{|x|} \right\} \) is bounded in \( W^{1,p}_0(\Omega) \), we deduce:
% % \begin{equation*}
% % \int_{\Omega} \frac{u_n \varphi}{|x|^2} \to \int_{\Omega} \frac{u_0 \varphi}{|x|^2} \quad \text{as } n \to \infty.
% % \end{equation*}
% % Since $u_n^+(x) = \frac{u_n(x) + |u_n(x)|}{2}$, we obtain $u_n^+(x) \to u_0^+(x)$ almost everywhere in $\om$ as $n \to \infty$. Again, weak convergence in \eqref{eq:weak_convergence} implies:
% % \begin{equation*}
% % \int_{\Omega} (u_n^+)^{2^*-1} \varphi \to \int_{\Omega} (u_0^+)^{2^*-1} \varphi
% % \end{equation*}
% % and 
% % \begin{equation*}
% % \int_{\Omega} (u_n^+)^{p-1} \varphi \to \int_{\Omega} (u_0^+)^{p-1} \varphi
% % \end{equation*}
% % as \( n \to \infty \). Thus, \( u_0 \) is a weak solution of problem \eqref{eq:critical_equation}. Hence
% % \begin{equation*}
% % \langle \Jplus(u_0), \varphi \rangle = 0.
% % \end{equation*}
% In addition, from the fact $u_n^+(x) = \frac{u_n(x) + |u_n(x)|}{2}$, Brezis-Lieb and \eqref{eq:PS_derivative}, we deduce that 
% \begin{equation*}
%     \begin{aligned}
%     &o_n(1) = \langle (\Jplus)^{\prime}(u_n),u_n\rangle \\
%     & = \lv u_n - u_0 \rv_{\X}^2 - \mu \int_{\om} \frac{|u_n -u_0|^2}{|x|^2} - \int_{\om} |u_n^+ - u_0^+|^{2^*} + \langle (\Jplus)^{\prime}(u_0),u_0\rangle + o_n(1).
%     \end{aligned}
% \end{equation*}
% That is equivalent to
% \begin{equation}\label{eq:terms_equivalence}
%      o_n(1) + \int_{\om} |u_n^+ - u_0^+|^{2^*}  = \lv u_n - u_0 \rv_{\X}^2 - \mu \int_{\om} \frac{|u_n -u_0|^2}{|x|^2} dx
% \end{equation}
% Applying the above estimate, Brezis-Lieb Lemma and weak convergence, we obtain
% \begin{equation}
% \Jplus(u_n) = o_n(1) + \Jplus(u_0) + \frac{1}{N} \lv u_n - u_0 \rv_{\X}^2 
% - \frac{\mu}{N} \int_{\Omega} \frac{|u_n - u_0|^2}{|x|^2}dx.
% \label{eq:Brezis_Lieb}
% \end{equation}
% Now, upto a subsequence, we have
% \begin{equation*}
% \lv u_n - u_0 \rv_{\X}^2 - \mu \int_{\Omega} \frac{|u_n - u_0|^2}{|x|^2} \to \ell, \quad \text{as } n \to \infty.
% \end{equation*}
% Equation \eqref{eq:terms_equivalence} implies that $\ell \leq \tilde{\ell},$ where $\int_{\Omega} |u_n - u_0|^{2^*} dx\to \tilde{\ell}$. 
% Using the definition of $S_{\mu}$ in \eqref{eq:spectral_ratio_frac}, we obtain
% \begin{equation*}
% \ell \geq S_{\mu,s}(\om) \tilde{\ell}^{\frac{2}{2^*}} = S_{\mu} \tilde{\ell}^{\frac{2}{2^*}} \geq S_{\mu} \ell^{\frac{2}{2^*}} 
% \end{equation*}
% This simplifies to either $\ell = 0$ or $ \ell \geq S_{\mu}^{N/2}$. Then $\langle (\Jplus)^{\prime}(u_0), u_0 \rangle = 0$ gives
% \begin{equation*}
% \Jplus(u_0) = \lambda \left( \frac{1}{2}-\frac{1}{p} \right) 
% \int_{\Omega} (u_0^+)^{p} + \left( \frac{1}{2}-\frac{1}{2^*} \right) \int_{\Omega} (u_0^+)^{2^*} \geq 0.
% \end{equation*}
% Thus, equation \eqref{eq:Brezis_Lieb} implies that 
% \begin{equation*}
%     c = \Jplus(u_0) + \frac{\ell}{N} \geq \frac{S_{\mu}^{N/2}}{N}
% \end{equation*}
% This forces $\ell = 0$ in view of our assumption on $c$. Hence, we have our result.
% \end{proof}

% Now, we estimate the following term with \( r = p+1 \) using the inequality:
% \begin{equation}
% (b - c)^p \geq b^p - p b^{p-1} c, \quad \text{for } b \geq c > 0 \text{ and } p \geq 2.
% \label{eq:ineq}
% \end{equation}
% Thus, we have
% \begin{equation*}
% \| u_{\varepsilon} \|_{L^r}^r \geq \int_{B_{1/m}} |u_{\varepsilon}|^r \, dx - r \int_{B_{1/m}} |u_{\varepsilon}|^{r-1} \frac{B_0 \varepsilon^{\sqrt{\bar{\mu}}}} {\left[ \varepsilon^2 \left(\frac{1}{m} \right)^{\frac{\ga^{\prime}}{ \sqrt{\bar{\mu}}}} + \left( \frac{1}{m} \right)^{\frac{\gamma}{ \sqrt{\bar{\mu}}}} \right]} \, dx.
% \end{equation*}
% Splitting the integral into different terms, we define
% \begin{equation*}
% I_1 = \int_{\mathbb{R}^N} |u_{\varepsilon}|^r \, dx, \quad
% I_2 = \int_{\mathbb{R}^N \setminus B_{1/m}} |u_{\varepsilon}|^r \, dx, \quad
% I_3 = r \int_{B_{1/m}} |u_{\varepsilon}|^{r-1} \frac{B_0 \varepsilon^{\sqrt{\bar{\mu}}}} {\left[ \varepsilon^2 \left(\frac{1}{m} \right)^{\frac{\ga^{\prime}}{ \sqrt{\bar{\mu}}}} + \left( \frac{1}{m} \right)^{\frac{\gamma}{ \sqrt{\bar{\mu}}}} \right]} \, dx.
% \end{equation*}
% Taking the transformation \( y = \frac{x}{\varepsilon^{ \sqrt{\bar{\mu}} /\sqrt{\bar{\mu} - \mu}}} \), we obtain:
% \begin{equation*}
% I_1 = \int_{\mathbb{R}^N} |u_{\varepsilon}|^r \, dx = \int_{\mathbb{R}^N} \frac{|B_0|^r (\varepsilon ^{\sqrt{\bar{\mu}}})^r}{\left(\varepsilon^2 |x|^{\gamma^{\prime} /\sqrt{\bar{\mu}}} + |x|^{\gamma /\sqrt{\bar{\mu}}} \right)^{r\sqrt{\bar{\mu}}}} \, dx = \varepsilon^{\frac{\sqrt{\bar{\mu}} (N - r \sqrt{\bar{\mu}})}{\sqrt{\bar{\mu} - \mu}}}  \| u_{1} \|_{L^r(\mathbb{R}^N)}^r.
% \end{equation*}
% By simplifying, we get
% \begin{equation*}
% I_2 = \int_{\mathbb{R}^N \setminus B_{1/m}} |u_{\varepsilon}|^r \, dx = C \int_{1/m}^{\infty} \frac{t^{N-1} (\varepsilon \sqrt{\bar{\mu}})^r}{\left[ \varepsilon^2 t^{\frac{\gamma^{\prime}}{ \sqrt{\bar{\mu}}}} + t^{\frac{\gamma}{ \sqrt{\bar{\mu}}}} \right]^{ r \sqrt{\bar{\mu}}}} \, dt \leq C \varepsilon^{r \sqrt{\bar{\mu}}} \int_{1/m}^{\infty} t^{N-1} t^{-r\gamma} \, dt = C \varepsilon^{r \sqrt{\bar{\mu}}} m^{-(N - \gamma r)}.
% \end{equation*}
% The last integral converges when \( N - r\gamma < 0 \), i.e., $\mu < \bar{\mu} - \left( \frac{N - r \sqrt{\bar{\mu}}}{r} \right)^2$.
% %#### **Evaluation of \( I_3 \):**
% Using similar bounding techniques
% \begin{equation*}
% I_3 \leq C m^{\gamma}\int_{0}^{1/m}  \frac{ t^{N-1} \varepsilon^{r \sqrt{\bar{\mu}}}} {\left[ \varepsilon^2 t^{\frac{\ga^{\prime}}{ \sqrt{\bar{\mu}}}} + t^{\frac{\gamma}{ \sqrt{\bar{\mu}}}} \right]^{(r-1)\sqrt{\bar{\mu}}}} \, dt \leq C \varepsilon^{r \sqrt{\bar{\mu}}} m^{-(N - r\gamma)},
% \end{equation*}
% where the last integral is finite as \( N - n(\gamma - 1) > 0 \) for \( N \geq 4 \). Combining all estimates, we get:
% \begin{equation}\label{eq:perturbed_term_bound}
% \| u_{\varepsilon} \|_{L^r}^r \geq \varepsilon^{\frac{\sqrt{\bar{\mu}} {(N - r\sqrt{\bar{\mu}})}}{{\sqrt{\bar{\mu} - \mu}}}} \| u_{1} \|_{L^r}^r - C \varepsilon^{r\sqrt{\bar{\mu}}} m^{-(N - r\gamma)}.
%\end{equation}


%%%%%%%%%%%%%%%%%%%%%%%%%%%%%%%%%%%%%%%%%%%%%%%%%%%%%%%%%%
%%%%% Blow up analysis
%%%%%%%%%%%%%%%%%%%%%%%%%%%%%%%%%%%%%%%%%%%%%%%%%%%%%%%%%%

Let us define the following parameters
\begin{equation*}
    \bar{\mu} = \left( \frac{N-2}{2} \right)^2, \quad \gamma = \sqrt{\bar{\mu}} + \sqrt{\bar{\mu} - \mu}, \quad \gamma^{\prime} = \sqrt{\bar{\mu}} - \sqrt{\bar{\mu} - \mu}.
\end{equation*}
For $\varepsilon > 0$, we define the function $U_{\varepsilon,\al}(x)$ as follows
\begin{equation*}
    U_{\varepsilon,\al}(x) =  \frac{B_0 \varepsilon^{ \alpha \sqrt{\bar{\mu}}}}{\left( \varepsilon^{2\alpha} |x|^{\gamma^{\prime}/\sqrt{\bar{\mu}}} + |x|^{\gamma/\sqrt{\bar{\mu}}} \right)^{\sqrt{\bar{\mu}}}},
\end{equation*}
where the constant $B_0 = \left( \frac{4N(\bar{\mu} - \mu)}{N-2} \right)^{(N-2)/4}$ and the parameter $\al > 0$ is chosen later.
The function $U_{\varepsilon,\al}\in D^{1,2}(\mathbb{R}^N)$ satisfies the following partial differential equation
\begin{equation*}
    -\Delta U - \frac{\mu}{|x|^2} U = |U|^{2^*-2} U.
\end{equation*}
Furthermore, $U_{\varepsilon,\al}$ achieves the Sobolev constant $S_{\mu}$, i.e.,
\begin{equation*}
    \int_{\mathbb{R}^N} \left( |\nabla U_{\varepsilon,\al}|^2 - \frac{\mu}{|x|^2} U_{\varepsilon,\al}^2 \right) dx = \int_{\mathbb{R}^N} |U_{\varepsilon,\al}|^{2^*} dx = S_{\mu}^{N/2}.
\end{equation*}
% Next, we define the function $u_{\varepsilon}^{m}(x)$ as follows:
% \begin{equation*}
%     u_{\varepsilon}^{m}(x) =
%     \begin{cases} 
%         U_{\varepsilon}(x) -  \frac{B_0 \varepsilon^{ \sqrt{\bar{\mu}}}}{\left( \varepsilon^2 \left( \frac{1}{m} \right)^{\gamma^{\prime}/\sqrt{\bar{\mu}}} + \left( \frac{1}{m} \right)^{\gamma/\sqrt{\bar{\mu}}} \right)^{\sqrt{\bar{\mu}}}}, & \text{if } x \in B_{1/m} \setminus \{ 0\}, \\ 
%         0, & \text{if } x \in \Omega \setminus B_{1/m}.
%     \end{cases}
% \end{equation*}
Next define $u_{\varepsilon,\al}(x) = \phi (x) U_{\varepsilon,\al}(x)$ where $\phi \in C_{c}^2(\om)$ is a cut-off function satisfying $0 \leq \phi(x) \leq 1$, $\phi \equiv 1$ in $B_{\delta} $, $Supp(\phi) \subset B_{2\delta} \subset \om$, and $|\grad \phi |\leq C$. Then from \cite{chen2003estimates}, we have the following estimates 
\begin{align}
\int_{\Omega} |\nabla u_{\varepsilon,\al}|^2 - \mu \int_{\Omega} \frac{u_{\varepsilon,\al}^2}{|x|^2} dx&= S_{\mu}^{N/2} + O(\varepsilon^{\al (N-2)}),
\label{eq:norm_upper_bound}\\
\| u_{\varepsilon,\al} \|^{2^*}_{L^{2^*}} &\geq S_{\mu}^{N/2} - O( \varepsilon^{\al N} ),
\label{eq:norm_lower_bound}\\
\int_{\om}|u_{\varepsilon,\al}|^{p}dx &= O(\varepsilon^{ \al (N - \sqrt{\bar\mu}p)\frac{\sqrt{\bar\mu}}{\sqrt{\bar\mu - \mu}}}),\label{eq:lower_bound_norm_r}
\end{align}
provided that $\mu < \bar{\mu} -1$ and $p > \frac{N}{\gamma}$. \\
%%%%%{Gagliardo norm estimate}

% Now we estimate the perturbation term. For \( |x| \leq \varepsilon^{\beta} \), with \( \beta = \frac{\sqrt{\bar{\mu}}}{\sqrt{\bar{\mu} - \mu}} \), we note
% \begin{equation}
%     \varepsilon^{2}  |x|^{\gamma^{\prime}/\sqrt{\bar{\mu}}} + |x|^{\gamma/\sqrt{\bar{\mu}}}
%     \leq 2 \varepsilon^{\gamma / \sqrt{\bar{\mu}-\mu}}.
% \label{eq:exponential_decay}
% \end{equation}
% Let \( q = r^{1 / \gamma^{\prime}} > 1 \), then for sufficiently small \( \varepsilon \), we have
% \begin{equation}
%     B_{\frac{\varepsilon^{\beta}}{q}} \subset B_{\frac{1}{qm}}, \quad \text{and} \quad U_\varepsilon(x) \geq U_{\varepsilon}\left(\frac{1}{qm}\right)> r \, U_\varepsilon\left(\frac{1}{m}\right).
% \label{eq:ball_inclusion}
% \end{equation}
% Consider
% \begin{equation}
% \begin{aligned}
%     \| u_\varepsilon^m \|_{L^r(B_{1/m})}^\eta 
%     &= \int_{B_{1/m}} |u_\varepsilon^m|^r \, dx \\
%     &\geq \int_{B_{\varepsilon^{\beta}/q}} |u_\varepsilon^m|^r \, dx \\
%     &= \int_{B_{\varepsilon^{\beta}/q}} 
%         \left| \frac{1}{r} \left( (r-1) U_\varepsilon(x) + U_\varepsilon(x) - r U_\varepsilon\left(\frac{1}{m}\right) \right) \right|^r dx \\
%     &\geq \int_{B_{\varepsilon^{\beta}/2}} \left( \frac{r - 1}{r} \right)^r |U_\varepsilon(x)|^r \, dx\\
%   &\geq \left( \frac{r - 1}{r} \right)^r 
%      \frac{\left(B_0 \, \varepsilon^{\sqrt{\bar{\mu}}}\right)^r}{(2 \varepsilon^{\gamma/\sqrt{\bar{\mu}-\mu} })^{r\sqrt{\bar{\mu}}}} 
%     \left| B_{\varepsilon^{\beta}/q} \right| \\
%     &= C \, \varepsilon^{r\sqrt{\bar{\mu}} - r \frac{\sqrt{\bar{\mu}}\gamma}{\sqrt{\bar{\mu} - \mu}} + \beta N} \\
%     &= C \, \varepsilon^{\frac{\sqrt{\bar{\mu}}(N - r \sqrt{\bar{\mu}})}{\sqrt{\bar{\mu} - \mu}}}.
% \end{aligned}
% \label{eq:lower_bound_norm_r}
% \end{equation}
We now estimate the Gagliardo norm of $u_{\varepsilon,\al}(x)$. Using the estimate $(6.12)$ along with Exercise $1.26$ of \cite{giovanni_fractional_book} and \eqref{eq:norm_upper_bound} we deduce
\begin{equation}\label{eq:fractional_norm_bound}
\begin{aligned}
\iint_{\mathbb{R}^{2N}} \frac{| u_{\varepsilon,\al}(x) - u_{\varepsilon,\al}(y)|^2}{|x - y|^{N+2s}} dx \, dy & \leq C \left( \int_{\rnn} |u_{\varepsilon,\al}|^2 \, dx \right)^{1-s}\left( \int_{\rnn} |\grad u_{\varepsilon,\al}|^2\,dx\right)^s\\
& \leq C \varepsilon^{2\al(1-s) \frac{\sqrt{\bar{\mu}}}{\sqrt{\bar \mu - \mu}}} \left( \int_{\rnn} |\grad u_{\varepsilon,\al}|^2\,dx - \mu \int_{\Omega} \frac{u_{\varepsilon,\al}^2}{|x|^2} + \mu \int_{\rnn} \frac{U_{\varepsilon,\al}^2}{|x|^2} \right)^s\\
& \leq C \varepsilon^{2\al(1-s) \frac{\sqrt{\bar{\mu}}}{\sqrt{\bar \mu - \mu}}} \left(  S_{\mu}^{N/2} + O(\varepsilon^{\al (N-2)}) + \mu \int_{\rnn} \frac{U_{1,\al}^2}{|x|^2} \right)^s\\
& \leq C \varepsilon^{2\al(1-s) \frac{\sqrt{\bar{\mu}}}{\sqrt{\bar \mu - \mu}}}
\end{aligned}
\end{equation}
% Applying substitution $x = \varepsilon^{\frac{\sqrt{\bar{\mu}}}{\sqrt{\bar{\mu} - \mu}}}  z, \quad y =\varepsilon^{\frac{\sqrt{\bar{\mu}}}{\sqrt{\bar{\mu} - \mu}}}  w$, we obtain
% \begin{equation}
% [u_{\varepsilon}^{m}]_{s}^{2} \leq \varepsilon ^{ \frac{\sqrt{\bar{\mu}} (N+2 - 4s)}{\sqrt{\bar{\mu} - \mu}}} \textcolor{red}{\varepsilon ^{ \frac{\sqrt{\bar{\mu}} (2-2s)}{\sqrt{\bar{\mu} - \mu}}}} [U_1]_{s}^{2}.
% \label{eq:fractional_norm_bound}
% \end{equation}



To conclude, we use the mountain pass theorem with the help of functions $u_{\varepsilon,\al}$ defined above. Fixing $\al =1$ and defining $u_{\varepsilon} = u_{\varepsilon, \alpha}$, we combine the estimates in \eqref{eq:norm_upper_bound}, \eqref{eq:norm_lower_bound}, \eqref{eq:lower_bound_norm_r} with $r = p$ and \eqref{eq:fractional_norm_bound} to obtain
\begin{equation*}
\begin{aligned}
\Jplus(t u_{\varepsilon}) &  \leq \frac{t^2}{2} \left[ S_{\mu}^{N/2} + C \varepsilon^{(N-2)}  \right] - \frac{t^{2^*}}{2^*}\left[ S_{\mu}^{N/2} - C \varepsilon^{ N}  \right] -\frac{\la t^{p}}{p} C \varepsilon^{\frac{ \sqrt{\bar{\mu}} {(N - p\sqrt{\bar{\mu}})}}{{\sqrt{\bar{\mu} - \mu}}}} + \frac{t^2}{2} C \varepsilon ^{ \frac{2  (1-s)\sqrt{\bar{\mu}}}{\sqrt{\bar{\mu} - \mu}}}\\
& \leq \frac{t^2}{2} \left[ S_{\mu}^{N/2} + C \varepsilon^{\beta_{\mu,N,s}}  \right] - \frac{t^{2^*}}{2^*}\left[ S_{\mu}^{N/2} - C \varepsilon^{ N}  \right] -\frac{\la t^{p}}{p} C \varepsilon^{\frac{ \sqrt{\bar{\mu}} {(N - p\sqrt{\bar{\mu}})}}{{\sqrt{\bar{\mu} - \mu}}}} := g(t)
\end{aligned}
\end{equation*}
where $\beta_{\mu,N,s}$ is defined in \eqref{eq:superlinear_parameter}.
Observing the behavior of \( g(t) \)
\begin{equation*}
g(0) = 0, \quad \text{and} \quad g(t) \to -\infty \text{ as } t \to \infty.
\end{equation*}
Thus, there exists \( t_{\varepsilon, \lambda} > 0 \) such that
\begin{equation*}
\sup_{t > 0} g(t) = g(t_{\varepsilon, \lambda}).
\end{equation*}
If \( t_{\varepsilon, \lambda} = 0 \), then there is nothing to prove. Otherwise, we take \( t_{\varepsilon, \lambda} > 0 \).
Differentiating \( g(t) \) and equating it to zero
\begin{equation*}
0 = g'(t_{\varepsilon, \lambda}) = t_{\varepsilon, \lambda} \left[ S_{\mu}^{N/2} + C \varepsilon^{\beta_{\mu,N,s}} \right] - t_{\varepsilon, \lambda}^{2^*-1} \left[ S_{\mu}^{N/2} - C \varepsilon^{ N} \right] - \lambda t_{\varepsilon, \lambda}^{p-1} C \varepsilon^{\frac{ \sqrt{\bar{\mu}} {(N - p\sqrt{\bar{\mu}})}}{{\sqrt{\bar{\mu} - \mu}}}}.
\end{equation*}
Rearranging terms
\begin{equation}\label{eq:g_derivative_0}
S_{\mu}^{N/2} + C \varepsilon^{ \beta_{\mu,N,s}} = t_{\varepsilon, \lambda}^{2^*-2} \left[ S_{\mu}^{N/2} - C \varepsilon^{ N} \right] + \lambda t_{\varepsilon, \lambda}^{p-2} C \varepsilon^{\frac{ \sqrt{\bar{\mu}} {(N - p\sqrt{\bar{\mu}})}}{{\sqrt{\bar{\mu} - \mu}}}}.
\end{equation}
This gives $\displaystyle\liminf_{\varepsilon \to 0} t_{\varepsilon,\la} > 0$. Indeed if $\displaystyle\liminf_{\varepsilon \to 0} t_{\varepsilon,\la} =0$, we arrive at $S_{\mu}^{N/2} = 0$ which is absurd. In addition, from \eqref{eq:g_derivative_0}, we derive an upper bound for \( t_{\varepsilon, \lambda} \) as
\begin{equation*}
t_{\varepsilon, \lambda} < \left( \frac{S_{\mu}^{N/2} + C \varepsilon^{\beta_{\mu,N,s}}}{S_{\mu}^{N/2} - C \varepsilon^{ N}} \right)^{\frac{1}{2^*-2}} = 1 + C \varepsilon^{\beta_{\mu,N,s}}= h(\varepsilon).
\end{equation*}
As \( \varepsilon \to 0 \), it follows that \( h(\varepsilon) \to 1 \) and $t_{\varepsilon, \lambda} \geq \theta_{\la} > 0$ for some $\theta_\la$. Furthermore, consider the function
\begin{equation*}
h_1(t) = \frac{t^2}{2} \left[ S_{\mu}^{N/2} + C \varepsilon^{\beta_{\mu,N,s}} \right] - \frac{t^{2^*}}{2^*} \left[ S_{\mu}^{N/2} - C \varepsilon^{ N} \right].
\end{equation*}
Then $h_1(t)$ is increasing on \( (0, h(\varepsilon)] \). Thus, we obtain
\begin{align*}
\sup_{t>0} g(t) & = g(t_{\varepsilon,\la})  \\
& \leq   S_{\mu}^{N/2} \left[ \frac{h(\varepsilon)^2}{2} - \frac{h(\varepsilon)^{2^*}}{2^*} \right] + \frac{C}{2} h(\varepsilon)^2 \varepsilon^{\beta_{\mu,N,s}} + \frac{h(\varepsilon)^{2^*} C \varepsilon^{ N}}{2^*}  -\frac{\la t_{\varepsilon,\la}^{p}}{p} C \varepsilon^{\frac{ \sqrt{\bar{\mu}} {(N - p\sqrt{\bar{\mu}})}}{{\sqrt{\bar{\mu} - \mu}}}}\\
& \leq   S_{\mu}^{N/2} \left(1 + \varepsilon^{\beta_{\mu,N,s}} \right)^{\frac{2}{2^*-2}} \left[ \frac{1}{2} - \frac{1 + \varepsilon^{\beta_{\mu,N,s}}}{2^*} \right] + C \varepsilon^{\beta_{\mu,N,s}}  -\la C \theta_{\la}^p \varepsilon^{\frac{ \sqrt{\bar{\mu}} {(N - p\sqrt{\bar{\mu}})}}{{\sqrt{\bar{\mu} - \mu}}}}\\
& = \frac{1}{N} S_{\mu}^{N/2} +  C \varepsilon^{\beta_{\mu,N,s}}  -\la \theta_{\la}^p C \varepsilon^{\frac{ \sqrt{\bar{\mu}} {(N - p\sqrt{\bar{\mu}})}}{{\sqrt{\bar{\mu} - \mu}}}}.
\end{align*}
\textbf{Case 1:} For the case $\beta_{\mu,N,s} > \frac{ \sqrt{\bar{\mu}} {(N - p\sqrt{\bar{\mu}})}}{{\sqrt{\bar{\mu} - \mu}}}$, it follows that for sufficiently small \( \varepsilon > 0 \), we obtain
\begin{equation*}
\sup_{t \geq 0} \Jplus(t u_{\varepsilon}) < \frac{1}{N} S_{\mu}^{N/2}.
\end{equation*}
\textbf{Case 2:} When $\beta_{\mu,N,s} \leq \frac{ \sqrt{\bar{\mu}} {(N - p\sqrt{\bar{\mu}})}}{{\sqrt{\bar{\mu} - \mu}}}$, we claim that $\displaystyle \lim_{\lambda \to \infty} t_{\varepsilon, \lambda} = 0$. On contrary
\begin{equation*}
l = \limsup_{\lambda \to \infty} t_{\varepsilon, \lambda} > 0.
\end{equation*}
Choose a sequence \( \lambda_k \to \infty \) as \( k \to \infty \) such that $t_{\varepsilon, \lambda_k} \to l$. Since \( \lambda_k \to \infty \), the right-hand side (RHS) of equation \eqref{eq:g_derivative_0} approaches infinity, which contradicts the left-hand side being finite. This contradiction establishes the claim.


\begin{proof}[Proof of Theorem \ref{thm:superlinear_pb}]
With the above discussion and Lemma \ref{lem:mountain_pass_superlinear}, the hypotheses of the mountain pass theorem hold. This gives the existence of a $(PS)_c$ sequence and Lemma \ref{lem:ps_holds_superlinear} gives that $(PS)_c$ condition is satisfied with $c < \frac{1}{N} S_{\mu}^{N/2}$. This completes the proof. 
\end{proof}






%%%%%%%%%%%%%%%%%%%%%%%%%%%%%%%%%%%%%%%%%%%%%%%%%%%%%%%%%%%%%%%%%%
%%%%%%%%%%%%%%%%    Sublinear Problem            %%%%%%%%%%%%%%%%%
%%%%%%%%%%%%%%%%%%%%%%%%%%%%%%%%%%%%%%%%%%%%%%%%%%%%%%%%%%%%%%%%%%
\section{Sublinear Case (\texorpdfstring{$1<p<2$})}\label{sec:sublinear}
This section is devoted to the sublinear case, that is, \( 1 < p < 2 \) in problem~\eqref{eq:main_problem_ep}. The presence of both convex and concave terms enriches the variational structure of the associated energy functional, allowing us to exploit its topology to establish the existence of two distinct positive solutions. Since we are concerned with positive solutions of \eqref{eq:main_problem_ep}, we follow in the same way as superlinear case and consider the following modified functional
\begin{equation*}
\Jep(u) = \frac{1}{2} \int_{\om} |\grad u|^2 + \frac{\varepsilon}{2} [u]_s^2 - \frac{\mu}{2} \int_{\Omega} \frac{|u|^2}{|x|^2} dx 
- \frac{1}{2^*} \int_{\Omega} (u^+)^{2^*} - \frac{\lambda}{p} \int_{\Omega} (u^+)^{p}.
\end{equation*}
\noi Any critical point of the functional \( \Jep \) is a solution of the following equation

\begin{equation}\tag{$\mathcal{P}_{\mu, \la}^{\varepsilon,+}$}
\begin{cases}
-\Delta u + \varepsilon(-\Delta)^s u - \mu \frac{u}{|x|^2} = (u^+)^{2^*-1} + \lambda (u^+)^{p-1}, & \text{in } \Omega, \\
u = 0, & \text{in } \mathbb{R}^N \setminus \Omega.
\end{cases}
\label{eq:maineq_plus_epsilon}
\end{equation}
% To link the solutions of equation \eqref{eq:maineq_plus_epsilon} with the solutions of problem \eqref{eq:main_problem_ep}, we use the following Weak Maximum Principle.
% %Theorem 1.2 (Weak Maximum Principle)}
% \begin{theorem}[Weak Maximum Principle]\label{thm:weak_maximum_principle}
% Let \( u \in \X(\Omega) \) be a supersolution of \eqref{eq:maineq_plus_epsilon}. Then \( u \geq 0 \) almost everywhere in \( \Omega \).
% \end{theorem}
% \begin{proof}
% On contrary assume that $u \not \geq 0 $ everywhere in $\om$. Then there exists a set \( E \) of positive measure such that \( u < 0 \) in \( E \). 
% Take \( u^- \) as a test function and Hardy inequality, we have
% \begin{equation}\label{eq:nonnegative_ip1}
% \begin{aligned}
% 0  &\leq \int_{\Omega} \nabla u  \nabla u^- + \langle u, u^- \rangle_{s} - \mu \int_{\Omega} \frac{u u^-}{|x|^2}  = -\int_{\Omega} |\nabla u^-|^2 + \mu \int_{\Omega} \frac{(u^-)^2}{|x|^2} + \langle u, u^-\rangle_s \\
% & \leq - \left(1 - \frac{\mu}{\bar{\mu}} \right) \int_{\Omega} |\nabla u^-|^2 +  \langle u, u^-\rangle_s \leq \langle u , u^- \rangle_s.
% \end{aligned}
% \end{equation}
% Noting that
% \begin{equation*}
% (u^+(x) - u^+(y))(u^-(x) - u^-(y)) = -u^+(x)u^-(y) - u^-(y)u^+(x) \leq 0,
% \end{equation*}
% Moreover, using inner product properties, we obtain
% \begin{equation*}
% \langle u^-, u^- \rangle_{s} = \iint_{\mathbb{R}^{2N}} \frac{|u^-(x) - u^-(y)|^2}{|x - y|^{N+2s}} dy\, dx
% \geq \iint_{E \times (\mathbb{R}^N \setminus \Omega)} \frac{|u^-(x)|^2}{|x - y|^{N+2s}} dy\, dx > 0.
% \end{equation*}
% Thus, we deduce
% \begin{equation*}
% \langle u, u^- \rangle_{s} = \langle u^+, u^- \rangle_{s} - \langle u^-, u^- \rangle_{s}
% < \langle u^+, u^- \rangle_{s} \leq 0.
% \end{equation*}
% Combining this with \eqref{eq:nonnegative_ip1}, we arrive at a contradiction. 
% \end{proof}

Applying the Weak Maximum Principle, we obtain \( u_0 \geq 0 \) almost everywhere in \( \mathbb{R}^N \), with $u_0$ being solution of equation \eqref{eq:maineq_plus_epsilon} holds. This implies $(u_0)_+ \equiv u_0$. Thus, \( u_0 \) is a solution to problem \eqref{eq:main_problem_ep}.
Hence, it is now sufficient to establish the existence of a nonzero critical point of the functional \( \Jep \).

We begin by proving the existence of the first positive solution to problem~\eqref{eq:main_problem_ep}. First we prove a lemma that gives the convergence of (PS) sequence. 
\begin{lemma}\label{lem:ps_sublinear}
If \( p \in (1, 2) \), then the (PS)\(_c\) condition for the functional \( \Jep \) holds for all    
\begin{equation*}
    c < \frac{1}{N} S^{N/2} 
    - |\Omega| \left(1 - \frac{ p}{2^*} \right) \left( \frac{pN}{2^*} \right)^{-p/(2^* - p)}
    \left[ \lambda \left( \frac{1}{p} - \frac{1}{2} \right) \right]^{2^*/(2^* - p)}.
\end{equation*}
\end{lemma}
\begin{proof}
% We begin by evaluating
% \begin{equation*}
%     \Jep(u_n) - \frac{1}{2^*} \langle (\Jep)^{\prime}(u_n), u_n \rangle = \frac{1}{N} \left( \int_{\Omega} |\nabla u_n|^2 \, dx + \varepsilon [u_n]^2_{s}  - \mu \int_{\Omega} \frac{|u_n|^2}{|x|^2} \, dx \right) 
%     - \lambda \left( \frac{1}{p} - \frac{1}{2^*} \right) \|u_n\|_{L^p}^p.
% \end{equation*}
% Since \( \{u_n\} \) is a (PS)\(_c\) sequence, we obtain
% \begin{equation*}
%     C \left(1 + \|u_n\|_{{\X,\varepsilon}}\right) \geq \frac{1}{N} \left(1 - \frac{\mu}{\bar{\mu}} \right) \|u_n\|_{\Xe}^2 
%     - \lambda \left( \frac{1}{p} - \frac{1}{2^*} \right) \|u_n\|_{\Xe}^p,
% \end{equation*}
% which implies that the sequence \( \{u_n\} \) is bounded. Thus, there exists \( u \in \X(\Omega) \) such that $u_n \rightharpoonup u$ weakly in $\X(\om)$, $ u_n \rightarrow u$ strongly in $L^r(\om)$ for $r \in [1, 2^*)$, and $u_n(x) \rightarrow u(x)$ pointwise a.e. in $\om$.
% \noi Now,
% \begin{align*}
%     &o_n(1)  = \langle (\Jep)^{\prime}(u_n), u_n - u \rangle \\
%     & = \left( \|\nabla u_n\|^2_2 - \|\nabla u\|^2_2 \right) + \varepsilon \left( [u_n]^2_{s} - [u]^2_{s} \right)
%     - \mu \int_{\Omega} \frac{|u_n|^2 - |u|^2}{|x|^2} \, dx - (\|u_n^+\|_{L^{2^*}}^{2^*} - \|u^+\|_{L^{2^*}}^{2^*}) + o_n(1) \\
%     & = \|\nabla u_n - \nabla u\|^2_2 + \varepsilon [u_n - u]^2_{s} 
%     - \mu \int_{\Omega} \frac{|u_n - u|^2}{|x|^2} \, dx - \|u_n^+ - u^+\|_{L^{2^*}}^{2^*} + o_n(1),
% \end{align*}
% where the final equality follows from the Brezis–Lieb Lemma. Thus, we conclude that
% \begin{equation*}
%     \lim_{n \to \infty} \left( \|u_n - u\|^2_{\X,\varepsilon} - \mu \int_{\Omega} \frac{|u_n - u|^2}{|x|^2} \, dx \right)
%     = \lim_{n \to \infty} \|u_n^+ - u^+\|_{L^{2^*}}^{2^*} = \ell  (\text{say}).
% \end{equation*}
% By \eqref{eq:spectral_ratio_frac}, we also have
% \begin{equation*}
%    S_{\mu} \|u_n^+ - u^+\|_{L^{2^*}}^2 \leq S_{\mu} \|u_n - u\|_{L^{2^*}}^2 \leq \|u_n - u\|^2_{\X,\varepsilon} - \mu \int_{\Omega} \frac{|u_n - u|^2}{|x|^2} \, dx.
% \end{equation*}
% This implies $ S_\mu \, \ell^{\frac{2}{2^*}} \leq \ell$.  If \( \ell > 0 \), then $\ell \geq S_\mu^{\frac{N}{2}}$.  \\
Following in the same way as in Lemma \ref{lem:ps_holds_superlinear}, we obtain
\begin{align*}
    \Jep(u_n) - \frac{1}{2} \langle (\Jep)^{\prime}(u_n), u_n \rangle 
    &= \left( \frac{1}{2} - \frac{1}{2^*} \right) \|u_n^+\|_{L^{2^*}}^{2^*} - \lambda \left( \frac{1}{p} - \frac{1}{2} \right)\int_{\Omega} |u_n^+|^p  dx.
\end{align*}
By the definition of (PS)\(_c\) sequence, H\"{o}lder's inequality, and Young’s inequality, we obtain
\begin{align*}
    c + o_n(1) &\geq \left( \frac{1}{2} - \frac{1}{2^*} \right) \left[ \ell + \|u^+\|_{L^{2^*}}^{2^*} \right]
    - \lambda \left( \frac{1}{p} - \frac{1}{2} \right) \|u^+\|_{L^p}^p \\
    &\geq \left( \frac{1}{2} - \frac{1}{2^*} \right) \left[ \ell + \|u^+\|_{L^{2^*}}^{2^*} \right]
    - \lambda \left( \frac{1}{p} - \frac{1}{2} \right) |\om|^{(2^*-p)/2^*} \|u^+\|_{L^{2^*}}^p \\
    &\geq \frac{\ell}{N}
    - |\Omega| \left(1 - \frac{ p}{2^*} \right) \left( \frac{pN}{2^*} \right)^{-p/(2^* - p)}
    \left[ \lambda \left( \frac{1}{p} - \frac{1}{2} \right) \right]^{2^*/(2^* - p)}.
\end{align*}
This contradicts the definition of \( c \), so \( \ell = 0 \). Hence, we obtained the desired result.
\end{proof}
Now, let us define 
\begin{equation}\label{eq:Lambda}
    \La = \sup \{ \la > 0 : \eqref{eq:main_problem_ep} \text{ has a weak solution} \}.
\end{equation}
Then we have the following lemma.
\begin{lemma}
Let $\La$ be defined as in \eqref{eq:Lambda}. Then $0 < \La < \infty$.
\end{lemma}
\begin{proof}
Take $\varphi \in \X(\om)$ such that $\|\varphi\|_{\X} = 1$. Since $\Jep(t \varphi) \to -\infty$ as $t \to \infty$, there exists $t_0 > 0$ independent of $\varepsilon$ and $\la$ (as term corresponding to $\la$ can be dropped) such that $\Jep(t_0 \varphi) < 0$. Then, for $0 < \rho \leq t_0$, we must have
\begin{equation*}
c_{\lambda} = \inf_{u \in \overline{B_{\rho}}} \Jep(u) < 0.
\end{equation*}
Then we obtain a (PS) sequence from the minimizing sequence via the Ekeland's Variational Principle \cite[Theorem 8.5]{willem2012minimax}. Thus, by the Lemma \ref{lem:ps_sublinear}, $\Jep$ achieves its minimum $c_{\lambda}$ at some nonnegative function $u_{\lambda}$ for $\la$ sufficiently small. Thus $\La > 0.$\\
Next we prove that $\La  < \infty$. Consider the eigenvalue problem 
\begin{equation*}
\left\{
\begin{aligned}
- \Delta \phi_1 + \varepsilon  (-\Delta)^s \phi_1 - \mu \frac{\phi_1}{|x|^2} &= {\lambda_{1,\varepsilon}} \phi_1, && \text{in } \Omega, \\
\phi_1 &= 0, && \text{in } \mathbb{R}^N \setminus \Omega,
\end{aligned}
\right. \tag{$\mathcal{P}_1$}
% \label{eq:auxillary_eq1} 
\end{equation*}
Let $u$ to be nonnegative solution of the problem \eqref{eq:main_problem_ep}. Then we have
\begin{equation} \label{eq:equality_of_soln}
\int_{\om} u^{2^*-1} \phi_1 dx + \la \int_{\om} u^{p-1} \phi_1 dx = \int_{\om} \grad u  \grad \phi_1 + \varepsilon \langle u , \phi_1 \rangle_s - \mu \int_{\om} \frac{u \phi_1}{|x|^2} dx = \lambda_{1,\varepsilon} \int_{\om} u \phi_1 dx.
\end{equation}
We note that there exists $\la_*$ such that for all $\la > \la_*$ we have 
\begin{equation*}
    t^{2^* - 1} + \la t^{p-1} > \la_{1,\varepsilon} t 
\end{equation*}
for all $t > 0$. Combining this with \eqref{eq:equality_of_soln} we deduce that $\La < \infty$.
\end{proof}
\begin{lemma}\label{lem:existence_each_lambda}
    The problem \eqref{eq:main_problem_ep} has at least one positive solution for each $\la \in (0, \La)$. In fact, the sequence of minimal solutions $\{u_\la \}$ are increasing with respect to $\la$. If $\la = \La$, then the problem \eqref{eq:main_problem_ep} admits at least one weak solution.
\end{lemma}
\begin{proof}
Let \(\tilde\lambda \in (0, \La)\) be fixed. By definition of $\La$, there exists \(\lambda_{*} \in (\tilde\lambda, \Lambda)\) such that \eqref{eq:main_problem_ep} with $\la = \la_*$ has a positive weak solution \(\bar u\).  
Then \(\bar  u\) is a supersolution of the problem \eqref{eq:main_problem_ep} with $\la = \tilde \la$. To obtain the subsolution, we define the following eigenvalue problem
\begin{equation}
\begin{cases}
-\Delta \varphi_1 + \varepsilon(-\Delta)^{s} \varphi_1 = \lambda_{\varepsilon}^1 \varphi_1 & \text{in } \Omega, \\[6pt]
\varphi_1 = 0 & \text{in } \mathbb{R}^N \setminus \Omega,
\end{cases}
\label{eq:P2}
\tag{P2}
\end{equation}
where \(\lambda_{\varepsilon}^1\) is the first eigenvalue of the mixed operator \(-\Delta + \varepsilon(-\Delta)^s\) with eigenfunction $\varphi_1 \in L^{\infty}(\om)$. Then \(\varphi_t = t \varphi_1\) also solves \eqref{eq:P2}.  
Choosing \(t > 0\) sufficiently small such that
\begin{equation*}
-\Delta \varphi_t + \varepsilon (-\Delta)^s \varphi_t = \lambda_{\varepsilon}^1 \varphi_t 
\leq \tilde\lambda \varphi_t^{p-1} \leq \tilde\lambda \varphi_t^{p-1} + \varphi_t^{2^*-1} + \mu \frac{\varphi_t}{|x|^2}.
% \label{eq:subsolution_estimate}
\end{equation*}
Hence, \(\varphi_t\) is a subsolution of \eqref{eq:main_problem_ep}. 
Next, taking \((\varphi_t - \bar u)^+\) as a test function, we obtain
\begin{equation*}
\begin{aligned}
0 &\leq \int_{\Omega} |\nabla (\varphi_t - \bar u)^+|^2 \, dx 
+  \varepsilon \big[ (\varphi_t - \bar u)^+ \big]^2_{s} \\
&\leq \int_{\Omega} \nabla (\varphi_t - \bar u)  \nabla (\varphi_t - \bar u)^+ \, dx 
+ \varepsilon \langle \varphi_t - \bar u, (\varphi_t - \bar u)^+ \rangle_{s} \\[6pt]
&= \lambda_{\varepsilon}^1 \int_{\Omega} \varphi_t (\varphi_t - \bar u)^+ \, dx
- \lambda_* \int_{\Omega} \bar u^{p-1} (\varphi_t - \bar u)^+ \, dx \\[6pt]
&\quad - \int_{\Omega} \bar u^{2^*-1} (\varphi_t - \bar u)^+ \, dx
- \mu \int_{\Omega} \bar u \, \frac{(\varphi_t - \bar u)^+}{|x|^2} \, dx\\
&\leq \lambda_* \int_{\Omega} \Big( \varphi_t^{p-1} - \bar u^{p-1} \Big) (\varphi_t - \bar u)^+ \, dx \leq 0,
\end{aligned}
% \label{eq:comparison_expansion}
\end{equation*}
which yields $\varphi_t \leq \bar u$. Now, let \(\{u_k\}\) denote the nonnegative sequence in \(\X(\Omega)\) of solutions to the iterated problem
\begin{equation*}
(P_k) \; \; \; 
\begin{cases}
-\Delta u_k + \varepsilon (-\Delta)^s u_k = \mu \dfrac{u_{k-1}}{|x|^2} + \tilde\lambda u_{k-1}^{p-1} + u_{k-1}^{2^*-1} & \text{in } \Omega, \\[8pt]
u_k = 0 & \text{in } \mathbb{R}^N \setminus \Omega
\end{cases}
% \label{eq:iterated_problem}
\end{equation*}
for $k \geq 1$ and $u_0 = \varphi_t$. Then we claim that
\[
\varphi_t \leq u_1 \leq u_2 \leq \cdots \leq \overline{u} .
\]
We prove this by principle of mathematical induction. Using the same idea as above, it is clear that \(\varphi_t \leq u_1\). Assume that the claim is true for $k$ and then we prove it for $k+1$, i.e., $u_{k} \leq u_{k+1}$ whenever $u_{k-1} \leq u_k$ for all $k \in \ntrl$. Take \((u_{k} - u_{k+1})^+\) as a test function in \((P_k)\) and \((P_{k+1})\) to get 
\begin{equation*}
\begin{aligned}
0 &\leq \int_{\Omega} |\nabla (u_{k} - u_{k+1})^+|^2 \, dx 
+ \varepsilon \big[ (u_{k} - u_{k+1})^+ \big]^2_{s} \\[6pt]
&\leq \mu \int_{\Omega} \frac{(u_{k-1} - u_{k})(u_{k} - u_{k+1})^+}{|x|^2} \, dx + \tilde\lambda \int_{\Omega} \big(u_{k-1}^{p-1} - u_{k}^{p-1}\big) (u_{k} - u_{k+1})^+ \, dx \\[6pt]
&\quad + \int_{\Omega} \big(u_{k-1}^{2^*-1} - u_{k}^{2^*-1}\big)(u_{k} - u_{k+1})^+ \, dx \leq 0,
\end{aligned}
% \label{eq:induction_step}
\end{equation*}
where the last inequality follows from induction hypothesis \(u_{k-1} \leq u_{k}\). Then it follows that $u_{k} \leq u_{k+1}$.
% Finally, consider the limit function \(\overline{u}\). We have
% \begin{equation*}
% \begin{aligned}
% 0 &\leq \int_{\Omega} |\nabla (u_{k} - \overline{u})^+|^2 \, dx 
% + \big[ (u_{k} - \overline{u})^+ \big]^2_{s} \\[6pt]
% &\leq \mu \int_{\Omega} \frac{(u_{k-1} - \overline{u})(u_{k} - \overline{u})^+}{|x|^2} \, dx + \lambda_* \int_{\Omega} \big(u_{k-1}^{p-1} - \overline{u}^{p-1}\big)(u_{k-1} - \overline{u})^+ \, dx \\
% &\quad + \int_{\Omega} \big(u_{k-1}^{2^*-1} - \overline{u}^{2^*-1}\big)(u_{k-1} - \overline{u})^+ \, dx \leq 0.
% \end{aligned}
% % \label{eq:upper_bound}
% \end{equation*}
Since \(u_{k} \leq \overline{u}\), it follows in same way that $u_{k+1} \leq \overline{u}$. Hence, our claim is proved. \\
Next, define
\[
u_{\tilde\lambda} = \lim_{k \to \infty} u_k \quad \text{in } L^1(\Omega).
\]
Then 

\begin{equation*}
\begin{aligned}
\|u_k\|^2_{\Xe} 
&= \mu \int_{\Omega} \frac{u_k u_{k-1}}{|x|^2} \, dx 
+ \int_{\Omega} u_k u_{k-1}^{2^*-1} \, dx 
+ \tilde\lambda \int_{\Omega} u_k u_{k-1}^{p-1} \, dx \\[6pt]
&\leq \mu \int_{\Omega} \frac{\overline{u}^2}{|x|^2} \, dx
+ \int_{\Omega} \overline{u}^{2^*} \, dx 
+ \tilde\lambda \int_{\Omega} \overline{u}^p \, dx \leq C.
\end{aligned}
% \label{eq:uk_energy_bound}
\end{equation*}
Thus, up to a subsequence, we know that \(u_k \rightharpoonup u_{\tilde\lambda}\) in \(\X(\Omega)\).  
Therefore, by passing to the limit in \((P_k)\), we conclude that \(u_{\tilde\lambda} \geq 0\) is the solution of the problem \eqref{eq:main_eq}. \\
For $\la = \La$, there exists a sequence of minimal solutions $\{ u _{\la_n}\}$ of the problem \eqref{eq:main_problem_ep}
 with $\la_n \uparrow \La$. By construction they are increasing with respect to $\la$. Let $\varphi_1$ be the solution of the problem \eqref{eq:P2}, then we have 
 \begin{equation}\label{eq:bound_minimal_seq}
 \begin{aligned}
 \mu \int_{\Omega} \frac{ u_{\la_n} \varphi_1}{|x|^2} \, dx 
+ \int_{\Omega} u_{\la_n}^{2^*-1} \varphi_1 \, dx 
+ \lambda_n \int_{\Omega} u_{\la_n}^{p-1} \varphi_1 \, dx = \la_{\varepsilon}^1 \int_{\om} u_{\la_n}\varphi_1  dx\leq \frac{1}{2} \int_{\Omega} u_{\la_n}^{2^*-1} \varphi_1 \, dx + C \int_{\om} \varphi_1 \, dx
\end{aligned}
 \end{equation}
 Let $\delta(x):= \text{dist}(x, \partial \Omega)$. Then by Lemma 2.2 of \cite{biswas2023boundary} there exists a bounded barrier function $\psi$ such that
\begin{equation*}
   \begin{cases}
       - \Delta \psi + \varepsilon (-\Delta)^s \psi \leq 0 &\text{ in } B_{4r} \setminus \bar{B}_r,\\
       0 \leq \psi \leq \tilde\ka r &\text{ in } B_r,\\
       \psi(x) \geq \frac{1}{\tilde \ka} \delta(x) & \text{ in } B_{4r} \setminus B_r,\\
       \psi \leq 0 & \text{ in } B_{4r}^C,
   \end{cases}
\end{equation*}
where $\tilde \ka$ is a constant independent of boundary of $B_{4r}$. 
 Using the comparison principle in \cite[Theorem 5.2]{biswas2025mixed} to compare $\varphi_1$ and $\psi$, we get ${\frac{\varphi_1}{\delta}(x_0)} \geq \frac{\psi}{\delta}(x_0)\geq \frac{1}{\tilde \ka}> 0$ at each $x_0 \in \partial \Omega$. But then Theorem 1.2 of \cite{biswas2023boundary} gives $\inf{\left( \frac{\varphi_1}{\delta}\right)} > 0$. So, using this and \eqref{eq:bound_minimal_seq}, we deduce that
 \begin{equation}\label{eq:min_seq_bd2}
 \begin{aligned}
     \inf{\left( \frac{\varphi_1}{\delta}\right)} &\left( \mu \int_{\Omega} \frac{ u_{\la_n} \delta}{|x|^2} \, dx + \int_{\Omega} u_{\la_n}^{2^*-1} \delta \, dx 
+ \lambda_n \int_{\Omega} u_{\la_n}^{p-1} \delta \, dx\right) \\
& \leq \mu \int_{\Omega} \frac{ u_{\la_n} \varphi_1}{|x|^2} \, dx 
+ \int_{\Omega} u_{\la_n}^{2^*-1} \varphi_1 \, dx 
+ \lambda_n \int_{\Omega} u_{\la_n}^{p-1} \varphi_1 \, dx \leq C .
\end{aligned}
 \end{equation}
 Let $\psi_1$ be the solution to the linear problem 
 \begin{equation*}
\begin{cases}
-\Delta \psi_1 + \varepsilon (-\Delta)^s \psi_1 = 1 & \text{in } \Omega, \\[8pt]
\psi_1 = 0 & \text{in } \mathbb{R}^N \setminus \Omega.
\end{cases}
\end{equation*}
Then taking $\psi_1$ as a test function in \eqref{eq:main_problem_ep} with $\la = \la_n$ and using Theorem $2.3$ of \cite{dhanya2026interior} and \eqref{eq:min_seq_bd2}, we get
\begin{align*}
    \int_{\om} u_{\la_n}dx & = \mu \int_{\Omega} \frac{ u_{\la_n} \psi_1}{|x|^2} \, dx 
+ \int_{\Omega} u_{\la_n}^{2^*-1} \psi_1 \, dx 
+ \lambda_n \int_{\Omega} u_{\la_n}^{p-1} \psi_1 \, dx \\
&  \leq C \left( \mu \int_{\Omega} \frac{ u_{\la_n} \delta}{|x|^2} \, dx + \int_{\Omega} u_{\la_n}^{2^*-1} \delta \, dx 
+ \lambda_n \int_{\Omega} u_{\la_n}^{p-1} \delta \, dx \right)\\
& \leq C.
\end{align*}
Hence the sequence $\{ u_{\la_n}\}$ is uniformly bounded in $L^1(\om)$ and converges to a function $u_{\La} \geq 0$. By monotone convergence theorem it follows that $u_{\La}$ is a weak solution.
 \end{proof}

 
\begin{lemma}
Let $\La$ be defined as in \eqref{eq:Lambda}. Then for $\la_0 \in (0,\La)$, either the problem \eqref{eq:main_problem_ep} has two distinct solutions or there exists a local minimum of functional $\Jepo$ in $\X(\om)$.
\end{lemma}
\begin{proof}
Let \(\lambda_0 \in (0,\Lambda)\) and fix \(\lambda_1 \in (\lambda_0, \Lambda)\) be fixed. Then by Lemma \ref{lem:existence_each_lambda} we get minimal solutions $u_{\lambda_0}$ and $u_{\la_1}$ of \eqref{eq:main_problem_ep} with $\lambda_0 $ and $\la_1$ respectively such that $u_{\lambda_0} \leq u_{\lambda_1}$. \\
Define the set $Z = \{ u \in \X(\Omega) : 0 \leq u \leq u_{\lambda_1} \}$. Then \(Z\) is a bounded, closed, and convex subset of \(\X(\Omega)\). Thus there exists $u_0 
\in Z$ such that
\begin{equation*}
\Jepo(u_0) = \inf_{u \in Z} \Jepo (u).
% \label{eq:J_lambda_min}
\end{equation*}
Since \(1 < p < 2\), we have $\Jepo(t u_{\lambda_1}) < 0$ for some sufficiently small $t > 0$. Moreover \(t u_{\lambda_1} \in Z\), it follows that \(u_0 \neq 0\).\\
If $u_0 \neq u_{\la_0}$, then we get two distinct solutions of \eqref{eq:main_problem_ep}. Otherwise, if $u_0 =  u_{\la_0}$, then we claim that \(u_0\) is a local minimum of the functional \(\Jepo\) in $\X(\om)$. We prove this by contradiction. Suppose, on the contrary, that \(u_0\) is not a local minimum in $\X(\om)$. Then there exists a sequence \(\{u_n\} \subset \X(\om)\) such that
\begin{equation*}
\lim_{n \to \infty} \|u_n - u_0\|_{\Xe} = 0
\quad \text{and} \quad
\Jepo(u_n) < \Jepo(u_0).
% \label{eq:contradiction_sequence}
\end{equation*}
Now, consider the auxiliary functions $w_n \in \X(\om)$, and $v_n \in Z$ as
\begin{equation*}
w_n = (u_n - u_{\lambda_1})^+ \quad \text{and} \quad v_n(x) =
\begin{cases}
0, & u_n(x) \leq 0, \\[6pt]
u_n(x), & 0 \leq u_n(x) \leq u_{\lambda_1}(x), \\[6pt]
u_{\lambda_1}(x), & u_{\lambda_1}(x) \leq u_n(x).
\end{cases}
% \label{eq:w_n_def}
\end{equation*}
Also, take the sets
\begin{equation*}
A_n = \{ x : v_n(x) = u_n(x) \}, 
\qquad 
B_n = \{ x : u_n(x) \geq u_{\lambda_1}(x) \},
\qquad 
\widetilde{A}_n = A_n \cap \Omega, 
\qquad 
\widetilde{B}_n = B_n \cap \Omega.
% \label{eq:Tn_Sn_def}
\end{equation*}

Next, consider the function
\begin{equation*}
H_{\lambda_0}(t) = \frac{1}{2^*} t^{2^*}_+ + \frac{\lambda_0}{p} t^p_+.
% \label{eq:F_lambda_def}
\end{equation*}
Moreover, to simplify notation, we set
\begin{equation*}
U_n(x,y) = \frac{|u_n(x) - u_n(y)|^2}{|x-y|^{N+2s}},
\qquad
W_n(x,y) = \frac{|w_n(x) - w_n(y)|^2}{|x-y|^{N+2s}},
% \label{eq:U_V_def}
\end{equation*}
\begin{equation*}
U_n^{\pm}(x,y) = \frac{|(u_n)^\pm(x) - (u_n)^\pm(y)|^2}{|x-y|^{N+2s}},
\qquad V_n(x,y) = \frac{|v_n(x) - v_n(y)|^2}{|x-y|^{N+2s}}.
% \label{eq:Uplus_Vplus_def}
\end{equation*}
Then using $U_n(x,y) = U_n^+(x,y) + U_n^-(x,y) + 2 \frac{u_n^+(x) u_n^-(y)+ u_n^+(y) u_n^-(x)}{|x-y|^{N+2s}}$, we obtain
\begin{equation*}
\begin{aligned}
\Jepo(u_n) \;\geq\; & 
\frac{1}{2} \int_{\Omega} |\grad u_n^+|^2 \, dx 
+ \frac{1}{2} \int_{\Omega} |\nabla u_n^-|^2 \, dx
+ \frac{\varepsilon}{2} \iint_{\mathbb{R}^{2N}} U_n^{+}(x,y) \, dx \, dy  \\
& \quad + \frac{\varepsilon}{2} \iint_{\mathbb{R}^{2N}} U_n^-(x,y) \, dx \, dy
- \frac{\mu}{2} \int_{\Omega} \frac{(u_n^+)^2}{|x|^2} \, dx
- \int_{\widetilde{A}_n} H_{\lambda_0}(v_n) \, dx
- \int_{\widetilde{B}_n} H_{\lambda_0}(u_n) \, dx.\\
\geq&\; \frac{\varepsilon}{2}\,[u_n^- ]_{s}^2 
+ \Jepo(v_n) 
+ \frac{1}{2}\int_{\Omega} |\nabla u_n^-|^2 \, dx 
+ \frac{1}{2}\int_{\widetilde{B}_n} \big(|\nabla u_n^+|^2 - |\nabla v_n|^2\big) dx \\
&\quad + \frac{\varepsilon}{2} \iint_{B_n \times B_n} 
 \big[U_n^+(x,y)-V_n(x,y)\big] \, dx dy +  \varepsilon\int_{B_n} \int_{B_n^C} 
\big[U_n^+(x,y)-V_n(x,y)\big]  \, dy dx \\
&\quad - \frac{\mu}{2} \int_{\widetilde{B}_n} \frac{((u_n^+)^2 - v_n^2)}{|x|^2} \, dx
+ \int_{\widetilde{B}_n} \big( H_{\lambda_0}(v_n) - H_{\lambda_0}(u_n) \big) dx .
\end{aligned}
% \label{eq:J_lambda_bound}
\end{equation*}
Since $w_n(x) = u_n(x) - u_{\lambda_1}(x)$ whenever $x \in B_n$, using $(37)-(40)$ of Lemma $7$ from \cite{shen_multiplicity_fractional} we obtain
\begin{equation*}
\begin{aligned}
\Jepo(u_n) 
&\;\geq\; \frac{\varepsilon}{2}[u_n^-]_s^2 
+ \frac{1}{2}\int_{\Omega} |\nabla u_n^-|^2 \, dx 
+ \Jepo(v_n) 
+ \frac{\varepsilon}{2}[w_n]_s^2 + \frac{1}{2}\int_{\widetilde{B}_n} |\nabla w_n|^2 \, dx\\
&\quad 
+  \int_{\widetilde{B}_n} \nabla w_n  \nabla u_{\lambda_1} \, dx + \varepsilon\iint_{B_n \times B_n} (w_n(x)-w_n(y))(u_{\lambda_1}(x)-u_{\lambda_1}(y)) \, d\nu \\
&\quad - 2\varepsilon \iint_{B_n \times B_n^c}
w_n(y)\big(u_n^+(x)-u_{\lambda_1}(y)\big) \, d\nu  - \frac{\mu}{2} \int_{\widetilde{B}_n} \frac{w_n^2(x) + 2 w_n(x) u_{\lambda_1}(x)}{|x|^2} \, dx \\
&\quad + \int_{\widetilde{B}_n} \big( H_{\lambda_0}(v_n) - H_{\lambda_0}(u_n) \big) dx \\
&\;\geq\; \frac{\varepsilon}{2}[u_n^-]_s^2 
+ \frac{1}{2}\int_{\Omega} |\nabla u_n^-|^2 \, dx 
+ \Jepo(v_n) 
+ \frac{\varepsilon}{2}[w_n]_s^2 + \frac{1}{2}\int_{\widetilde{B}_n} |\nabla w_n|^2 \, dx\\
&\quad 
+  \int_{\widetilde{B}_n} \nabla w_n  \nabla u_{\lambda_1} \, dx + \varepsilon \langle u_{\la_1}, w_n \rangle_s  - \frac{\mu}{2} \int_{\widetilde{B}_n} \frac{w_n^2(x) + 2 w_n(x) u_{\lambda_1}(x)}{|x|^2} \, dx \\
&\quad + \int_{\widetilde{B}_n} \big( H_{\lambda_0}(v_n) - H_{\lambda_0}(u_n) \big) dx,
\end{aligned}
% \label{eq:J_lambda0_secondbound}
\end{equation*}
where the last inequality follows from the fact that $Supp (w_n) = B_n$ and $u_n^+(x) \leq u_{\la_1}(x)$ for $x \in B_n^c$.
Since $u_{\lambda_1}$ is a supersolution of $(P_{\mu, \lambda_0}^{\varepsilon})$, we deduce
\begin{equation*}
\int_{\Omega} \nabla u_{\lambda_1}  \nabla w_n \, dx 
+ \varepsilon \langle u_{\lambda_1}, w_n \rangle_s 
\;\geq\; 
\mu \int_{\Omega} \frac{u_{\lambda_1} w_n}{|x|^2} \, dx 
+ \int_{\widetilde{B}_n} w_n H'_{\lambda_0}(u_{\lambda_1}) \, dx .
% \label{eq:super_solution_condition}
\end{equation*}
Thus we have
\begin{align}
\Jepo(u_n) \;\geq\;& \Jepo(v_n) + \frac{\varepsilon}{2} [w_n]_{s}^{2} 
+ \frac{1}{2} \int_{\widetilde{B}_n} |\nabla w_n|^2 
 +\int_{\widetilde{B}_n} \Bigg[ \frac{1}{2^*} \Big( u_{\lambda_1}^{2^*} - (u_{\lambda_1}+w_n)^{2^*} \Big) 
+ w_n \, u_{\lambda_1}^{2^*-1} \Bigg] dx \notag \\
& \quad - \frac{\mu}{2} \int_{\widetilde{B}_n} \frac{w_n^2(x)}{|x|^2} \, dx  + \lambda_0 \int_{\widetilde{B}_n} 
\frac{u_{\lambda_1}^p - (u_{\lambda_1}+w_n)^p}{p} 
+ w_n \, u_{\lambda_1}^{p-1} \, dx .\label{eq:J_lambda_expanded} 
\end{align}
Using Talyor's expansion and the fact that \( p \in (1,2) \), we obtain
\begin{equation*}
0 \;\leq\; \frac{1}{p} \Big[ (u_{\lambda_1}+w_n)^p - u_{\lambda_1}^p \Big] - w_n u_{\lambda_1}^{p-1} 
\;\leq\; \frac{(p-1)}{2} \, \frac{w_n^2}{u_{\lambda_1}^{2-p}}.
% \label{eq:p_range_ineq}
\end{equation*}
Moreover, using AM-GM inequality
\begin{equation*}
\big(u_{\lambda_1}(x) - u_{\lambda_1}(y)\big) 
\left( \frac{w_n^2(x)}{u_{\lambda_1}(x)} - \frac{w_n^2(y)}{u_{\lambda_1}(y)} \right) \leq |w_n(x) - w_n(y)|^2 .
% \label{eq:diff_ineq}
\end{equation*}
Since \( u_{\lambda_1} \) is a solution of \eqref{eq:main_problem_ep} with $\la = \la_1$, using the above inequality together with Picone’s identity, we deduce
\begin{equation*}
\begin{aligned}
\lv w_n \rv_{\Xe}^2\;\geq\; \varepsilon \Big\langle u_{\lambda_1}, \frac{w_n^2}{u_{\lambda_1}} \Big\rangle_s 
+ \int_{\Omega} \nabla \!\left( \frac{w_n^2}{u_{\lambda_1}} \right) \nabla u_{\lambda_1} \, dx \;\geq\; \mu \int_{\Omega} \frac{w_n^2}{|x|^2} \, dx 
+ \lambda_0 \int_{\Omega} \frac{w_n^2}{u_{\lambda_1}^{2-p}} \, dx .
\end{aligned}
% \label{eq:picone_application}
\end{equation*}
Thus, we obtain
\begin{equation}
\begin{aligned}
\lv w_n \rv_{\Xe}^2 - \mu \int_{\Omega} \frac{w_n^2}{|x|^2} \geq \lambda_0 \int_{\Omega} \frac{w_n^2}{u_{\lambda_1}^{2-p}}  \;\geq\; \frac{2 \la_0}{p-1} 
\int_{\widetilde{B}_n} 
\left[ \frac{(u_{\lambda_1}+w_n)^p - u_{\lambda_1}^p}{p} 
- w_n u_{\lambda_1}^{p-1} \right].
\end{aligned}
\label{eq:ineq46}
\end{equation}
Furthermore, we have the following inequality
\begin{equation}
0 \;\leq\; \frac{1}{2^*} \Big[ (u_{\lambda_1}+w_n)^{2^*} - u_{\lambda_1}^{2^*} \Big] - w_n u_{\lambda_1}^{2^*-1}
\;\leq\; C \left( u_{\lambda_1}^{2^*-2} w_n^2 + w_n^{2^*} \right).
\label{eq:ineq47}
\end{equation}
Plugging \eqref{eq:ineq46} and \eqref{eq:ineq47} into \eqref{eq:J_lambda_expanded}, we deduce
\begin{equation*}
\begin{aligned}
\Jepo(u_n) 
&\geq\Jepo(v_n) 
+ \frac{2-p}{2} 
\left[ \varepsilon [w_n]_s^2 + \int_{\widetilde{B}_n} |\nabla w_n|^2  
- \mu \int_{\widetilde{B}_n} \frac{w_n^2(x)}{|x|^2} \right] - C \int_{\Omega} \Big( u_{\lambda_1}^{2^*-2} w_n^2 + w_n^{2^*} \Big)\\
& \geq \Jepo(v_n) + C \|w_n\|_{\Xe}^2 
- C \int_{\widetilde{B}_n} \Big( u_{\lambda_1}^{2^*-2} w_n^2 + w_n^{2^*} \Big) dx .
\end{aligned}
% \label{eq:J_lambda_lowerbound}
\end{equation*}
Moreover, using the idea of Lemma $7$ of \cite{shen_multiplicity_fractional} and $u_0 = u_{\la_0}$, we obtain
\begin{equation*}
\lim_{n \to \infty} |\widetilde{B}_n| = 0 \quad \text{and} \quad \int_{\widetilde{B}_n} \Big( u_{\lambda_1}^{2^*-2} w_n^2 + w_n^{2^*} \Big) dx \leq o_n(1) \lv w_n\rv_{\Xe}^2.
% \label{eq:limit_Sn}
\end{equation*}
Hence, we finally conclude
\begin{equation*}
\begin{aligned}
\Jepo(u_n) 
&\;\geq\; \Jepo(v_n) + \big(C - o_n(1)\big)\|w_n\|_{\Xe}^2 \geq  \Jepo(u_0),
\end{aligned}
% \label{eq:J_lambda_final}
\end{equation*}
for large $n$. This gives us the contradiction. Thus, $u_0$ is a local minimum for $\Jepo$.
\end{proof}
Next, we examine the existence of a second positive solution of the form
\begin{equation*}
u = u_{\lambda} + v.
\end{equation*}
The equation for $v$ becomes
\begin{equation}
\left\{
\begin{aligned}
- \Delta v + \varepsilon  (-\Delta)^s v - \mu \frac{v}{|x|^2} &= \lambda (u_{\lambda} + v)^{p-1} + (u_{\lambda} + v)^{2^* - 1} - \lambda u_{\lambda}^{p-1} - u_{\lambda}^{2^* - 1}, && \text{in } \Omega, \\
v &= 0, && \text{in } \mathbb{R}^N \setminus \Omega,
\end{aligned}
\right.
\label{eq:second_solution_equation}
\end{equation}
Define the function \( g(x,t) \) as
\begin{equation*}
g(x,t) = 
\begin{cases}
(u_{\lambda} + t)^{2^*-1} - u_{\lambda}^{2^*-1} + \lambda \left[ (u_{\lambda} + t)^{p-1} - u_{\lambda}^{p-1} \right], & \text{if } t > 0, \\
0, & \text{if } t \leq 0,
\end{cases}
\end{equation*}
and its primitive
\begin{equation*}
G(v) = \int_0^v g(x,t) \, dt.
\end{equation*}
The associated energy functional becomes
\begin{equation*}
\Ie(v) = \frac{1}{2} \left[ \int_{\Omega} |\nabla v|^2 + \varepsilon  [v]_{s}^2 - \mu \int_{\Omega} \frac{v^2}{|x|^2} \right] - \int_{\Omega} G(v) \, dx.
\end{equation*}
There is a one-to-one correspondence between the critical points of $\Ie$ in $\X$ and the weak solutions of problem \eqref{eq:second_solution_equation}.
We now prove the existence of a nontrivial solution \( v \) using a contradiction argument. Assume that \( v = 0 \) is the only solution of \eqref{eq:second_solution_equation}.

\begin{lemma}
The functional \( \Ie \) has a local minimum at $v =0$ in \( \X(\om) \).
\end{lemma}
\begin{proof}
Let \( w \in \X (\om)\). Then using the fact that \( \langle (\Je)^{\prime}(u_{\lambda}), w^+ \rangle = 0 \) and rearranging terms, we deduce
\begin{equation*}
\begin{aligned}
\Ie(w) &= \frac{1}{2} \int_{\Omega} |\nabla w|^2 - \mu \int_{\Omega} \frac{w^2}{|x|^2} \, dx + \frac{\varepsilon }{2} [w]^2_s - \frac{1}{2^*} \int_{\Omega} \left[ (u_{\lambda} + w^+)^{2^*} - u_{\lambda}^{2^*} - 2^* u_{\lambda}^{2^* - 1} w^+ \right] \\
&\quad - \frac{\lambda}{ p} \int_{\Omega} \left[ (u_{\lambda} + w^+)^{p} - u_{\lambda}^{p} - p u_{\lambda}^{p-1} w^+ \right] \\
&  = \Je(u_{\lambda} + w^+) - \Je(u_{\lambda})  + \frac{1}{2} \int_{\Omega} |\nabla w^-|^2 - \mu \int_{\Omega} \frac{(w^-)^2}{|x|^2} + \frac{\varepsilon }{2} [w^-]_s^2 + \varepsilon  \langle w^+, w^- \rangle_s.
\end{aligned}
\end{equation*}
Since \( u_{\lambda} \) is a local minimizer and
\begin{equation*}
\frac{1}{2} [w^-]_s^2 + \langle w^+, w^- \rangle_s = \frac{1}{2} [w]_s^2 - \frac{1}{2} [w^+]_s^2 \geq 0,
\end{equation*}
we obtain the inequality
\begin{equation*}
\Ie(w) \geq \frac{1}{2} \left( \int_{\Omega} |\nabla w^-|^2 - \mu \int_{\Omega} \frac{(w^-)^2}{|x|^2} \right),
\end{equation*}
for \( \|w\|_{\X, \varepsilon} \) small enough. This concludes the proof.
\end{proof}

\begin{lemma}
If \( v = 0 \) is the only critical point of  \( \Ie \), then \( \Ie \) satisfies the Palais–Smale condition (PS)\(_c\) for any \( c < \frac{1}{N} S_{\mu}^{N/2} \).
\end{lemma}
\begin{proof}
Let \( \{v_n\} \subset \X(\om) \) be a (PS)\(_c\) sequence. Then
\begin{equation*}
\Ie(v_n) \to c < \frac{1}{N} S_{\mu}^{N/2}, \quad  (\Ie)^{\prime}(v_n) \to 0 \quad \text{in } (\X)'.
\end{equation*}
Then using Hardy inequality and H\"{o}lder inequality, we compute
\begin{equation*}
\begin{aligned}
c + o_n(1)& \|u_{\lambda} + v_n\|_{\X, \varepsilon}  \geq 2^* \Ie(v_n) - \langle (\Ie)^{\prime}(v_n),u_{\la}+v_n\rangle\\
&= \left( \frac{2^*}{2} - 1 \right) \left[ \int_{\Omega} |\nabla v_n|^2 - \mu \int_{\Omega} \frac{v_n^2}{|x|^2} + \varepsilon [v_n]^2_s \right] + (2^* - 2) \int_{\Omega} u_{\lambda}^{2^* - 1} v_n^+   \\
&\quad +\int_{\Omega} u_{\lambda}^{2^* - 1} v_n^-+ \la\left( \frac{2^*}{ p} - 1 \right) \int_{\om} u_{\la}^{p} + \la(2^*-2) \int_{\Omega} u_{\lambda}^{p-1} v_n^+ \\
& \quad+ \lambda \left( \frac{p - 2^*}{ p} \right) \int_{\Omega} (u_{\lambda} + v_n^+)^{p} + \lambda \int_{\Omega} v_n^{p-1} v_n^-.\\
&\geq \left( \frac{2}{N - 2} \right) \left( 1 - \frac{\mu}{\bar{\mu}} \right) \left( \int_{\Omega} |\nabla v_n|^2 + \varepsilon  [v_n]_s^2 \right) 
+ C \lambda \frac{(p  - 2^*)}{ p} \|u_{\lambda} + v_n^+\|_{\X}^{p }.
\end{aligned}
\end{equation*}
Since \( 1 < p < 2 \), this implies that \( \{v_n\} \) is bounded in \( \X \). Thus, passing to a subsequence if necessary, we have 
% $v_n \rightharpoonup v_0$ weakly in $\X(\om)$, $v_n(x) \to v_0(x)$ pointwise a.e. in $\om$, and $v_n \to v_0 $ in $L^{t}(\om)$ for all $ 1 < t < 2^*$
\begin{equation}\label{convergence_properties}
\begin{aligned}
v_n &\rightharpoonup v_0 \quad \quad \;\text{in } \X,\\
v_n(x) &\to v_0(x) \quad \text{a.e. in } \Omega,\\
v_n &\to v_0 \quad \quad \;\text{in } L^t(\Omega) \text{ for all } 1 < t < 2^*.
\end{aligned}
\end{equation}
Clearly, \( v_0 \) is a critical point of \( \Ie \). But our assumption implies \( v_0 = 0 \).\\
Using the Brézis–Lieb Lemma and \eqref{convergence_properties}, we deduce
\begin{equation}\label{eq:reduced_I_functional}
\Ie(v_n) = \frac{1}{2} \int_{\Omega} |\nabla v_n|^2 - \frac{\mu}{2} \int_{\Omega} \frac{v_n^2}{|x|^2} + \frac{\varepsilon }{2} [v_n]_s^2 - \frac{1}{2^*} \int_{\Omega} (v_n^+)^{2^*} + o_n(1).
\end{equation}
Also, we compute
\begin{equation*} 
\left\langle (\Ie)^{\prime}(v_n), u_\la +v_n \right\rangle = \int_{\Omega} |\nabla v_n|^2 + \varepsilon  [v_n]_s^2 - \mu \int_{\Omega} \frac{v_n^2}{|x|^2} + \int_{\Omega} (v_n^+)^{2^*} + o_n(1).
\end{equation*}
Now, we can assume that
\begin{equation*}
\int_{\Omega} |\nabla v_n|^2 + \varepsilon  [v_n]_s^2 - \mu \int_{\Omega} \frac{v_n^2}{|x|^2} \to b,
\end{equation*}
and
\begin{equation*}
\int_{\Omega} (v_n^+)^{2^*} \to b, \quad \text{as } n \to \infty.
\end{equation*}

If \( b = 0 \), the proof is complete. If \( b \neq 0 \), then by the definition of the Sobolev constant \( S_{\mu} \), we have
\begin{equation*} 
b \geq S_{\mu} b^{2 / 2^*},
\end{equation*}
i.e., $b \geq S_{\mu}^{N/2}$. Thus, using \eqref{eq:reduced_I_functional} we obtain the inequality
\begin{equation*}
c = o_n(1) + \Ie(v_n) = \frac{b}{N} + o_n(1) \geq \frac{S_{\mu}^{N/2}}{N} + o_n(1),
\end{equation*}
which contradicts our earlier assumption on the energy level \( c \).
\end{proof}



\begin{proof}[Proof of Theorem \ref{thm:sublinear_two_solutions}]
We define the min–max value as
\begin{equation*}
c^*_{\lambda} = \inf_{h \in \Gamma} \max_{t \in [0,1]} \Ie(h(t)),
\end{equation*}
where $\Gamma = \left\{ h \in C([0,1], \X(\om)) : h(0) = 0,\, h(1) = t_0 u_{\varepsilon,\al} \right\}$, and \( t_0 \) is chosen, independent of $\varepsilon$, such that $\Ie(t_0 u_{\varepsilon,\al}) < 0$.

We aim to show the existence of a nontrivial critical point of \( \Ie \) using the Mountain Pass Theorem (see \cite[Theorem $1$]{ghoussoub_mountain_pass_theorem}). To accomplish this, it remains only to verify that
\begin{equation*}
c^*_{\lambda} < \frac{1}{N} S_{\mu}^{N/2}.
\end{equation*}
Using the standard inequality $(b + d)^n \geq b^n + d^n + n b^{n-1} d$, for $n > 1; \; b,d > 0$, 
% \begin{equation*}
% (b + d)^n \geq b^n + d^n + n b^{n-1} d, \quad \text{for } n > 1; \; b,d > 0,
% \end{equation*}
we deduce that
\begin{equation*}
g(x, a) \geq a^p + p  u_\lambda^{p-1} a.
\end{equation*}
Choosing \( \delta \) small enough such that
\[
u_{\lambda} \geq M_0 > 0 \quad \text{in } B_{2 \de} \setminus \{0\}
\]
for some $M_0>0$ independent of $\varepsilon$. Then using the functions $u_{\varepsilon,\al}$ defined in subsection \ref{sec:superlinear}, we obtain
\begin{equation*} 
\begin{aligned}
\Ie(t u_{\varepsilon,\al}) &\leq \frac{t^2}{2} \left( \int_{\Omega} |\nabla u_{\varepsilon,\al}|^2 - \mu \int_{\Omega} \frac{u_{\varepsilon,\al}^2}{|x|^2} + \varepsilon[u_{\varepsilon,\al}]_s^2 - M_0^{2^*-2}(2^*-1) \int_{\Omega} u_{\varepsilon,\al}^2 \, dx \right) 
- \frac{t^{2^*}}{2^*} \int_{\Omega} u_{\varepsilon,\al}^{2^*} \, dx\\
& = \frac{t^2}{2}A - \frac{t^{2^*}}{2^*}B\\
& \leq \frac{1}{N}   \frac{A^{2^* / (2^* - 2)}}{B^{2 / (2^* - 2)}}
\end{aligned}
\end{equation*}
Using \eqref{eq:norm_lower_bound}, \eqref{eq:norm_upper_bound}, \eqref{eq:fractional_norm_bound}, and \eqref{eq:lower_bound_norm_r} with $r = 2$, we have
\begin{equation*} 
\begin{aligned}
\Ie(t u_{\varepsilon,\al}) &\leq \frac{1}{N} \frac{\left( S_{\mu}^{N/2} + C \varepsilon^{\al (N - 2)} + C \varepsilon^{1+2 \al (1-s)\sqrt{\bar{\mu}}/\sqrt{\bar{\mu} - \mu}} - 
 C \varepsilon^{2 \al \sqrt{\bar{\mu}}/\sqrt{\bar{\mu} - \mu}}  \right)^{2^*/(2^*-2)}}{\left(S_{\mu}^{N/2}-C \varepsilon^{\al N}\right)^{2/(2^*-2)}}\\
& = \frac{1}{N} S_{\mu}^{N/2}  \frac{(1 - C\varepsilon^{2 \al \sqrt{\bar{\mu}}/\sqrt{\bar{\mu}-\mu}})^{N/2}}{(1 - C \varepsilon^{\al N})^{(N - 2)/2}} < \frac{1}{N} S_{\mu}^{N/2}.
\end{aligned}
\end{equation*}
provided that \( \mu < \bar{\mu} - 1 \) and we choose $\al$ such that $\min\{\al(N-2), 1+2 \al (1-s)\sqrt{\bar{\mu}}/\sqrt{\bar{\mu} - \mu}\} > 2 \al \sqrt{\bar{\mu}}/\sqrt{\bar{\mu} - \mu}$.
Therefore, we conclude
\begin{equation*} 
c^*_{\lambda} \leq \max_{t > 0} \Ie(t u_{\varepsilon,\al}) < \frac{1}{N} S_{\mu}^{N/2}. \qedhere
\end{equation*}
\end{proof}
%%%%%%%%%%%%%%%%%%%%%%%%%%%%%%%%%%%%%%%%%%%%%%%%%%%%%%%%%%
%%%%%%%%%%%%%%%%%%%%%%%%%%%%% End of first 2 solutions
%%%%%%%%%%%%%%%%%%%%%%%%%%%%%%%%%%%%%%%%%%%%%%%%%%%%%%%%%%

%%%%%%%%%%%%%%%%%%%%%%%%%%%%%%%%%%%%%%%%%%%%%%%%%%%%%%%%%%
%%%%%%%%%%%%%%%%%%%%%%%%%%% Infinitely many solutions 
%%%%%%%%%%%%%%%%%%%%%%%%%%%%%%%%%%%%%%%%%%%%%%%%%%%%%%%%%%

% \subsection{Infinitely many solutions}

% In this section, we prove Theorem \ref{thm:sublinear_infinite_solutions}. For this, we assume along the section that $q\in(1,2)$ without further mentioning. Firstly, we recall the definition of genus and some its fundamental properties; see \cite{rabinowitz1986minimax} for more details.

% Let $E$ be a Banach space and $A$ a subset of $E$. We say that $A$ is symmetric if $u\in A$ implies that $-u\in A$. For a closed symmetric set $A$ which does not contain the origin, we define the genus $\gamma(A)$ of $A$ as
% \begin{equation*}\label{genus}
% \begin{split}
% \gamma^* (A) = 
% \begin{cases}
% 0 & \mbox{ if } A = \emptyset, \\
% \inf\{ m \in \ntrl \cup \{ 0 \} : \mbox{ there exists } h \in C(A;\mathbb{R}^m \setminus \{ 0 \}), h(-u) = h(u) \}  & \mbox{ if } A \neq \emptyset,
% \end{cases}
% \end{split}
% \end{equation*}
% if no such infimum exists then we take $\gamma^*(A) = \infty $.

% Then we have the following result.
% \begin{proposition}\label{prop:genus_properties}
% Let $A$ and $B$ be closed symmetric subsets of $E$ which do not contain the origin. Then we have:
% \begin{enumerate}
% \item[$(i)$] if there exists an odd continuous mapping from $A$ to $B$, then $\gamma(A)\leq \gamma(B)$;
% \item[$(ii)$] if there is an odd homeomorphism from $A$ onto $B$, then $\gamma(A)= \gamma(B)$;
% \item[$(iii)$] if $\gamma(B)<\infty$, then $\gamma(A\setminus B)\geq \gamma(A)-\gamma(B)$;
% \item[$(iv)$] the $k$-dimensional sphere $\mathbb{S}^{k}$ has a genus of $k+1$ by the Borsuk-Ulam Theorem;
% \item[$(v)$] if $A$ is compact, then $\gamma(A)<\infty$ and there exist $\delta>0$ and a closed and symmetric neighborhood $N_{\delta}(A)=\{x\in E:\,\, \text{dist}(x,A)\leq \delta \}$ of $A$ such that $\gamma(N_{\delta}(A))=\gamma(A)$.
% \end{enumerate}
% \end{proposition}


% We also need the following deformation lemma, as given in  \cite[Lemma 1.3]{ambrosetti1973dual}. For this, considering  $\mathcal{I}\in C^{1}(E;\mathbb R)$, we set for any $c\in\mathbb R$
% \begin{equation}\label{kappac}
% K_c=\left\{u\in E:\,\,\mathcal{I}^{\prime}(u)=0\mbox{ and }\mathcal{I}(u)=c\right\}.
% \end{equation}
% Also, we use the following notation $\mathcal{I}^d=\left\{u\in E:\,\, \mathcal{I}(u)\leq d\right\}$ for $d\in\mathbb R$.

% \begin{lemma}\label{prop:deformation_lemma}
% Let $E$ be an infinite-dimensional Banach space and let $\mathcal{I}\in C^{1}(E;\mathbb R)$ be a functional satisfying the $(PS)_c$, with $c\in \mathbb{R}$. Let $U$ be any neighborhood of $K_c$. Then, there exists $\eta\in C^0([0,1]\times E; E)$, with  $\eta_t(x)=\eta (t,x)$, and constants $d_1$, $\varepsilon>0$ with $|c|>d_1$, such that:
% \begin{itemize}
% \item[$(i)$] $\eta_0= \mathcal{I}d_E$;

% \item[$(ii)$]$\eta_t(u)=u$  for any   $u\not\in \mathcal{I}^{-1}([c-\varepsilon,c+\varepsilon])$  and  $t\in [0,1]$;

% \item[$(iii)$] $\eta_t$ is a homeomorphism for any $t\in [0,1]$;

% \item[$(iv)$] $\mathcal{I}(\eta_t (u))\leq \mathcal{I}(u)$ for any $u\in E$  and $t\in [0,1]$;

% \item[$(v)$] $\eta_1 (\mathcal{I}^{c+d_1}\setminus U)\subseteq \mathcal{I}^{c-d_1}$;

% \item[$(vi)$] if $K_c =\emptyset$, then $\eta_1 (\mathcal{I}^{c+d_1})\subseteq \mathcal{I}^{c-d_1}$;

% \item[$(vii)$] if $\mathcal{I}$ is even, then $\eta_t$ is odd for any $t\in [0,1]$.
% \end{itemize}
% \end{lemma}


% \medskip

% \begin{lemma}\label{lem:sublinear_ps_bdd}
% Let \( \lambda > 0 \). If \( \{u_n\} \subset \X(\Omega) \) is a sequence satisfying the Palais–Smale (PS)\(_c\) condition, then it is bounded in $\X(\om)$. 
% \end{lemma}
% \begin{proof}
% We begin by evaluating
% \begin{equation*}
%     \J(u_n) - \frac{1}{2^*} \langle \J^{\prime}(u_n), u_n \rangle = \frac{1}{N} \left( \int_{\Omega} |\nabla u_n|^2 + [u_n]^2_{s} \, dx - \mu \int_{\Omega} \frac{|u_n|^2}{|x|^2} \, dx \right) 
%     - \lambda \left( \frac{1}{p} - \frac{1}{2^*} \right) \|u_n\|_{L^q}^q.
% \end{equation*}
% Since \( \{u_n\} \) is a (PS)\(_c\) sequence, we obtain the inequality:
% \begin{equation}\label{eq:ps-boundedness}
%     C \left(1 + \|u_n\|_{\X}\right) \geq \frac{1}{N} \left(1 - \frac{\mu}{\bar{\mu}} \right) \|u_n\|_{\X}^2 
%     - \lambda \left( \frac{1}{p} - \frac{1}{2^*} \right) \|u_n\|_{\X}^q,
% \end{equation}
% which implies that the sequence \( \{u_n\} \) is bounded.
% \end{proof}

% \begin{lemma}\label{lem:ps_threshold_sublinear}
% If \( q \in (1, 2) \), then the (PS)\(_c\) condition for the functional \( \J \) holds for all    
% \begin{equation}\label{eq:ps-critical}
%     c < \frac{1}{N} S^{N/2} - |\Omega| \left( \frac{1}{2} - \frac{1}{2^*} \right)^{-\frac{q}{2^* - q}} 
%     \left[ \lambda \left( \frac{1}{p} - \frac{1}{2^*} \right) \right]^{\frac{2^*}{2^* - q}}.
% \end{equation}
% \end{lemma}


% \begin{proof}
% Since \( \{u_n\} \) is a (PS)\(_c\) sequence, by Lemma \ref{lem:sublinear_ps_bdd}, it is bounded in \( \X(\Omega) \). Thus, there exists \( u \in \X(\Omega) \) such that
% \begin{align*}
%     u_n &\rightharpoonup u \quad \text{in } \X(\Omega), \\
%     u_n &\rightarrow u \quad \text{in } L^r(\Omega) \text{ for } r \in [1, 2^*), \\
%     u_n(x) &\rightarrow u(x) \quad \text{pointwise a.e. in } \Omega.
% \end{align*}
% \noi Consider 
% \begin{align*}
%     o_n(1) & = \langle \J^{\prime}(u_n), u_n - u \rangle \\
%     & = \left( \|\nabla u_n\|^2_2 - \|\nabla u\|^2_2 \right) + \left( [u_n]^2_{s} - [u]^2_{s} \right)
%     - \mu \int_{\Omega} \frac{|u_n|^2 - |u|^2}{|x|^2} \, dx - (\|u_n\|_{L^{2^*}}^{2^*} - \|u\|_{L^{2^*}}^{2^*}) + o_n(1) \\
%     & = \|\nabla u_n - \nabla u\|^2_2 + [u_n - u]^2_{s} 
%     - \mu \int_{\Omega} \frac{|u_n - u|^2}{|x|^2} \, dx - \|u_n - u\|_{L^{2^*}}^{2^*} + o_n(1),
% \end{align*}
% where the final equality follows from the Brezis–Lieb Lemma. Thus, we conclude that:
% \begin{equation*}
%     \lim_{n \to \infty} \left( \|u_n - u\|^2_{\X} - \mu \int_{\Omega} \frac{|u_n - u|^2}{|x|^2} \, dx \right)
%     = \lim_{n \to \infty} \|u_n - u\|_{L^{2^*}}^{2^*} = \ell  (\text{say}).
% \end{equation*}

% \noi By \eqref{eq:spectral_ratio_frac}, we also have
% \begin{equation*}
%     S_{\mu} \|u_n - u\|_{L^{2^*}}^2 \leq \|u_n - u\|^2_{\X} - \mu \int_{\Omega} \frac{|u_n - u|^2}{|x|^2} \, dx.
% \end{equation*}
% \noi This implies
% \begin{equation*}
%     S_\mu \, \ell^{\frac{2}{2^*}} \leq \ell.
% \end{equation*}
% \noi If \( \ell > 0 \), then
% \begin{equation*}
%     \ell \geq S_\mu^{\frac{N}{2}}.
% \end{equation*}
% \noi Now, consider
% \begin{align*}
%     \J(u_n) - \frac{1}{2} \langle \J^{\prime}(u_n), u_n \rangle 
%     &= \left( \frac{1}{2} - \frac{1}{2^*} \right) \|u_n\|_{L^{2^*}}^{2^*} - \lambda \int_{\Omega} |u_n|^q \left( \frac{1}{p} - \frac{1}{2} \right) dx.
% \end{align*}

% \noi By the (PS)\(_c\) condition, H\"{o}lder's inequality, and Young’s inequality, we deduce
% \begin{align*}
%     C &\geq \left( \frac{1}{2} - \frac{1}{2^*} \right) \left[ \ell + \|u\|_{L^{2^*}}^{2^*} \right]
%     - \lambda \left( \frac{1}{p} - \frac{1}{2} \right) \|u\|_{L^q}^q \\
%     &\geq \left( \frac{1}{2} - \frac{1}{2^*} \right) \left[ \ell + \|u\|_{L^{2^*}}^{2^*} \right]
%     - \lambda \left( \frac{1}{p} - \frac{1}{2} \right) |\om|^{(2^*-q)/2^*} \|u\|_{L^{2^*}}^2 \\
%     &\geq \left( \frac{1}{2} - \frac{1}{2^*} \right) \left[ \ell + \|u\|_{L^{2^*}}^{2^*} \right]
%     - |\Omega| \left( \frac{1}{2} - \frac{1}{2^*} \right)^{-q/(2^* - q)}
%     \left[ \lambda \left( \frac{1}{p} - \frac{1}{2} \right) \right]^{2^*/(2^* - q)}.
% \end{align*}
% This contradicts the definition of \( c \), so \( \ell = 0 \). Hence, we have the desired result.
% \end{proof}

% \noi Since the functional \( \J \) is not bounded below, we apply a truncation argument to define another functional \( \Jtilde \) with desirable properties.

% Consider:
% \begin{equation*}
%     \J(u) \geq \frac{1}{2} \left( 1 - \frac{\mu}{\bar{\mu}} \right) \|u\|_{\X}^2 - \frac{\lambda C_q}{2} \|u\|_{\X}^q 
%     - \frac{S^{2^*/2}}{2^*} \|u\|_{\X}^{2^*} := g_\lambda(\|u\|_{\X}),
% \end{equation*}
% where \( g_\lambda : [0, \infty) \to \mathbb{R} \) is strictly decreasing with respect to \( \lambda \). Now fix \( R_1 > 0 \) such that
% \begin{equation*}
%     \frac{1}{2} \left( 1 - \frac{\mu}{\bar{\mu}} \right) R_1^2 - \frac{S^{2^*/2}}{2^*} R_1^{2^*} > 0.
% \end{equation*}

% \noi Also, define
% \begin{equation*}
%     \lambda_0 = \frac{q}{2 C_q R_1^2}   
%     \left( \frac{1}{2} \left( 1 - \frac{\mu}{\bar{\mu}} \right) R_1^2 - \frac{S^{2^*/2}}{2^*} R_1^{2^*} \right) > 0.
% \end{equation*}
% Then, \( g_{\lambda_0}(R_1) = \lambda_0 > 0 \). Define
% \begin{equation*}
%     R_0 = \max \left\{ t \in (0, R_1) : g_{\lambda_0}(t) \leq 0 \right\}.
% \end{equation*}

% \noi Now, consider a cut-off function \( \psi \in C^1([0,\infty)) \) such that
% \[
% 0 \leq \psi(t) \leq 1, \quad 
% \psi(t) = 1 \text{ for } t \in [0, R_0], \quad 
% \psi(t) = 0 \text{ for } t \in [R_1, \infty).
% \]

% \noi Then define the truncated functional \( \Jtilde : \X(\Omega) \to \mathbb{R} \) associated with \( \J \) as:
% \begin{equation}\label{def:truncated_functional}
% \Jtilde(u) = \frac{1}{2} \|u\|_{\X}^2 - \frac{\mu}{2} \int_{\Omega} \frac{|u|^2}{|x|^2} dx
% - \frac{\lambda}{q} \|u\|_q^q - \psi(\|u\|_{\X}^2) \frac{\|u\|_{2^*}^{2^*}}{2^*}
% \end{equation}
% Then clearly, \( \Jtilde \) is bounded and coercive.

% \begin{lemma}\label{lem:ps_condn_sublinear}
% There exists \( \lambda_* > 0 \) such that for any \( \lambda \in (0, \lambda_*) \), we have:
% \begin{enumerate}
%     \item[$(i)$] If \( \Jtilde(u) \leq 0 \), then \( \|u\| \leq R_0 \). In addition, there exists a neighborhood \( U \) of \( u \) such that \( \Jtilde(v) = \J(v) \) for all \( v \in U \).
%     \item[$(ii)$] \( \Jtilde \) satisfies the (PS)\(_c\) condition for any \( c < 0 \).
% \end{enumerate}
% \end{lemma}

% \begin{proof}
% $(i)$ Let \( \Jtilde(u) \leq 0 \).

% \textbf{Subcase 1:} If \( \|u\|_{\X} \geq R_1 \), then
% \[
% \Jtilde(u) \geq \frac{1}{2} \left(1 - \frac{\mu}{\bar{\mu}} \right) \|u\|_{\X}^2 - \frac{\lambda C_q}{q} \|u\|_{\X}^q := h_\lambda(\|u\|_{\X}).
% \]

% Then \( h_\lambda \) has two roots, namely \( t_0 = 0 \) and
% \[
% t_{1,\lambda} = \left[ \frac{2 \lambda C_q}{q} \left(1 - \frac{\mu}{\bar{\mu}} \right)^{-1} \right]^{1/(2 - q)}.
% \]

% So, if we take
% \[
% \lambda_1^* = R_1^{2 - q} \left(1 - \frac{\mu}{\bar{\mu}} \right) \frac{q}{2 C_q},
% \]
% then \( t_{1,\lambda} < R_1 \) for all \( \lambda < \min \{ \lambda_0, \lambda_1^* \} \). This implies that
% \[
% 0 \geq \Jtilde(u) \geq h_\lambda(\|u\|_{\X}) > 0,
% \]
% which is a contradiction.

% \textbf{Subcase 2:} When \( \|u\|_{\X} < R_1 \), we have
% \[
% 0 \geq \Jtilde(u) \geq g_\lambda(\|u\|_{\X}) > g_{\lambda_0}(\|u\|_{\X}),
% \]
% implying \( \|u\|_{\X} < R_0 \), by the definition of \( g_{\lambda_0} \).

% Moreover, by the continuity of \( \Jtilde \), there exists a neighborhood \( U \) of \( u \) such that \( U \subset B(0, R_0) \) and \( \Jtilde(v) = \J(v) \) for all \( v \in U \).
% \medskip
% $(ii)$ To prove this, we choose \( \lambda_2^* \) sufficiently small such that
% \[
% \frac{1}{N} S_{\mu}^{N/2} - |\Omega| \left( \frac{1}{2} - \frac{1}{2^*} \right)^{-q/(2^* - q)} 
% \left[ \lambda_2^* \left( \frac{1}{p} - \frac{1}{2} \right) \right]^{2^*/(2^* - q)} > 0.
% \]
% Then taking \( \lambda^* = \min \{ \lambda_0, \lambda_1^*, \lambda_2^* \} \), and letting \( \lambda \in (0, \lambda^*) \), if \( (u_n) \subset \X(\Omega) \) is a (PS)\(_c\) sequence for large \( n \), we have
% \[
% \Jtilde(u_n) = \J(u_n), \quad \text{and} \quad \Jtilde'(u_n) = \J^{\prime}(u_n).
% \]
% Now, we can apply Lemma \ref{lem:ps_threshold_sublinear} to deduce our result.
% \end{proof}

% \begin{lemma}\label{lem:existence_of_k_genus_set}
% Let \( \lambda > 0 \) and \( k \in \mathbb{N} \). Then there exists \( \varepsilon = \varepsilon(\lambda, k) > 0 \) such that
% \[
% \gamma(\Jtilde^{-\varepsilon}) \geq k,
% \]
% where
% \[
% \Jtilde^{-\varepsilon} = \{ u \in \X(\Omega) : \Jtilde(u) \leq -\varepsilon \}.
% \]
% \end{lemma}

% \begin{proof}
% To prove this, we consider a \( k \)-dimensional subspace \( V_k \subset \X(\Omega) \). Then \( \| \|_{\X} \) and \( \| \|_q \) are equivalent on \( V_k \).
% Thus, there exists \( C(k) > 0 \) such that
% \[
% C(k) \|u\|_{\X}^q \leq \|u\|_q^q \quad \text{for any } u \in V_k.
% \]

% For any \( u \in V_k \) with \( \|u\|_{\X} \leq R_0 \), we get:
% \[
% \Jtilde(u) = \J(u)
% \leq \frac{1}{2} \|u\|_{\X}^2 - \frac{\mu}{2} \int_{\om} \frac{|u|^2}{|x|^2} dx - \frac{\lambda C(k)}{q} \|u\|_{\X}^q
% \leq \frac{1}{2} \|u\|_{\X}^2 - \frac{\lambda C(k)}{q} \|u\|_{\X}^q.
% \]

% Let \( \beta > 0 \) be such that
% \[
% \beta < \min \left\{ R_0, \left( \frac{2 \lambda C(k)}{q} \right)^{1/(2 - q)} \right\}.
% \]
% Take \( S_\beta = \{ u \in V_k : \|u\|_{\X} = \beta \} \). Then \( S_\beta \) is homeomorphic to the \((k-1)\)-dimensional unit sphere \( \mathbb{S}^{k-1} \). As a consequence, for any \( u \in S_\beta \), we have:
% \[
% \Jtilde(u) \leq \beta^q \left( \frac{1}{2} \beta^{2 - q} - \frac{\lambda C(k)}{q} \right) < 0.
% \]
% Hence, there exists \( \varepsilon = \varepsilon(k, \lambda) > 0 \) such that:
% \[
% S_\beta \subset \Jtilde^{-\varepsilon} \setminus \{0\}.
% \]
% Thus, by Proposition \ref{prop:genus_properties}, we have
% \[
% \gamma(\Jtilde^{-\varepsilon}) \geq \gamma(S_\beta) = k.
% \]
% \end{proof}

% Now, set the numbers
% \[
% c_k = \inf_{A \in \Sigma_k} \sup_{u \in A} \Jtilde(u),
% \]
% where
% \[
% \Sigma_k = \{ A \subset \X(\Omega) : A \text{ closed, symmetric, } 0 \notin A, \text{ and } \gamma(A) \geq k \}.
% \]
% Clearly, \( c_k \leq c_{k+1} \).

% \begin{lemma}\label{lem:ck_negative}
% For \( \lambda > 0 \), \( k \in \mathbb{N} \), we have \( c_k < 0 \).
% \end{lemma}

% \begin{proof}
% From Lemma \ref{lem:existence_of_k_genus_set}, there exists \( \varepsilon = \varepsilon(k, \lambda) > 0 \) such that
% \[
% \gamma(\Jtilde^{-\varepsilon}) \geq k.
% \]

% Moreover, as \( \Jtilde \) is continuous and even, it implies \( \Jtilde^{-\varepsilon} \) is closed and symmetric. In addition, \( 0 \notin \Jtilde^{-\varepsilon} \).
% Also, as \( \Jtilde \) is bounded from below, we deduce \( c_k > -\infty \).  
% Therefore,
% \[
% c_k \leq \sup_{u \in \Jtilde^{-\varepsilon}} \Jtilde(u) \leq -\varepsilon < 0.
% \]
% \end{proof}

% \noi Next, we define the set of critical points of \( \Jtilde \) at level \( c \):
% \[
% K_c = \left\{ u \in \X(\Omega) \,\middle|\, \Jtilde'(u) = 0 \text{ and } \Jtilde(u) = c \right\}.
% \]

% \begin{lemma}\label{lem:genus_lower_bound}
% Let \( \lambda \in (0, \lambda^*) \) and \( k \in \mathbb{N} \). If \( c = c_k =  s = c_{k+m} \) for some \( m \in \mathbb{N} \), then
% \[
% \gamma(K_c) \geq m+1.
% \]
% \end{lemma}

% \begin{proof}
% Proof follows as in Lemma 3.6 of \cite{silva2024mixed}.
% \end{proof}

% \begin{proof}[Proof of Theorem \ref{thm:sublinear_infinite_solutions}]
%   Let $\La$ as in Lemma \ref{lem:ps_condn_sublinear}, then fix any $\la \in (0,\La)$. Since $c_k<0$ from Lemma \ref{lem:ck_negative}, part $(ii)$ of Lemma \ref{lem:ps_condn_sublinear} implies that the functional satisfies the $(PS)_{c_k}$ condition. Thus $c_k$ is a critical value of $\Jtilde$ for any $k \in \ntrl$ (see \cite{rabinowitz1986minimax}). \\
%   Now, if all $c_k$ are distinct there is nothing to prove. So assume there exist $k, m \in \ntrl$ such that $c = c_k =  s = c_{k+m}$, then by Lemma \ref{lem:genus_lower_bound}, $\gamma(K_c) \geq m+1 \geq 2$. Thus the set $K_c$ has infinitely many critical values for $\Jtilde$ by part $(ii)$ of Lemma \ref{lem:ps_condn_sublinear}.\\
%   Finally part $(i)$ implies the existence of infinitely many critical points of $\J = \Jtilde$. This gives our result. 
% \end{proof}






%%%%%%%%%%%%%%%%%%%%%%%%%%%%%%%%%%%%%%%%%%%%%%%%%%%%%%%%%%
%%%%%%%%%%%%%%%%%%%%%%%%%%% Acknowledgment
%%%%%%%%%%%%%%%%%%%%%%%%%%%%%%%%%%%%%%%%%%%%%%%%%%%%%%%%%%
\vspace{.2cm}
{\textbf{Acknowledgment:}}\\
The author, Shammi Malhotra, is supported by the Prime Minister’s Research Fellowship (PMRF ID - 1403226). The author, Sarika Goyal, would like to thank the Anusandhan National Research Foundation, Department of Science and Technology, Government of India for the financial support under the grant SPG/2022/002068.



% This can be decomposed as:
% \begin{equation*}
%     \iint_{B_{1/m} \times B_{1/m}} \frac{|U_{\varepsilon}(x) - U_{\varepsilon}(y)|^2}{|x-y|^{N+2s}} dy \, dx \quad \text{(denoted as $I_1$)}
% \end{equation*}
% \begin{equation*}
%     + 2 \iint_{B_{1/m} \times B_{1/m}} \frac{|U_{\varepsilon}(x)|^2}{|x-y|^{N+2s}} dy \, dx \quad \text{(denoted as $I_2$)}.
% \end{equation*}

% We further analyze $I_1$:
% \begin{equation*}
%     I_1 = \iint_{B_{1/m} \times B_{1/m}} \frac{|U_{\varepsilon}(x) - U_{\varepsilon}(y)|^2}{|x-y|^{N+2s}} dy \, dx.
% \end{equation*}
% Using the transformation:
% \begin{equation*}
%     x = \varepsilon^{\sqrt{\bar{\mu}}/\sqrt{\bar{\mu}-\mu}} z, \quad y = \varepsilon^{\sqrt{\bar{\mu}}/\sqrt{\bar{\mu}-\mu}} w,
% \end{equation*}
% we deduce:
% \begin{equation*}
%     I_1 \leq \varepsilon^{2\sqrt{\bar{\mu}}} \left( -2s \frac{\sqrt{\bar{\mu}}}{\sqrt{\bar{\mu}-\mu}} - 2 \frac{\sqrt{\bar{\mu}}}{\sqrt{\bar{\mu}-\mu}} + 2N \frac{\sqrt{\bar{\mu}}}{\sqrt{\bar{\mu}-\mu}} \right) \iint_{\mathbb{R}^{2N}} \frac{|U_1(z) - U_1(w)|^2}{|z-w|^2} dz dw.
% \end{equation*}

% To estimate $I_2$, we define two sets:
% \begin{equation*}
%     \mathcal{D} = \{(x,y) \in \mathbb{R}^{2N} \mid x \in B_{1/m}, y \in (B_{1/m})^c, |x-y| > \frac{1}{2m} \},
% \end{equation*}
% \begin{equation*}
%     \mathcal{E} = \{(x,y) \in \mathbb{R}^{2N} \mid x \in B_{1/m}, y \in (B_{1/m})^c, |x-y| \leq \frac{1}{2m} \}.
% \end{equation*}

% Claim 1:
% If $x \in (B_{1/2m})^c$, then
% \begin{equation*}
%     |U_{\varepsilon}^{m}(x)| \leq B_0 \varepsilon^{\sqrt{\bar{\mu}}} (2m)^{\delta}.
% \end{equation*}

% Proof of Claim 1:
% \begin{equation*}
%     |U_{\varepsilon}^{m}(x)| \leq |U_{\varepsilon}(x)| = B_0 \varepsilon^{\sqrt{\bar{\mu}}} \frac{1}{\left( \varepsilon^2 |x|^{\gamma^{\prime}/\sqrt{\bar{\mu}}} + |x|^{x/\sqrt{\bar{\mu}}} \right) \sqrt{\bar{\mu}}}.
% \end{equation*}
% From this, we obtain the bound:
% \begin{equation*}
%     |U_{\varepsilon}^{m}(x)| \leq B_0 \varepsilon^{\sqrt{\bar{\mu}}} |x|^{-\delta} \leq C \varepsilon^{\sqrt{\bar{\mu}}} m^{\delta}.
% \end{equation*}

% Claim 2:
% If $x \in (B_{1/2m})^c$, then...

% \subsection*{Proof of Claim 2}
% We have
% \begin{equation*}
%     |\nabla u_\varepsilon^m (x)| = |\nabla u_\varepsilon (x)| \leq C \left[ \varepsilon^{2 \sqrt{\bar{\mu}}} m^{1+ \gamma} + 2 \sqrt{\bar{\mu}} \varepsilon^{\sqrt{\bar{\mu}}} m^{1+ \gamma} + \varepsilon^{\frac{\sqrt{\bar{\mu}}}{4\sqrt{\bar{\mu} - \mu}}} \right].
% \end{equation*}

% From the definition of $u_\varepsilon (x)$,
% \begin{equation*}
%     \nabla u_\varepsilon (x) = B_0 \varepsilon \left( -\sqrt{\bar{\mu}} \right) \frac{ \left[ \varepsilon^2 |x|^{\gamma} / \sqrt{\bar{\mu}} - 2 \right] x + \gamma \frac{x}{\sqrt{\bar{\mu}}} }{\left( \varepsilon^2 |x|^{\gamma} / \sqrt{\bar{\mu}} + |x|^{\gamma} / \sqrt{\bar{\mu}} \right) \sqrt{\bar{\mu}}}.
% \end{equation*}

% Taking the modulus and simplifying,
% \begin{equation*}
%     |\nabla u_\varepsilon (x)| \leq C \varepsilon \sqrt{\bar{\mu}} \left( \frac{ \varepsilon^2 |x|^{\gamma} / \sqrt{\bar{\mu}} - 1 }{ \left( \varepsilon^2 |x|^{\gamma} / \sqrt{\bar{\mu}} + |x|^{\gamma} / \sqrt{\bar{\mu}} \right) \sqrt{\bar{\mu}} } \right).
% \end{equation*}

% Further simplifying,
% \begin{equation*}
%     \leq C \varepsilon \sqrt{\bar{\mu}} \left( \varepsilon^2 |x|^{\gamma} / \sqrt{\bar{\mu}} + |x|^{\gamma} / \sqrt{\bar{\mu}} \right)^{-1}.
% \end{equation*}

% \subsection*{Claim 3}
% If $x \in \mathbb{R}^N$, $y \in B_{1/m}^c$, and $|x - y| \leq \frac{1}{2m}$, then
% \begin{equation*}
%     |u_\varepsilon (x) - u_\varepsilon (y)| \leq C \left[ \varepsilon^{2 \sqrt{\bar{\mu}}} m^{1+ \gamma} + 2 \sqrt{\bar{\mu}} \varepsilon^{\sqrt{\bar{\mu}}} m^{1+ \gamma} \right] |x - y|.
% \end{equation*}

% \subsection*{Proof of Claim 3}
% If $\xi = tx + (1 - t) y$, then
% \begin{equation*}
%     |\xi| = |y + t(x - y)| \geq |y| - t |x - y| \geq \frac{1}{2m}.
% \end{equation*}

% Using Taylor expansions and Claim 2, we get
% \begin{equation*}
%     |u_\varepsilon (x) - u_\varepsilon (y)| \leq |\nabla u_\varepsilon (\xi)| |x - y| \leq C \left[ \varepsilon^{2 \sqrt{\bar{\mu}}} m^{1+ \gamma} + 2 \sqrt{\bar{\mu}} \varepsilon^{\sqrt{\bar{\mu}}} m^{1+ \gamma} \right] |x - y|.
% \end{equation*}

% \subsection*{Estimating $I_2$}
% \begin{equation*}
%     I_2 = \iint_E \frac{|u_\varepsilon^m (x) - u_\varepsilon^m (y)|^2}{|x - y|^{N + 2s}} dy dx + \iint_D \frac{|u_\varepsilon^m (x) - u_\varepsilon^m (y)|^2}{|x - y|^{N + 2s}} dy dx.
% \end{equation*}

% Defining $I_3$ and $L_1$ as follows,
% \begin{equation*}
%     I_3 = \iint_E \frac{|u_\varepsilon^m (x) - u_\varepsilon^m (y)|^2}{|x - y|^{N + 2s}} dy dx,
% \end{equation*}

% and using Claim 3,
% \begin{equation*}
%     I_3 \leq C \frac{\varepsilon^{2 \sqrt{\bar{\mu}}}}{1 - s} m^{-N + 4s - 2 \sqrt{\bar{\mu}}}.
% \end{equation*}

% For $L_1$, using integral estimates,
% \begin{equation*}
%     L_1 \leq \int_{B_{1/m}} |u_\varepsilon (x)|^2 dx \int_{(B_{1/2m})^c} \frac{1}{|y|^{N + 2s}} dy \leq C \varepsilon^{2 \sqrt{\bar{\mu}} / \sqrt{\bar{\mu} - \mu}} ||u_\varepsilon||_{L^2}^2 m^{-2s}.
% \end{equation*}

\bibliographystyle{abbrv}
\bibliography{references}
\end{document}
